% 高一10_复附浦东_周末作业1.tex —— 高一上数学“2026-2027 年复附浦东高一上周末作业 1”（试卷版/答案版）
% 试卷版：xelatex 高一10_复附浦东_周末作业1.tex
% 答案版：xelatex "\def\isanswer{1}% 高一10_复附浦东_周末作业1.tex —— 高一上数学“2026-2027 年复附浦东高一上周末作业 1”（试卷版/答案版）
% 试卷版：xelatex 高一10_复附浦东_周末作业1.tex
% 答案版：xelatex "\def\isanswer{1}% 高一10_复附浦东_周末作业1.tex —— 高一上数学“2026-2027 年复附浦东高一上周末作业 1”（试卷版/答案版）
% 试卷版：xelatex 高一10_复附浦东_周末作业1.tex
% 答案版：xelatex "\def\isanswer{1}% 高一10_复附浦东_周末作业1.tex —— 高一上数学“2026-2027 年复附浦东高一上周末作业 1”（试卷版/答案版）
% 试卷版：xelatex 高一10_复附浦东_周末作业1.tex
% 答案版：xelatex "\def\isanswer{1}\input{高一10_复附浦东_周末作业1.tex}"
% 紧凑版：xelatex "\def\iscompact{1}\input{高一10_复附浦东_周末作业1.tex}"
%
% 原始材料是微信公众号图文（公众号“魔都数学加油站”），正文为 9 张 PNG 页；原文本身就是一份
% **讲评版**——题目下方直接印出【答案】或【解析】。试题文字、答案与解析均照原卷录入，未改内容。
%
% 与原卷的排版差异（均已在 README“说明”中交代）：
%   ① 原文自**第 3 题**起（第 1、2 题未在该文中发布），故本稿题号亦自 3 起；
%   ② 原卷本就印有“一、填空题”“二、选择题”“三、解答题”三节标题（分别在原图 01/04/06.png
%      上，逐页放大核对过），本稿照原卷分节，只把标题改排成全稿统一的“大节标题”样式；
%   ③ 原卷填空、选择的答案直接印在题目下方，本稿统一改为试卷版隐去、答案版以【答案】给出，
%      使两版彻底分离；
%   ④ 原卷的实数集符号时作正体 $R$、时作粗体 $\mathbf{R}$（第 21(3) 与 (3) 末句均为粗体），
%      本稿按全稿惯例统一写作 $\R$.
%
% 原卷笔误一处（本稿已补，见该处上方的 `%` 注）：
%   ① 第 10 题【解析】首句原卷印作“这个元不是 $1$ 就是 $m$”，“元素”漏了“素”字——
%      本稿按 高一b 第 13 题“原卷漏排、本稿补上”的先例补出“素”字，其余照原卷。
% 除此一处外，本卷 9 页逐页放大核对，未发现别的笔误（题 16 的 $1995=15\times 133$ 与
% “$15^2$ 的倍数有 8 个”两处最易被误判为笔误，实测 $133$、$8$ 均正确）。

% 本篇在公众号“魔都数学加油站”连载里的编号是“27学年高一10”，本地编号与之对齐（高一10）；
% 对应关系见 README“与公众号连载号的对照表”。
\documentclass[a4paper,11pt]{article}
\usepackage{common}

\begin{document}

\examtitlest[3]{2026--2027 年复附浦东高一上周末作业 1}{集合与常用逻辑用语}

\bigsec{一、填空题}

\begin{enumerate}
\setcounter{enumi}{2}

\item 已知集合 $M=\{1,2m+1\}$，$N=\{-1,m^2\}$，且 $M=N$，则 $m=$ \blankline[2.4em].

\ans{$-1$}

\ifanswer
\solbegin
\step{若 $M=N$，必有 $\begin{cases}2m+1=-1\\ m^2=1\end{cases}$，解得 $m=-1$.}
\solend

\fi

\item 集合 $\{(x,y)\mid x+2y=2,\ x,y\in\N\}$ 用列举法表示为 \blankline[4.4em].

\ans{$\{(0,1),(2,0)\}$}

\ifanswer
\solbegin
\step{由题意得 $x+2y=2$ 可化为 $y=1-\dfrac{x}{2}$，}
\step{由 $y\in\N$，有 $y=1-\dfrac{x}{2}\geqslant 0$，解得 $x\leqslant 2$，}
\step{又由 $x\in\N$，得 $x$ 可能取值为 $0$，$1$，$2$，}
\step{当 $x=0$ 时，得 $y=1$，满足条件，}
\step{当 $x=1$ 时，得 $y=\dfrac{1}{2}\notin\N$，不满足条件，}
\step{当 $x=2$ 时，得 $y=0$，满足条件，}
\step{综上，集合用列举法表示为 $\{(0,1),(2,0)\}$.}
\solend

\fi

\item 已知集合 $A=\{1\}$，集合 $B=\{0,1,2\}$，则满足关系 $A\sub P\sub B$ 的所有集合 $P$ 为 \blankline[4.4em].

\ans{$\{1\},\{0,1\},\{1,2\},\{0,1,2\}$}

\ifanswer
\solbegin
\step{由题意，集合 $P$ 可以是 $\{1\},\{0,1\},\{1,2\},\{0,1,2\}$.}
\solend

\fi

\item 已知集合 $A=\{x\mid a-1\leqslant x\leqslant a+2\}$，$B=\{x\mid 3<x<5\}$，则使 $A\supseteq B$ 成立的实数 $a$ 的取值范围是 \blankline[4.4em].

\ans{$[3,4]$}

\ifanswer
\solbegin
\step{因为 $A\supseteq B$，所以 $\begin{cases}a+2\geqslant 5\\ a-1\leqslant 3\end{cases}$，解得 $3\leqslant a\leqslant 4$.}
\solend

\fi

\item 若集合 $\{x\mid mx^2-4x+1=0\}=\{n\}$，则 $m+n=$ \blankline[2.4em].

\ans{$\dfrac{1}{4}$ 或 $\dfrac{9}{2}$}

\ifanswer
\solbegin
\step{若集合 $\{x\mid mx^2-4x+1=0\}=\{n\}$，则 $mx^2-4x+1=0$ 有唯一实数根 $n$，}
\step{当 $m=0$ 时，得 $n=\dfrac{1}{4}$，$m+n=\dfrac{1}{4}$，}
\step{当 $m\neq 0$ 时，$\Delta=16-4m=0$，即 $m=4$，}
\step{此时 $4x^2-4x+1=0$ 的根为 $n=\dfrac{1}{2}$，$m+n=\dfrac{9}{2}$.}
\step{故 $m+n=\dfrac{1}{4}$ 或 $\dfrac{9}{2}$.}
\solend

\fi

\item 集合 $M=\{-1,2,3,5,6\}$ 的全部非空子集的元素和等于 \blankline[4.4em].

\ans{$240$}

\ifanswer
\solbegin
\step{集合 $M=\{-1,2,3,5,6\}$ 中共有 $5$ 个元素，}
\step{在其所有子集中，元素 $-1$ 出现的次数为 $2^{5-1}=2^4=16$，}
\step{同理 $2$、$3$、$5$、$6$ 都出现 $16$ 次，}
\step{则集合 $M$ 的全部非空子集的元素和为 $(-1+2+3+5+6)\times 16=240$.}
\solend

\fi

\item 设 $A=\{x\mid x^2-3x-4=0\}$，$B=\{x\mid ax-1=0\}$，若 $A\cap B=B$，则实数 $a$ 的值可以为 \blankline[4.4em].

\ans{$0$、$-1$、$\dfrac{1}{4}$}

\ifanswer
\solbegin
\step{$A=\{-1,4\}$，因为 $A\cap B=B$，所以 $B\sub A$，}
\step{①$a=0$ 时，$B=\e$，满足 $B\sub A$；}
\step{②$a\neq 0$ 时，$B=\left\{\dfrac{1}{a}\right\}$，则 $\dfrac{1}{a}=-1$ 或 $\dfrac{1}{a}=4$，所以 $a=-1$ 或 $\dfrac{1}{4}$；}
\step{所以 $a$ 的值可以为 $0$、$-1$、$\dfrac{1}{4}$.}
\solend

\fi

\item 设集合 $A=\{1,m\}$，其中 $m$ 为实数. 令集合 $B=\{a+b\mid a\in A,b\in A\}$，若 $A\cap B$ 恰有一个元素，则 $A\cup B$ 的元素之和为 \blankline[2.4em].

\ans{$5$ 或 $10$}

\ifanswer
\solbegin
% 注：原卷本句印作“这个元不是 $1$ 就是 $m$”——“元素”漏了“素”字。按 高一b 第 13 题
%     “原卷漏排、本稿补上”的先例，本稿补出“素”字，其余照原卷，见文件头“原卷笔误”。
\step{由题意得，若 $A\cap B$ 恰有一个元素，这个元素不是 $1$ 就是 $m$，}
\stepm{（1）}{当 $1\in B$ 时，}
\step{①若 $1=1+m$，则 $m=0$，此时 $A=\{1,0\},B=\{2,1,0\},A\cap B=\{0,1\}$，不满足题意，舍去；}
\step{②若 $1=2m$，则 $m=\dfrac{1}{2}$，此时 $A=\left\{1,\dfrac{1}{2}\right\}$，$B=\left\{2,\dfrac{3}{2},1\right\}$，$A\cap B=\{1\}$，}
\step{满足题意，所以 $A\cup B=\left\{1,\dfrac{1}{2},2,\dfrac{3}{2}\right\}$，其元素之和为 $5$，}
\stepm{（2）}{当 $m\in B$ 时，}
\step{①若 $m=1+m$，则 $m$ 无解，}
\step{②若 $m=2m$，则 $m=0$，此时 $A=\{1,0\},B=\{2,1,0\},A\cap B=\{0,1\}$，不满足题意，}
\step{③若 $m=2$，则 $A=\{1,2\},B=\{2,3,4\},A\cap B=\{2\}$，满足题意，}
\step{所以 $A\cup B=\{1,2,3,4\}$，其元素之和为 $10$，}
\step{综上所述，$A\cup B$ 的元素之和为 $5$ 或 $10$.}
\solend

\fi

\item 设非空集合 $Q\sub M$，当 $Q$ 中所有元素和为偶数时（集合为单元素时和为元素本身），称 $Q$ 是 $M$ 的偶子集. 若集合 $M=\{1,2,3,4,5,6,7\}$，则其偶子集 $Q$ 的个数为 \blankline[2.4em].

\ans{$63$}

\ifanswer
\solbegin
\step{偶子集可以分为只有奇数、偶数及偶数个奇数和偶数的集合，}
\step{奇数有 $4$ 个，偶数有 $3$ 个，}
\step{只有奇数的有 $\mathrm{C}_4^2+1=7$ 个，只有偶数的有 $\mathrm{C}_3^1+\mathrm{C}_3^2+1=7$ 个，}
\step{偶数个奇数和偶数的：两奇一偶 $\mathrm{C}_4^2\cdot\mathrm{C}_3^1=18$ 个，两奇二偶 $\mathrm{C}_4^2\cdot\mathrm{C}_3^2=18$ 个，}
\step{两奇三偶 $\mathrm{C}_4^2\cdot\mathrm{C}_3^3=6$ 个，四奇一偶 $\mathrm{C}_4^4\cdot\mathrm{C}_3^1=3$ 个，四奇二偶 $\mathrm{C}_4^4\cdot\mathrm{C}_3^2=3$ 个，}
\step{四奇三偶 $\mathrm{C}_4^4\cdot\mathrm{C}_3^3=1$ 个，所以共有 $7+7+18+18+6+3+3+1=63$ 个.}
\solend

\fi

\item 已知非空集合 $S=\left\{x\mid -\dfrac{1}{2}\leqslant x\leqslant m\right\}$ 满足：当 $x\in S$ 时，有 $x^2\in S$，则实数 $m$ 的取值范围是 \blankline[4.4em].

\ans{$\left[\dfrac{1}{4},1\right]$}

\ifanswer
\solbegin
\step{因为集合 $S=\left\{x\mid -\dfrac{1}{2}\leqslant x\leqslant m\right\}$ 是非空集合，所以 $m\geqslant -\dfrac{1}{2}$，}
\step{又当 $x\in S$ 时，有 $x^2\in S$，所以 $m^2\leqslant m$ 且 $m\geqslant \left(-\dfrac{1}{2}\right)^2=\dfrac{1}{4}$，所以 $\dfrac{1}{4}\leqslant m\leqslant 1$.}
\solend

\fi

\end{enumerate}

\bigsec{二、选择题}

\begin{enumerate}
\setcounter{enumi}{12}

\item 已知集合 $U=\{1+x,1\}$，$x\in U$，则 $x=$\mbox{（\quad）}

\keepwithprev
A. $-1$\qquad B. $0$\qquad C. $1$\qquad D. $2$

\ans{C}

\ifanswer
\solbegin
\step{因为集合 $U=\{1+x,1\}$，$x\in U$，所以 $1+x=x$ 或 $x=1$，}
\step{当 $x=1+x$ 时，即 $0=1$，此时显然不成立，舍去，}
\step{当 $x=1$，此时集合 $U=\{1,2\}$，满足互异性，}
\step{综上所述，$x=1$. 故选 C.}
\solend

\fi

\item 已知集合 $M=\{m\mid m=a+b\sqrt2,\ a,b\in\Q\}$，则下列四个元素中属于 $M$ 的元素的个数是\mbox{（\quad）}

①$1+\sqrt2\pi$；②$\sqrt{11+6\sqrt2}$；③$\dfrac{1}{2+\sqrt2}$；④$\sqrt{2-\sqrt3}+\sqrt{2+\sqrt3}$

\keepwithprev
A. $4$\qquad B. $3$\qquad C. $2$\qquad D. $1$

\ans{C}

\ifanswer
\solbegin
\step{①$1+\sqrt2\pi$，由于 $\sqrt2$、$\pi$ 都为无理数，因此 $1+\sqrt2\pi\notin M$；}
\step{②$\sqrt{11+6\sqrt2}=\sqrt{11+2\sqrt{18}}=3+\sqrt2\in M$；}
\step{③$\dfrac{1}{2+\sqrt2}=1-\dfrac{1}{2}\sqrt2\in M$；}
\step{④$\sqrt{2-\sqrt3}+\sqrt{2+\sqrt3}=\sqrt{\dfrac{8-2\sqrt{12}}{4}}+\sqrt{\dfrac{8+2\sqrt{12}}{4}}=\dfrac{\sqrt6-\sqrt2}{2}+\dfrac{\sqrt6+\sqrt2}{2}=\sqrt6=\sqrt3\cdot\sqrt2\notin M$；}
\step{则上述四个元素中属于 $M$ 的元素的个数是 $2$. 故选 C.}
\solend

\fi

\item 集合 $A=\{2,3,7,8,9,11\}$，集合 $B=\{1,2,4,6,8\}$，下面说法正确的是\mbox{（\quad）}

\keepwithprev
A. 两个集合的交集是 $\{2,8\}$

B. 两个集合没有并集

C. $B$ 集合的元素无法用 $\{x\mid x=f(k),k\in D\}$ 这种形式表示出来

D. 集合 $A$ 有 $6$ 个子集

\ans{A}

\ifanswer
\solbegin
\step{对于 A，$A\cap B=\{2,8\}$，故 A 正确；}
\step{对于 B，$A\cup B=\{1,2,3,4,6,7,8,9,11\}$，故 B 错误；}
\step{对于 C，集合 $B$ 可以表示为 $B=\{x\mid x=f(k),k\in\{1,2,3,4,5\}\}$，}
\step{其中 $f(1)=1,f(2)=2,f(3)=4,f(4)=6,f(5)=8$，}
\step{即集合 $B$ 可描述为 $\{x\mid x=f(k),k\in D\}$ 形式，故 C 错误；}
\step{对于 D，集合 $A$ 有 $6$ 个元素，其子集总数为 $2^6=64\neq 6$，故 D 错误.}
\step{故选 A.}
\solend

\fi

\item 设 $M=\{1,2,\cdots,1995\}$，$A$ 是 $M$ 的子集且满足条件：当 $x\in A$ 时，$15x\notin A$，则 $A$ 中元素的个数的最大值为\mbox{（\quad）}

\keepwithprev
A. $1862$\qquad B. $1866$\qquad C. $1868$\qquad D. $1870$

\ans{D}

\ifanswer
\solbegin
\step{$1995=15\times 133$. 故取出所有不是 $15$ 的倍数的数，共 $1862$ 个，}
\step{这些数均符合要求.}
\step{在所有 $15$ 的倍数的数中，$15^2$ 的倍数有 $8$ 个，这些数又可以取出，}
\step{这样共取出了 $1870$ 个，即 $|A|\geqslant 1870$.}
\step{又 $\{k,15k\}(k=9,10,11,\cdots,133)$ 中的两个元素不能同时取出，}
\step{故 $|A|\leqslant 1995-133+8=1870$，故选 D.}
\solend

\fi

\end{enumerate}

\bigsec{三、解答题}

\begin{enumerate}
\setcounter{enumi}{16}

\item 已知集合 $A=\{x\mid a\leqslant x\leqslant a+4\}$，$B=\{x\mid x\geqslant 6\}$.

\begin{enumerate}[label=(\arabic*)]
\item 若 $a=3$，求 $A\cup B$；
\item 若 $A\cap B=A$，求实数 $a$ 的取值范围.
\end{enumerate}

\ans{（1）$\{x\mid x\geqslant 3\}$；（2）$[6,+\infty)$}

\ifanswer
\solbegin
\stepm{（1）}{当 $a=3$ 时，$A=\{x\mid 3\leqslant x\leqslant 7\}$，所以 $A\cup B=\{x\mid x\geqslant 3\}$；}
\stepm{（2）}{若 $A\cap B=A$，则 $A\sub B$，所以 $a\geqslant 6$，所以实数 $a$ 的取值范围是 $[6,+\infty)$.}
\solend

\fi

\ansspace[7]

\item 已知集合 $M=\{x\mid (x-a)(x^2-ax+a-1)=0\}$.

\begin{enumerate}[label=(\arabic*)]
\item 若 $a=3$，求集合 $M$；
\item 若集合 $M$ 中各元素之和等于 $3$，求实数 $a$ 的值，并用列举法表示集合 $M$.
\end{enumerate}

\ans{（1）$\{1,2,3\}$；（2）$a=2$ 或 $\dfrac{3}{2}$，此时 $M=\{1,2\}$ 或 $M=\left\{\dfrac{1}{2},1,\dfrac{3}{2}\right\}$}

\ifanswer
\solbegin
\stepm{（1）}{$\{1,2,3\}$；}
\stepm{（2）}{当 $(x-a)(x^2-ax+a-1)=0$ 有重根时，重根只能算作集合的一个元素，}
\step{$(x-a)(x^2-ax+a-1)=(x-a)(x-1)[x-(a-1)]=0$，}
\step{当 $a=1$ 时，得 $M=\{0,1\}$，不符合题意；}
\step{当 $a-1=1$ 时，即 $a=2$ 时，得 $M=\{1,2\}$，符合题意；}
\step{当 $a\neq 1$ 且 $a\neq 2$ 时，此时 $M=\{1,a,a-1\}$，得 $1+a+a-1=3$，}
\step{解得 $a=\dfrac{3}{2}$，此时 $M=\left\{\dfrac{1}{2},1,\dfrac{3}{2}\right\}$，符合题意，}
\step{综上，实数 $a$ 的值为 $2$ 或 $\dfrac{3}{2}$，}
\step{当 $a=2$ 时，$M=\{1,2\}$；当 $a=\dfrac{3}{2}$ 时，$M=\left\{\dfrac{1}{2},1,\dfrac{3}{2}\right\}$.}
\solend

\fi

\ansspace[8]

\item 已知集合 $A=\{x\mid x^2-2x-3=0\}$，集合 $B=\{x\mid ax-1=0\}$. 若 $B\sub A$，求实数 $a$ 的值.

\ans{$a=0$ 或 $a=-1$ 或 $a=\dfrac{1}{3}$}

\ifanswer
\solbegin
\step{当 $B=\e$ 时，$a=0$，满足题意；}
\step{当 $B\neq\e$，即 $a\neq 0$ 时，$B=\left\{\dfrac{1}{a}\right\}$，}
\step{又 $A=\{x\mid x^2-2x-3=0\}=\{x\mid x=-1\text{ 或 }x=3\}$，$B\sub A$，}
\step{所以 $\dfrac{1}{a}=-1$ 或 $\dfrac{1}{a}=3$，所以 $a=-1$ 或 $a=\dfrac{1}{3}$.}
\step{综上所述，$a=0$ 或 $a=-1$ 或 $a=\dfrac{1}{3}$.}
\solend

\fi

\ansspace[8]

\item 设集合 $A=\{x\mid -2\leqslant x\leqslant 5\}$，$B=\{x\mid m+1\leqslant x\leqslant 2m-1\}$.

\begin{enumerate}[label=(\arabic*)]
\item 若 $B\sub A$，求实数 $m$ 的取值范围；
\item 当 $x\in\Z$ 时，求 $A$ 的非空真子集个数；
\item 当 $x\in\R$ 时，不存在元素 $x$ 使 $x\in A$ 与 $x\in B$ 同时成立，求实数 $m$ 的取值范围.
\end{enumerate}

\ans{（1）$\{m\mid m\leqslant 3\}$；（2）$254$；（3）$\{m\mid m<2\text{ 或 }m>4\}$}

\ifanswer
\solbegin
\stepm{（1）}{当 $m+1>2m-1$，即 $m<2$ 时，$B=\e$，满足 $B\sub A$.}
\step{当 $m+1\leqslant 2m-1$，即 $m\geqslant 2$ 时，要使 $B\sub A$ 成立，}
\step{需 $\begin{cases}m+1\geqslant -2\\ 2m-1\leqslant 5\end{cases}$，得 $2\leqslant m\leqslant 3$，}
\step{综上，$m$ 的范围为 $\{m\mid m\leqslant 3\}$；}
\stepm{（2）}{$A=\{x\mid -2\leqslant x\leqslant 5,x\in\Z\}=\{-2,-1,0,1,2,3,4,5\}$，}
\step{则 $A$ 的非空真子集个数为 $2^8-1-1=254$；}
\stepm{（3）}{因为没有元素 $x$ 使 $x\in A$ 与 $x\in B$ 同时成立，所以 $A$ 与 $B$ 交集为空集.}
\step{①若 $B=\e$，即 $m+1>2m-1$，得 $m<2$ 时满足条件；}
\step{②若 $B\neq\e$，则要满足的条件是 $\begin{cases}m+1\leqslant 2m-1\\ m+1>5\end{cases}$ 或 $\begin{cases}m+1\leqslant 2m-1\\ 2m-1<-2\end{cases}$，}
\step{解得 $m>4$.}
\step{综上，有 $m<2$ 或 $m>4$，所以 $m$ 的范围为 $\{m\mid m<2\text{ 或 }m>4\}$.}
\solend

\fi

\ansspace[10]

% 压轴题：其后已无内容，故作答区全用可丢弃胶水（\ansspace*），并在题目前先量一下版面。
\ansneed{7cm}

\item 已知集合 $A=\{x\mid ax=1\}$，$B=\{x\mid x^2-ax+b=0\}$.

\begin{enumerate}[label=(\arabic*)]
\item 若 $a=2$，且 $A\sub B$，求 $b$ 的值；
\item 当 $a^2\geqslant 4b$ 时，若 $B\sub A$，求 $a$、$b$ 的值；
\item 若 $A\sub B$，讨论 $a$、$b$ 的取值范围.
\end{enumerate}

\ans{（1）$\dfrac{3}{4}$；（2）$a=\pm\sqrt2$，$b=\dfrac{1}{2}$；（3）$a=0$ 时 $b\in\R$，$a\neq 0$ 时 $b\in(-\infty,1)$}

\ifanswer
\solbegin
\stepm{（1）}{当 $a=2$ 时，$A=\left\{\dfrac{1}{2}\right\}$，}
\step{由 $A\sub B$ 得 $\dfrac{1}{2}$ 是方程 $x^2-2x+b=0$ 的解，代入得 $\left(\dfrac{1}{2}\right)^2-2\times\dfrac{1}{2}+b=0$，}
\step{解得 $b=\dfrac{3}{4}$.}
\stepm{（2）}{当 $a^2\geqslant 4b$ 时，方程 $x^2-ax+b=0$ 有实根.}
\step{若 $a=0$，则 $A=\e$，要 $B\sub A$ 需 $B=\e$，}
\step{但此时方程 $x^2+b=0$ 在 $a^2\geqslant 4b$（即 $b\leqslant 0$）时有实根，矛盾，舍去.}
\step{若 $a\neq 0$，则 $A=\left\{\dfrac{1}{a}\right\}$，$B\sub A$ 且方程有实根，故方程有且仅有一个根 $\dfrac{1}{a}$.}
\step{由韦达定理，根的和为 $a=\dfrac{1}{a}+\dfrac{1}{a}$，根的积为 $b=\dfrac{1}{a}\times\dfrac{1}{a}$.}
\step{由根的和得 $a=\dfrac{2}{a}$，即 $a^2=2$，$a=\pm\sqrt2$，$b=\dfrac{1}{2}$.}
\stepm{（3）}{当 $a=0$ 时，$A=\e$，故 $A\sub B$ 恒成立，此时 $b\in\R$.}
\step{当 $a\neq 0$ 时，$A=\left\{\dfrac{1}{a}\right\}$，将 $x=\dfrac{1}{a}$ 代入 $x^2-ax+b=0$，得 $b=1-\dfrac{1}{a^2}$.}
\step{此时方程判别式 $\Delta=a^2-4b=\dfrac{(a^2-2)^2}{a^2}\geqslant 0$，}
\step{若 $\Delta=0$，则 $a=\pm\sqrt2$，$b=\dfrac{1}{2}$；}
\step{若 $\Delta>0$，则 $a\neq 0$ 且 $a\neq\pm\sqrt2$，此时 $b=1-\dfrac{1}{a^2}<1$.}
\step{综上，当 $a=0$ 时，$b\in\R$；当 $a\neq 0$ 时，$b\in(-\infty,1)$.}
\solend

\fi

\ansspace*[12]

\end{enumerate}

\end{document}
"
% 紧凑版：xelatex "\def\iscompact{1}% 高一10_复附浦东_周末作业1.tex —— 高一上数学“2026-2027 年复附浦东高一上周末作业 1”（试卷版/答案版）
% 试卷版：xelatex 高一10_复附浦东_周末作业1.tex
% 答案版：xelatex "\def\isanswer{1}\input{高一10_复附浦东_周末作业1.tex}"
% 紧凑版：xelatex "\def\iscompact{1}\input{高一10_复附浦东_周末作业1.tex}"
%
% 原始材料是微信公众号图文（公众号“魔都数学加油站”），正文为 9 张 PNG 页；原文本身就是一份
% **讲评版**——题目下方直接印出【答案】或【解析】。试题文字、答案与解析均照原卷录入，未改内容。
%
% 与原卷的排版差异（均已在 README“说明”中交代）：
%   ① 原文自**第 3 题**起（第 1、2 题未在该文中发布），故本稿题号亦自 3 起；
%   ② 原卷本就印有“一、填空题”“二、选择题”“三、解答题”三节标题（分别在原图 01/04/06.png
%      上，逐页放大核对过），本稿照原卷分节，只把标题改排成全稿统一的“大节标题”样式；
%   ③ 原卷填空、选择的答案直接印在题目下方，本稿统一改为试卷版隐去、答案版以【答案】给出，
%      使两版彻底分离；
%   ④ 原卷的实数集符号时作正体 $R$、时作粗体 $\mathbf{R}$（第 21(3) 与 (3) 末句均为粗体），
%      本稿按全稿惯例统一写作 $\R$.
%
% 原卷笔误一处（本稿已补，见该处上方的 `%` 注）：
%   ① 第 10 题【解析】首句原卷印作“这个元不是 $1$ 就是 $m$”，“元素”漏了“素”字——
%      本稿按 高一b 第 13 题“原卷漏排、本稿补上”的先例补出“素”字，其余照原卷。
% 除此一处外，本卷 9 页逐页放大核对，未发现别的笔误（题 16 的 $1995=15\times 133$ 与
% “$15^2$ 的倍数有 8 个”两处最易被误判为笔误，实测 $133$、$8$ 均正确）。

% 本篇在公众号“魔都数学加油站”连载里的编号是“27学年高一10”，本地编号与之对齐（高一10）；
% 对应关系见 README“与公众号连载号的对照表”。
\documentclass[a4paper,11pt]{article}
\usepackage{common}

\begin{document}

\examtitlest[3]{2026--2027 年复附浦东高一上周末作业 1}{集合与常用逻辑用语}

\bigsec{一、填空题}

\begin{enumerate}
\setcounter{enumi}{2}

\item 已知集合 $M=\{1,2m+1\}$，$N=\{-1,m^2\}$，且 $M=N$，则 $m=$ \blankline[2.4em].

\ans{$-1$}

\ifanswer
\solbegin
\step{若 $M=N$，必有 $\begin{cases}2m+1=-1\\ m^2=1\end{cases}$，解得 $m=-1$.}
\solend

\fi

\item 集合 $\{(x,y)\mid x+2y=2,\ x,y\in\N\}$ 用列举法表示为 \blankline[4.4em].

\ans{$\{(0,1),(2,0)\}$}

\ifanswer
\solbegin
\step{由题意得 $x+2y=2$ 可化为 $y=1-\dfrac{x}{2}$，}
\step{由 $y\in\N$，有 $y=1-\dfrac{x}{2}\geqslant 0$，解得 $x\leqslant 2$，}
\step{又由 $x\in\N$，得 $x$ 可能取值为 $0$，$1$，$2$，}
\step{当 $x=0$ 时，得 $y=1$，满足条件，}
\step{当 $x=1$ 时，得 $y=\dfrac{1}{2}\notin\N$，不满足条件，}
\step{当 $x=2$ 时，得 $y=0$，满足条件，}
\step{综上，集合用列举法表示为 $\{(0,1),(2,0)\}$.}
\solend

\fi

\item 已知集合 $A=\{1\}$，集合 $B=\{0,1,2\}$，则满足关系 $A\sub P\sub B$ 的所有集合 $P$ 为 \blankline[4.4em].

\ans{$\{1\},\{0,1\},\{1,2\},\{0,1,2\}$}

\ifanswer
\solbegin
\step{由题意，集合 $P$ 可以是 $\{1\},\{0,1\},\{1,2\},\{0,1,2\}$.}
\solend

\fi

\item 已知集合 $A=\{x\mid a-1\leqslant x\leqslant a+2\}$，$B=\{x\mid 3<x<5\}$，则使 $A\supseteq B$ 成立的实数 $a$ 的取值范围是 \blankline[4.4em].

\ans{$[3,4]$}

\ifanswer
\solbegin
\step{因为 $A\supseteq B$，所以 $\begin{cases}a+2\geqslant 5\\ a-1\leqslant 3\end{cases}$，解得 $3\leqslant a\leqslant 4$.}
\solend

\fi

\item 若集合 $\{x\mid mx^2-4x+1=0\}=\{n\}$，则 $m+n=$ \blankline[2.4em].

\ans{$\dfrac{1}{4}$ 或 $\dfrac{9}{2}$}

\ifanswer
\solbegin
\step{若集合 $\{x\mid mx^2-4x+1=0\}=\{n\}$，则 $mx^2-4x+1=0$ 有唯一实数根 $n$，}
\step{当 $m=0$ 时，得 $n=\dfrac{1}{4}$，$m+n=\dfrac{1}{4}$，}
\step{当 $m\neq 0$ 时，$\Delta=16-4m=0$，即 $m=4$，}
\step{此时 $4x^2-4x+1=0$ 的根为 $n=\dfrac{1}{2}$，$m+n=\dfrac{9}{2}$.}
\step{故 $m+n=\dfrac{1}{4}$ 或 $\dfrac{9}{2}$.}
\solend

\fi

\item 集合 $M=\{-1,2,3,5,6\}$ 的全部非空子集的元素和等于 \blankline[4.4em].

\ans{$240$}

\ifanswer
\solbegin
\step{集合 $M=\{-1,2,3,5,6\}$ 中共有 $5$ 个元素，}
\step{在其所有子集中，元素 $-1$ 出现的次数为 $2^{5-1}=2^4=16$，}
\step{同理 $2$、$3$、$5$、$6$ 都出现 $16$ 次，}
\step{则集合 $M$ 的全部非空子集的元素和为 $(-1+2+3+5+6)\times 16=240$.}
\solend

\fi

\item 设 $A=\{x\mid x^2-3x-4=0\}$，$B=\{x\mid ax-1=0\}$，若 $A\cap B=B$，则实数 $a$ 的值可以为 \blankline[4.4em].

\ans{$0$、$-1$、$\dfrac{1}{4}$}

\ifanswer
\solbegin
\step{$A=\{-1,4\}$，因为 $A\cap B=B$，所以 $B\sub A$，}
\step{①$a=0$ 时，$B=\e$，满足 $B\sub A$；}
\step{②$a\neq 0$ 时，$B=\left\{\dfrac{1}{a}\right\}$，则 $\dfrac{1}{a}=-1$ 或 $\dfrac{1}{a}=4$，所以 $a=-1$ 或 $\dfrac{1}{4}$；}
\step{所以 $a$ 的值可以为 $0$、$-1$、$\dfrac{1}{4}$.}
\solend

\fi

\item 设集合 $A=\{1,m\}$，其中 $m$ 为实数. 令集合 $B=\{a+b\mid a\in A,b\in A\}$，若 $A\cap B$ 恰有一个元素，则 $A\cup B$ 的元素之和为 \blankline[2.4em].

\ans{$5$ 或 $10$}

\ifanswer
\solbegin
% 注：原卷本句印作“这个元不是 $1$ 就是 $m$”——“元素”漏了“素”字。按 高一b 第 13 题
%     “原卷漏排、本稿补上”的先例，本稿补出“素”字，其余照原卷，见文件头“原卷笔误”。
\step{由题意得，若 $A\cap B$ 恰有一个元素，这个元素不是 $1$ 就是 $m$，}
\stepm{（1）}{当 $1\in B$ 时，}
\step{①若 $1=1+m$，则 $m=0$，此时 $A=\{1,0\},B=\{2,1,0\},A\cap B=\{0,1\}$，不满足题意，舍去；}
\step{②若 $1=2m$，则 $m=\dfrac{1}{2}$，此时 $A=\left\{1,\dfrac{1}{2}\right\}$，$B=\left\{2,\dfrac{3}{2},1\right\}$，$A\cap B=\{1\}$，}
\step{满足题意，所以 $A\cup B=\left\{1,\dfrac{1}{2},2,\dfrac{3}{2}\right\}$，其元素之和为 $5$，}
\stepm{（2）}{当 $m\in B$ 时，}
\step{①若 $m=1+m$，则 $m$ 无解，}
\step{②若 $m=2m$，则 $m=0$，此时 $A=\{1,0\},B=\{2,1,0\},A\cap B=\{0,1\}$，不满足题意，}
\step{③若 $m=2$，则 $A=\{1,2\},B=\{2,3,4\},A\cap B=\{2\}$，满足题意，}
\step{所以 $A\cup B=\{1,2,3,4\}$，其元素之和为 $10$，}
\step{综上所述，$A\cup B$ 的元素之和为 $5$ 或 $10$.}
\solend

\fi

\item 设非空集合 $Q\sub M$，当 $Q$ 中所有元素和为偶数时（集合为单元素时和为元素本身），称 $Q$ 是 $M$ 的偶子集. 若集合 $M=\{1,2,3,4,5,6,7\}$，则其偶子集 $Q$ 的个数为 \blankline[2.4em].

\ans{$63$}

\ifanswer
\solbegin
\step{偶子集可以分为只有奇数、偶数及偶数个奇数和偶数的集合，}
\step{奇数有 $4$ 个，偶数有 $3$ 个，}
\step{只有奇数的有 $\mathrm{C}_4^2+1=7$ 个，只有偶数的有 $\mathrm{C}_3^1+\mathrm{C}_3^2+1=7$ 个，}
\step{偶数个奇数和偶数的：两奇一偶 $\mathrm{C}_4^2\cdot\mathrm{C}_3^1=18$ 个，两奇二偶 $\mathrm{C}_4^2\cdot\mathrm{C}_3^2=18$ 个，}
\step{两奇三偶 $\mathrm{C}_4^2\cdot\mathrm{C}_3^3=6$ 个，四奇一偶 $\mathrm{C}_4^4\cdot\mathrm{C}_3^1=3$ 个，四奇二偶 $\mathrm{C}_4^4\cdot\mathrm{C}_3^2=3$ 个，}
\step{四奇三偶 $\mathrm{C}_4^4\cdot\mathrm{C}_3^3=1$ 个，所以共有 $7+7+18+18+6+3+3+1=63$ 个.}
\solend

\fi

\item 已知非空集合 $S=\left\{x\mid -\dfrac{1}{2}\leqslant x\leqslant m\right\}$ 满足：当 $x\in S$ 时，有 $x^2\in S$，则实数 $m$ 的取值范围是 \blankline[4.4em].

\ans{$\left[\dfrac{1}{4},1\right]$}

\ifanswer
\solbegin
\step{因为集合 $S=\left\{x\mid -\dfrac{1}{2}\leqslant x\leqslant m\right\}$ 是非空集合，所以 $m\geqslant -\dfrac{1}{2}$，}
\step{又当 $x\in S$ 时，有 $x^2\in S$，所以 $m^2\leqslant m$ 且 $m\geqslant \left(-\dfrac{1}{2}\right)^2=\dfrac{1}{4}$，所以 $\dfrac{1}{4}\leqslant m\leqslant 1$.}
\solend

\fi

\end{enumerate}

\bigsec{二、选择题}

\begin{enumerate}
\setcounter{enumi}{12}

\item 已知集合 $U=\{1+x,1\}$，$x\in U$，则 $x=$\mbox{（\quad）}

\keepwithprev
A. $-1$\qquad B. $0$\qquad C. $1$\qquad D. $2$

\ans{C}

\ifanswer
\solbegin
\step{因为集合 $U=\{1+x,1\}$，$x\in U$，所以 $1+x=x$ 或 $x=1$，}
\step{当 $x=1+x$ 时，即 $0=1$，此时显然不成立，舍去，}
\step{当 $x=1$，此时集合 $U=\{1,2\}$，满足互异性，}
\step{综上所述，$x=1$. 故选 C.}
\solend

\fi

\item 已知集合 $M=\{m\mid m=a+b\sqrt2,\ a,b\in\Q\}$，则下列四个元素中属于 $M$ 的元素的个数是\mbox{（\quad）}

①$1+\sqrt2\pi$；②$\sqrt{11+6\sqrt2}$；③$\dfrac{1}{2+\sqrt2}$；④$\sqrt{2-\sqrt3}+\sqrt{2+\sqrt3}$

\keepwithprev
A. $4$\qquad B. $3$\qquad C. $2$\qquad D. $1$

\ans{C}

\ifanswer
\solbegin
\step{①$1+\sqrt2\pi$，由于 $\sqrt2$、$\pi$ 都为无理数，因此 $1+\sqrt2\pi\notin M$；}
\step{②$\sqrt{11+6\sqrt2}=\sqrt{11+2\sqrt{18}}=3+\sqrt2\in M$；}
\step{③$\dfrac{1}{2+\sqrt2}=1-\dfrac{1}{2}\sqrt2\in M$；}
\step{④$\sqrt{2-\sqrt3}+\sqrt{2+\sqrt3}=\sqrt{\dfrac{8-2\sqrt{12}}{4}}+\sqrt{\dfrac{8+2\sqrt{12}}{4}}=\dfrac{\sqrt6-\sqrt2}{2}+\dfrac{\sqrt6+\sqrt2}{2}=\sqrt6=\sqrt3\cdot\sqrt2\notin M$；}
\step{则上述四个元素中属于 $M$ 的元素的个数是 $2$. 故选 C.}
\solend

\fi

\item 集合 $A=\{2,3,7,8,9,11\}$，集合 $B=\{1,2,4,6,8\}$，下面说法正确的是\mbox{（\quad）}

\keepwithprev
A. 两个集合的交集是 $\{2,8\}$

B. 两个集合没有并集

C. $B$ 集合的元素无法用 $\{x\mid x=f(k),k\in D\}$ 这种形式表示出来

D. 集合 $A$ 有 $6$ 个子集

\ans{A}

\ifanswer
\solbegin
\step{对于 A，$A\cap B=\{2,8\}$，故 A 正确；}
\step{对于 B，$A\cup B=\{1,2,3,4,6,7,8,9,11\}$，故 B 错误；}
\step{对于 C，集合 $B$ 可以表示为 $B=\{x\mid x=f(k),k\in\{1,2,3,4,5\}\}$，}
\step{其中 $f(1)=1,f(2)=2,f(3)=4,f(4)=6,f(5)=8$，}
\step{即集合 $B$ 可描述为 $\{x\mid x=f(k),k\in D\}$ 形式，故 C 错误；}
\step{对于 D，集合 $A$ 有 $6$ 个元素，其子集总数为 $2^6=64\neq 6$，故 D 错误.}
\step{故选 A.}
\solend

\fi

\item 设 $M=\{1,2,\cdots,1995\}$，$A$ 是 $M$ 的子集且满足条件：当 $x\in A$ 时，$15x\notin A$，则 $A$ 中元素的个数的最大值为\mbox{（\quad）}

\keepwithprev
A. $1862$\qquad B. $1866$\qquad C. $1868$\qquad D. $1870$

\ans{D}

\ifanswer
\solbegin
\step{$1995=15\times 133$. 故取出所有不是 $15$ 的倍数的数，共 $1862$ 个，}
\step{这些数均符合要求.}
\step{在所有 $15$ 的倍数的数中，$15^2$ 的倍数有 $8$ 个，这些数又可以取出，}
\step{这样共取出了 $1870$ 个，即 $|A|\geqslant 1870$.}
\step{又 $\{k,15k\}(k=9,10,11,\cdots,133)$ 中的两个元素不能同时取出，}
\step{故 $|A|\leqslant 1995-133+8=1870$，故选 D.}
\solend

\fi

\end{enumerate}

\bigsec{三、解答题}

\begin{enumerate}
\setcounter{enumi}{16}

\item 已知集合 $A=\{x\mid a\leqslant x\leqslant a+4\}$，$B=\{x\mid x\geqslant 6\}$.

\begin{enumerate}[label=(\arabic*)]
\item 若 $a=3$，求 $A\cup B$；
\item 若 $A\cap B=A$，求实数 $a$ 的取值范围.
\end{enumerate}

\ans{（1）$\{x\mid x\geqslant 3\}$；（2）$[6,+\infty)$}

\ifanswer
\solbegin
\stepm{（1）}{当 $a=3$ 时，$A=\{x\mid 3\leqslant x\leqslant 7\}$，所以 $A\cup B=\{x\mid x\geqslant 3\}$；}
\stepm{（2）}{若 $A\cap B=A$，则 $A\sub B$，所以 $a\geqslant 6$，所以实数 $a$ 的取值范围是 $[6,+\infty)$.}
\solend

\fi

\ansspace[7]

\item 已知集合 $M=\{x\mid (x-a)(x^2-ax+a-1)=0\}$.

\begin{enumerate}[label=(\arabic*)]
\item 若 $a=3$，求集合 $M$；
\item 若集合 $M$ 中各元素之和等于 $3$，求实数 $a$ 的值，并用列举法表示集合 $M$.
\end{enumerate}

\ans{（1）$\{1,2,3\}$；（2）$a=2$ 或 $\dfrac{3}{2}$，此时 $M=\{1,2\}$ 或 $M=\left\{\dfrac{1}{2},1,\dfrac{3}{2}\right\}$}

\ifanswer
\solbegin
\stepm{（1）}{$\{1,2,3\}$；}
\stepm{（2）}{当 $(x-a)(x^2-ax+a-1)=0$ 有重根时，重根只能算作集合的一个元素，}
\step{$(x-a)(x^2-ax+a-1)=(x-a)(x-1)[x-(a-1)]=0$，}
\step{当 $a=1$ 时，得 $M=\{0,1\}$，不符合题意；}
\step{当 $a-1=1$ 时，即 $a=2$ 时，得 $M=\{1,2\}$，符合题意；}
\step{当 $a\neq 1$ 且 $a\neq 2$ 时，此时 $M=\{1,a,a-1\}$，得 $1+a+a-1=3$，}
\step{解得 $a=\dfrac{3}{2}$，此时 $M=\left\{\dfrac{1}{2},1,\dfrac{3}{2}\right\}$，符合题意，}
\step{综上，实数 $a$ 的值为 $2$ 或 $\dfrac{3}{2}$，}
\step{当 $a=2$ 时，$M=\{1,2\}$；当 $a=\dfrac{3}{2}$ 时，$M=\left\{\dfrac{1}{2},1,\dfrac{3}{2}\right\}$.}
\solend

\fi

\ansspace[8]

\item 已知集合 $A=\{x\mid x^2-2x-3=0\}$，集合 $B=\{x\mid ax-1=0\}$. 若 $B\sub A$，求实数 $a$ 的值.

\ans{$a=0$ 或 $a=-1$ 或 $a=\dfrac{1}{3}$}

\ifanswer
\solbegin
\step{当 $B=\e$ 时，$a=0$，满足题意；}
\step{当 $B\neq\e$，即 $a\neq 0$ 时，$B=\left\{\dfrac{1}{a}\right\}$，}
\step{又 $A=\{x\mid x^2-2x-3=0\}=\{x\mid x=-1\text{ 或 }x=3\}$，$B\sub A$，}
\step{所以 $\dfrac{1}{a}=-1$ 或 $\dfrac{1}{a}=3$，所以 $a=-1$ 或 $a=\dfrac{1}{3}$.}
\step{综上所述，$a=0$ 或 $a=-1$ 或 $a=\dfrac{1}{3}$.}
\solend

\fi

\ansspace[8]

\item 设集合 $A=\{x\mid -2\leqslant x\leqslant 5\}$，$B=\{x\mid m+1\leqslant x\leqslant 2m-1\}$.

\begin{enumerate}[label=(\arabic*)]
\item 若 $B\sub A$，求实数 $m$ 的取值范围；
\item 当 $x\in\Z$ 时，求 $A$ 的非空真子集个数；
\item 当 $x\in\R$ 时，不存在元素 $x$ 使 $x\in A$ 与 $x\in B$ 同时成立，求实数 $m$ 的取值范围.
\end{enumerate}

\ans{（1）$\{m\mid m\leqslant 3\}$；（2）$254$；（3）$\{m\mid m<2\text{ 或 }m>4\}$}

\ifanswer
\solbegin
\stepm{（1）}{当 $m+1>2m-1$，即 $m<2$ 时，$B=\e$，满足 $B\sub A$.}
\step{当 $m+1\leqslant 2m-1$，即 $m\geqslant 2$ 时，要使 $B\sub A$ 成立，}
\step{需 $\begin{cases}m+1\geqslant -2\\ 2m-1\leqslant 5\end{cases}$，得 $2\leqslant m\leqslant 3$，}
\step{综上，$m$ 的范围为 $\{m\mid m\leqslant 3\}$；}
\stepm{（2）}{$A=\{x\mid -2\leqslant x\leqslant 5,x\in\Z\}=\{-2,-1,0,1,2,3,4,5\}$，}
\step{则 $A$ 的非空真子集个数为 $2^8-1-1=254$；}
\stepm{（3）}{因为没有元素 $x$ 使 $x\in A$ 与 $x\in B$ 同时成立，所以 $A$ 与 $B$ 交集为空集.}
\step{①若 $B=\e$，即 $m+1>2m-1$，得 $m<2$ 时满足条件；}
\step{②若 $B\neq\e$，则要满足的条件是 $\begin{cases}m+1\leqslant 2m-1\\ m+1>5\end{cases}$ 或 $\begin{cases}m+1\leqslant 2m-1\\ 2m-1<-2\end{cases}$，}
\step{解得 $m>4$.}
\step{综上，有 $m<2$ 或 $m>4$，所以 $m$ 的范围为 $\{m\mid m<2\text{ 或 }m>4\}$.}
\solend

\fi

\ansspace[10]

% 压轴题：其后已无内容，故作答区全用可丢弃胶水（\ansspace*），并在题目前先量一下版面。
\ansneed{7cm}

\item 已知集合 $A=\{x\mid ax=1\}$，$B=\{x\mid x^2-ax+b=0\}$.

\begin{enumerate}[label=(\arabic*)]
\item 若 $a=2$，且 $A\sub B$，求 $b$ 的值；
\item 当 $a^2\geqslant 4b$ 时，若 $B\sub A$，求 $a$、$b$ 的值；
\item 若 $A\sub B$，讨论 $a$、$b$ 的取值范围.
\end{enumerate}

\ans{（1）$\dfrac{3}{4}$；（2）$a=\pm\sqrt2$，$b=\dfrac{1}{2}$；（3）$a=0$ 时 $b\in\R$，$a\neq 0$ 时 $b\in(-\infty,1)$}

\ifanswer
\solbegin
\stepm{（1）}{当 $a=2$ 时，$A=\left\{\dfrac{1}{2}\right\}$，}
\step{由 $A\sub B$ 得 $\dfrac{1}{2}$ 是方程 $x^2-2x+b=0$ 的解，代入得 $\left(\dfrac{1}{2}\right)^2-2\times\dfrac{1}{2}+b=0$，}
\step{解得 $b=\dfrac{3}{4}$.}
\stepm{（2）}{当 $a^2\geqslant 4b$ 时，方程 $x^2-ax+b=0$ 有实根.}
\step{若 $a=0$，则 $A=\e$，要 $B\sub A$ 需 $B=\e$，}
\step{但此时方程 $x^2+b=0$ 在 $a^2\geqslant 4b$（即 $b\leqslant 0$）时有实根，矛盾，舍去.}
\step{若 $a\neq 0$，则 $A=\left\{\dfrac{1}{a}\right\}$，$B\sub A$ 且方程有实根，故方程有且仅有一个根 $\dfrac{1}{a}$.}
\step{由韦达定理，根的和为 $a=\dfrac{1}{a}+\dfrac{1}{a}$，根的积为 $b=\dfrac{1}{a}\times\dfrac{1}{a}$.}
\step{由根的和得 $a=\dfrac{2}{a}$，即 $a^2=2$，$a=\pm\sqrt2$，$b=\dfrac{1}{2}$.}
\stepm{（3）}{当 $a=0$ 时，$A=\e$，故 $A\sub B$ 恒成立，此时 $b\in\R$.}
\step{当 $a\neq 0$ 时，$A=\left\{\dfrac{1}{a}\right\}$，将 $x=\dfrac{1}{a}$ 代入 $x^2-ax+b=0$，得 $b=1-\dfrac{1}{a^2}$.}
\step{此时方程判别式 $\Delta=a^2-4b=\dfrac{(a^2-2)^2}{a^2}\geqslant 0$，}
\step{若 $\Delta=0$，则 $a=\pm\sqrt2$，$b=\dfrac{1}{2}$；}
\step{若 $\Delta>0$，则 $a\neq 0$ 且 $a\neq\pm\sqrt2$，此时 $b=1-\dfrac{1}{a^2}<1$.}
\step{综上，当 $a=0$ 时，$b\in\R$；当 $a\neq 0$ 时，$b\in(-\infty,1)$.}
\solend

\fi

\ansspace*[12]

\end{enumerate}

\end{document}
"
%
% 原始材料是微信公众号图文（公众号“魔都数学加油站”），正文为 9 张 PNG 页；原文本身就是一份
% **讲评版**——题目下方直接印出【答案】或【解析】。试题文字、答案与解析均照原卷录入，未改内容。
%
% 与原卷的排版差异（均已在 README“说明”中交代）：
%   ① 原文自**第 3 题**起（第 1、2 题未在该文中发布），故本稿题号亦自 3 起；
%   ② 原卷本就印有“一、填空题”“二、选择题”“三、解答题”三节标题（分别在原图 01/04/06.png
%      上，逐页放大核对过），本稿照原卷分节，只把标题改排成全稿统一的“大节标题”样式；
%   ③ 原卷填空、选择的答案直接印在题目下方，本稿统一改为试卷版隐去、答案版以【答案】给出，
%      使两版彻底分离；
%   ④ 原卷的实数集符号时作正体 $R$、时作粗体 $\mathbf{R}$（第 21(3) 与 (3) 末句均为粗体），
%      本稿按全稿惯例统一写作 $\R$.
%
% 原卷笔误一处（本稿已补，见该处上方的 `%` 注）：
%   ① 第 10 题【解析】首句原卷印作“这个元不是 $1$ 就是 $m$”，“元素”漏了“素”字——
%      本稿按 高一b 第 13 题“原卷漏排、本稿补上”的先例补出“素”字，其余照原卷。
% 除此一处外，本卷 9 页逐页放大核对，未发现别的笔误（题 16 的 $1995=15\times 133$ 与
% “$15^2$ 的倍数有 8 个”两处最易被误判为笔误，实测 $133$、$8$ 均正确）。

% 本篇在公众号“魔都数学加油站”连载里的编号是“27学年高一10”，本地编号与之对齐（高一10）；
% 对应关系见 README“与公众号连载号的对照表”。
\documentclass[a4paper,11pt]{article}
\usepackage{common}

\begin{document}

\examtitlest[3]{2026--2027 年复附浦东高一上周末作业 1}{集合与常用逻辑用语}

\bigsec{一、填空题}

\begin{enumerate}
\setcounter{enumi}{2}

\item 已知集合 $M=\{1,2m+1\}$，$N=\{-1,m^2\}$，且 $M=N$，则 $m=$ \blankline[2.4em].

\ans{$-1$}

\ifanswer
\solbegin
\step{若 $M=N$，必有 $\begin{cases}2m+1=-1\\ m^2=1\end{cases}$，解得 $m=-1$.}
\solend

\fi

\item 集合 $\{(x,y)\mid x+2y=2,\ x,y\in\N\}$ 用列举法表示为 \blankline[4.4em].

\ans{$\{(0,1),(2,0)\}$}

\ifanswer
\solbegin
\step{由题意得 $x+2y=2$ 可化为 $y=1-\dfrac{x}{2}$，}
\step{由 $y\in\N$，有 $y=1-\dfrac{x}{2}\geqslant 0$，解得 $x\leqslant 2$，}
\step{又由 $x\in\N$，得 $x$ 可能取值为 $0$，$1$，$2$，}
\step{当 $x=0$ 时，得 $y=1$，满足条件，}
\step{当 $x=1$ 时，得 $y=\dfrac{1}{2}\notin\N$，不满足条件，}
\step{当 $x=2$ 时，得 $y=0$，满足条件，}
\step{综上，集合用列举法表示为 $\{(0,1),(2,0)\}$.}
\solend

\fi

\item 已知集合 $A=\{1\}$，集合 $B=\{0,1,2\}$，则满足关系 $A\sub P\sub B$ 的所有集合 $P$ 为 \blankline[4.4em].

\ans{$\{1\},\{0,1\},\{1,2\},\{0,1,2\}$}

\ifanswer
\solbegin
\step{由题意，集合 $P$ 可以是 $\{1\},\{0,1\},\{1,2\},\{0,1,2\}$.}
\solend

\fi

\item 已知集合 $A=\{x\mid a-1\leqslant x\leqslant a+2\}$，$B=\{x\mid 3<x<5\}$，则使 $A\supseteq B$ 成立的实数 $a$ 的取值范围是 \blankline[4.4em].

\ans{$[3,4]$}

\ifanswer
\solbegin
\step{因为 $A\supseteq B$，所以 $\begin{cases}a+2\geqslant 5\\ a-1\leqslant 3\end{cases}$，解得 $3\leqslant a\leqslant 4$.}
\solend

\fi

\item 若集合 $\{x\mid mx^2-4x+1=0\}=\{n\}$，则 $m+n=$ \blankline[2.4em].

\ans{$\dfrac{1}{4}$ 或 $\dfrac{9}{2}$}

\ifanswer
\solbegin
\step{若集合 $\{x\mid mx^2-4x+1=0\}=\{n\}$，则 $mx^2-4x+1=0$ 有唯一实数根 $n$，}
\step{当 $m=0$ 时，得 $n=\dfrac{1}{4}$，$m+n=\dfrac{1}{4}$，}
\step{当 $m\neq 0$ 时，$\Delta=16-4m=0$，即 $m=4$，}
\step{此时 $4x^2-4x+1=0$ 的根为 $n=\dfrac{1}{2}$，$m+n=\dfrac{9}{2}$.}
\step{故 $m+n=\dfrac{1}{4}$ 或 $\dfrac{9}{2}$.}
\solend

\fi

\item 集合 $M=\{-1,2,3,5,6\}$ 的全部非空子集的元素和等于 \blankline[4.4em].

\ans{$240$}

\ifanswer
\solbegin
\step{集合 $M=\{-1,2,3,5,6\}$ 中共有 $5$ 个元素，}
\step{在其所有子集中，元素 $-1$ 出现的次数为 $2^{5-1}=2^4=16$，}
\step{同理 $2$、$3$、$5$、$6$ 都出现 $16$ 次，}
\step{则集合 $M$ 的全部非空子集的元素和为 $(-1+2+3+5+6)\times 16=240$.}
\solend

\fi

\item 设 $A=\{x\mid x^2-3x-4=0\}$，$B=\{x\mid ax-1=0\}$，若 $A\cap B=B$，则实数 $a$ 的值可以为 \blankline[4.4em].

\ans{$0$、$-1$、$\dfrac{1}{4}$}

\ifanswer
\solbegin
\step{$A=\{-1,4\}$，因为 $A\cap B=B$，所以 $B\sub A$，}
\step{①$a=0$ 时，$B=\e$，满足 $B\sub A$；}
\step{②$a\neq 0$ 时，$B=\left\{\dfrac{1}{a}\right\}$，则 $\dfrac{1}{a}=-1$ 或 $\dfrac{1}{a}=4$，所以 $a=-1$ 或 $\dfrac{1}{4}$；}
\step{所以 $a$ 的值可以为 $0$、$-1$、$\dfrac{1}{4}$.}
\solend

\fi

\item 设集合 $A=\{1,m\}$，其中 $m$ 为实数. 令集合 $B=\{a+b\mid a\in A,b\in A\}$，若 $A\cap B$ 恰有一个元素，则 $A\cup B$ 的元素之和为 \blankline[2.4em].

\ans{$5$ 或 $10$}

\ifanswer
\solbegin
% 注：原卷本句印作“这个元不是 $1$ 就是 $m$”——“元素”漏了“素”字。按 高一b 第 13 题
%     “原卷漏排、本稿补上”的先例，本稿补出“素”字，其余照原卷，见文件头“原卷笔误”。
\step{由题意得，若 $A\cap B$ 恰有一个元素，这个元素不是 $1$ 就是 $m$，}
\stepm{（1）}{当 $1\in B$ 时，}
\step{①若 $1=1+m$，则 $m=0$，此时 $A=\{1,0\},B=\{2,1,0\},A\cap B=\{0,1\}$，不满足题意，舍去；}
\step{②若 $1=2m$，则 $m=\dfrac{1}{2}$，此时 $A=\left\{1,\dfrac{1}{2}\right\}$，$B=\left\{2,\dfrac{3}{2},1\right\}$，$A\cap B=\{1\}$，}
\step{满足题意，所以 $A\cup B=\left\{1,\dfrac{1}{2},2,\dfrac{3}{2}\right\}$，其元素之和为 $5$，}
\stepm{（2）}{当 $m\in B$ 时，}
\step{①若 $m=1+m$，则 $m$ 无解，}
\step{②若 $m=2m$，则 $m=0$，此时 $A=\{1,0\},B=\{2,1,0\},A\cap B=\{0,1\}$，不满足题意，}
\step{③若 $m=2$，则 $A=\{1,2\},B=\{2,3,4\},A\cap B=\{2\}$，满足题意，}
\step{所以 $A\cup B=\{1,2,3,4\}$，其元素之和为 $10$，}
\step{综上所述，$A\cup B$ 的元素之和为 $5$ 或 $10$.}
\solend

\fi

\item 设非空集合 $Q\sub M$，当 $Q$ 中所有元素和为偶数时（集合为单元素时和为元素本身），称 $Q$ 是 $M$ 的偶子集. 若集合 $M=\{1,2,3,4,5,6,7\}$，则其偶子集 $Q$ 的个数为 \blankline[2.4em].

\ans{$63$}

\ifanswer
\solbegin
\step{偶子集可以分为只有奇数、偶数及偶数个奇数和偶数的集合，}
\step{奇数有 $4$ 个，偶数有 $3$ 个，}
\step{只有奇数的有 $\mathrm{C}_4^2+1=7$ 个，只有偶数的有 $\mathrm{C}_3^1+\mathrm{C}_3^2+1=7$ 个，}
\step{偶数个奇数和偶数的：两奇一偶 $\mathrm{C}_4^2\cdot\mathrm{C}_3^1=18$ 个，两奇二偶 $\mathrm{C}_4^2\cdot\mathrm{C}_3^2=18$ 个，}
\step{两奇三偶 $\mathrm{C}_4^2\cdot\mathrm{C}_3^3=6$ 个，四奇一偶 $\mathrm{C}_4^4\cdot\mathrm{C}_3^1=3$ 个，四奇二偶 $\mathrm{C}_4^4\cdot\mathrm{C}_3^2=3$ 个，}
\step{四奇三偶 $\mathrm{C}_4^4\cdot\mathrm{C}_3^3=1$ 个，所以共有 $7+7+18+18+6+3+3+1=63$ 个.}
\solend

\fi

\item 已知非空集合 $S=\left\{x\mid -\dfrac{1}{2}\leqslant x\leqslant m\right\}$ 满足：当 $x\in S$ 时，有 $x^2\in S$，则实数 $m$ 的取值范围是 \blankline[4.4em].

\ans{$\left[\dfrac{1}{4},1\right]$}

\ifanswer
\solbegin
\step{因为集合 $S=\left\{x\mid -\dfrac{1}{2}\leqslant x\leqslant m\right\}$ 是非空集合，所以 $m\geqslant -\dfrac{1}{2}$，}
\step{又当 $x\in S$ 时，有 $x^2\in S$，所以 $m^2\leqslant m$ 且 $m\geqslant \left(-\dfrac{1}{2}\right)^2=\dfrac{1}{4}$，所以 $\dfrac{1}{4}\leqslant m\leqslant 1$.}
\solend

\fi

\end{enumerate}

\bigsec{二、选择题}

\begin{enumerate}
\setcounter{enumi}{12}

\item 已知集合 $U=\{1+x,1\}$，$x\in U$，则 $x=$\mbox{（\quad）}

\keepwithprev
A. $-1$\qquad B. $0$\qquad C. $1$\qquad D. $2$

\ans{C}

\ifanswer
\solbegin
\step{因为集合 $U=\{1+x,1\}$，$x\in U$，所以 $1+x=x$ 或 $x=1$，}
\step{当 $x=1+x$ 时，即 $0=1$，此时显然不成立，舍去，}
\step{当 $x=1$，此时集合 $U=\{1,2\}$，满足互异性，}
\step{综上所述，$x=1$. 故选 C.}
\solend

\fi

\item 已知集合 $M=\{m\mid m=a+b\sqrt2,\ a,b\in\Q\}$，则下列四个元素中属于 $M$ 的元素的个数是\mbox{（\quad）}

①$1+\sqrt2\pi$；②$\sqrt{11+6\sqrt2}$；③$\dfrac{1}{2+\sqrt2}$；④$\sqrt{2-\sqrt3}+\sqrt{2+\sqrt3}$

\keepwithprev
A. $4$\qquad B. $3$\qquad C. $2$\qquad D. $1$

\ans{C}

\ifanswer
\solbegin
\step{①$1+\sqrt2\pi$，由于 $\sqrt2$、$\pi$ 都为无理数，因此 $1+\sqrt2\pi\notin M$；}
\step{②$\sqrt{11+6\sqrt2}=\sqrt{11+2\sqrt{18}}=3+\sqrt2\in M$；}
\step{③$\dfrac{1}{2+\sqrt2}=1-\dfrac{1}{2}\sqrt2\in M$；}
\step{④$\sqrt{2-\sqrt3}+\sqrt{2+\sqrt3}=\sqrt{\dfrac{8-2\sqrt{12}}{4}}+\sqrt{\dfrac{8+2\sqrt{12}}{4}}=\dfrac{\sqrt6-\sqrt2}{2}+\dfrac{\sqrt6+\sqrt2}{2}=\sqrt6=\sqrt3\cdot\sqrt2\notin M$；}
\step{则上述四个元素中属于 $M$ 的元素的个数是 $2$. 故选 C.}
\solend

\fi

\item 集合 $A=\{2,3,7,8,9,11\}$，集合 $B=\{1,2,4,6,8\}$，下面说法正确的是\mbox{（\quad）}

\keepwithprev
A. 两个集合的交集是 $\{2,8\}$

B. 两个集合没有并集

C. $B$ 集合的元素无法用 $\{x\mid x=f(k),k\in D\}$ 这种形式表示出来

D. 集合 $A$ 有 $6$ 个子集

\ans{A}

\ifanswer
\solbegin
\step{对于 A，$A\cap B=\{2,8\}$，故 A 正确；}
\step{对于 B，$A\cup B=\{1,2,3,4,6,7,8,9,11\}$，故 B 错误；}
\step{对于 C，集合 $B$ 可以表示为 $B=\{x\mid x=f(k),k\in\{1,2,3,4,5\}\}$，}
\step{其中 $f(1)=1,f(2)=2,f(3)=4,f(4)=6,f(5)=8$，}
\step{即集合 $B$ 可描述为 $\{x\mid x=f(k),k\in D\}$ 形式，故 C 错误；}
\step{对于 D，集合 $A$ 有 $6$ 个元素，其子集总数为 $2^6=64\neq 6$，故 D 错误.}
\step{故选 A.}
\solend

\fi

\item 设 $M=\{1,2,\cdots,1995\}$，$A$ 是 $M$ 的子集且满足条件：当 $x\in A$ 时，$15x\notin A$，则 $A$ 中元素的个数的最大值为\mbox{（\quad）}

\keepwithprev
A. $1862$\qquad B. $1866$\qquad C. $1868$\qquad D. $1870$

\ans{D}

\ifanswer
\solbegin
\step{$1995=15\times 133$. 故取出所有不是 $15$ 的倍数的数，共 $1862$ 个，}
\step{这些数均符合要求.}
\step{在所有 $15$ 的倍数的数中，$15^2$ 的倍数有 $8$ 个，这些数又可以取出，}
\step{这样共取出了 $1870$ 个，即 $|A|\geqslant 1870$.}
\step{又 $\{k,15k\}(k=9,10,11,\cdots,133)$ 中的两个元素不能同时取出，}
\step{故 $|A|\leqslant 1995-133+8=1870$，故选 D.}
\solend

\fi

\end{enumerate}

\bigsec{三、解答题}

\begin{enumerate}
\setcounter{enumi}{16}

\item 已知集合 $A=\{x\mid a\leqslant x\leqslant a+4\}$，$B=\{x\mid x\geqslant 6\}$.

\begin{enumerate}[label=(\arabic*)]
\item 若 $a=3$，求 $A\cup B$；
\item 若 $A\cap B=A$，求实数 $a$ 的取值范围.
\end{enumerate}

\ans{（1）$\{x\mid x\geqslant 3\}$；（2）$[6,+\infty)$}

\ifanswer
\solbegin
\stepm{（1）}{当 $a=3$ 时，$A=\{x\mid 3\leqslant x\leqslant 7\}$，所以 $A\cup B=\{x\mid x\geqslant 3\}$；}
\stepm{（2）}{若 $A\cap B=A$，则 $A\sub B$，所以 $a\geqslant 6$，所以实数 $a$ 的取值范围是 $[6,+\infty)$.}
\solend

\fi

\ansspace[7]

\item 已知集合 $M=\{x\mid (x-a)(x^2-ax+a-1)=0\}$.

\begin{enumerate}[label=(\arabic*)]
\item 若 $a=3$，求集合 $M$；
\item 若集合 $M$ 中各元素之和等于 $3$，求实数 $a$ 的值，并用列举法表示集合 $M$.
\end{enumerate}

\ans{（1）$\{1,2,3\}$；（2）$a=2$ 或 $\dfrac{3}{2}$，此时 $M=\{1,2\}$ 或 $M=\left\{\dfrac{1}{2},1,\dfrac{3}{2}\right\}$}

\ifanswer
\solbegin
\stepm{（1）}{$\{1,2,3\}$；}
\stepm{（2）}{当 $(x-a)(x^2-ax+a-1)=0$ 有重根时，重根只能算作集合的一个元素，}
\step{$(x-a)(x^2-ax+a-1)=(x-a)(x-1)[x-(a-1)]=0$，}
\step{当 $a=1$ 时，得 $M=\{0,1\}$，不符合题意；}
\step{当 $a-1=1$ 时，即 $a=2$ 时，得 $M=\{1,2\}$，符合题意；}
\step{当 $a\neq 1$ 且 $a\neq 2$ 时，此时 $M=\{1,a,a-1\}$，得 $1+a+a-1=3$，}
\step{解得 $a=\dfrac{3}{2}$，此时 $M=\left\{\dfrac{1}{2},1,\dfrac{3}{2}\right\}$，符合题意，}
\step{综上，实数 $a$ 的值为 $2$ 或 $\dfrac{3}{2}$，}
\step{当 $a=2$ 时，$M=\{1,2\}$；当 $a=\dfrac{3}{2}$ 时，$M=\left\{\dfrac{1}{2},1,\dfrac{3}{2}\right\}$.}
\solend

\fi

\ansspace[8]

\item 已知集合 $A=\{x\mid x^2-2x-3=0\}$，集合 $B=\{x\mid ax-1=0\}$. 若 $B\sub A$，求实数 $a$ 的值.

\ans{$a=0$ 或 $a=-1$ 或 $a=\dfrac{1}{3}$}

\ifanswer
\solbegin
\step{当 $B=\e$ 时，$a=0$，满足题意；}
\step{当 $B\neq\e$，即 $a\neq 0$ 时，$B=\left\{\dfrac{1}{a}\right\}$，}
\step{又 $A=\{x\mid x^2-2x-3=0\}=\{x\mid x=-1\text{ 或 }x=3\}$，$B\sub A$，}
\step{所以 $\dfrac{1}{a}=-1$ 或 $\dfrac{1}{a}=3$，所以 $a=-1$ 或 $a=\dfrac{1}{3}$.}
\step{综上所述，$a=0$ 或 $a=-1$ 或 $a=\dfrac{1}{3}$.}
\solend

\fi

\ansspace[8]

\item 设集合 $A=\{x\mid -2\leqslant x\leqslant 5\}$，$B=\{x\mid m+1\leqslant x\leqslant 2m-1\}$.

\begin{enumerate}[label=(\arabic*)]
\item 若 $B\sub A$，求实数 $m$ 的取值范围；
\item 当 $x\in\Z$ 时，求 $A$ 的非空真子集个数；
\item 当 $x\in\R$ 时，不存在元素 $x$ 使 $x\in A$ 与 $x\in B$ 同时成立，求实数 $m$ 的取值范围.
\end{enumerate}

\ans{（1）$\{m\mid m\leqslant 3\}$；（2）$254$；（3）$\{m\mid m<2\text{ 或 }m>4\}$}

\ifanswer
\solbegin
\stepm{（1）}{当 $m+1>2m-1$，即 $m<2$ 时，$B=\e$，满足 $B\sub A$.}
\step{当 $m+1\leqslant 2m-1$，即 $m\geqslant 2$ 时，要使 $B\sub A$ 成立，}
\step{需 $\begin{cases}m+1\geqslant -2\\ 2m-1\leqslant 5\end{cases}$，得 $2\leqslant m\leqslant 3$，}
\step{综上，$m$ 的范围为 $\{m\mid m\leqslant 3\}$；}
\stepm{（2）}{$A=\{x\mid -2\leqslant x\leqslant 5,x\in\Z\}=\{-2,-1,0,1,2,3,4,5\}$，}
\step{则 $A$ 的非空真子集个数为 $2^8-1-1=254$；}
\stepm{（3）}{因为没有元素 $x$ 使 $x\in A$ 与 $x\in B$ 同时成立，所以 $A$ 与 $B$ 交集为空集.}
\step{①若 $B=\e$，即 $m+1>2m-1$，得 $m<2$ 时满足条件；}
\step{②若 $B\neq\e$，则要满足的条件是 $\begin{cases}m+1\leqslant 2m-1\\ m+1>5\end{cases}$ 或 $\begin{cases}m+1\leqslant 2m-1\\ 2m-1<-2\end{cases}$，}
\step{解得 $m>4$.}
\step{综上，有 $m<2$ 或 $m>4$，所以 $m$ 的范围为 $\{m\mid m<2\text{ 或 }m>4\}$.}
\solend

\fi

\ansspace[10]

% 压轴题：其后已无内容，故作答区全用可丢弃胶水（\ansspace*），并在题目前先量一下版面。
\ansneed{7cm}

\item 已知集合 $A=\{x\mid ax=1\}$，$B=\{x\mid x^2-ax+b=0\}$.

\begin{enumerate}[label=(\arabic*)]
\item 若 $a=2$，且 $A\sub B$，求 $b$ 的值；
\item 当 $a^2\geqslant 4b$ 时，若 $B\sub A$，求 $a$、$b$ 的值；
\item 若 $A\sub B$，讨论 $a$、$b$ 的取值范围.
\end{enumerate}

\ans{（1）$\dfrac{3}{4}$；（2）$a=\pm\sqrt2$，$b=\dfrac{1}{2}$；（3）$a=0$ 时 $b\in\R$，$a\neq 0$ 时 $b\in(-\infty,1)$}

\ifanswer
\solbegin
\stepm{（1）}{当 $a=2$ 时，$A=\left\{\dfrac{1}{2}\right\}$，}
\step{由 $A\sub B$ 得 $\dfrac{1}{2}$ 是方程 $x^2-2x+b=0$ 的解，代入得 $\left(\dfrac{1}{2}\right)^2-2\times\dfrac{1}{2}+b=0$，}
\step{解得 $b=\dfrac{3}{4}$.}
\stepm{（2）}{当 $a^2\geqslant 4b$ 时，方程 $x^2-ax+b=0$ 有实根.}
\step{若 $a=0$，则 $A=\e$，要 $B\sub A$ 需 $B=\e$，}
\step{但此时方程 $x^2+b=0$ 在 $a^2\geqslant 4b$（即 $b\leqslant 0$）时有实根，矛盾，舍去.}
\step{若 $a\neq 0$，则 $A=\left\{\dfrac{1}{a}\right\}$，$B\sub A$ 且方程有实根，故方程有且仅有一个根 $\dfrac{1}{a}$.}
\step{由韦达定理，根的和为 $a=\dfrac{1}{a}+\dfrac{1}{a}$，根的积为 $b=\dfrac{1}{a}\times\dfrac{1}{a}$.}
\step{由根的和得 $a=\dfrac{2}{a}$，即 $a^2=2$，$a=\pm\sqrt2$，$b=\dfrac{1}{2}$.}
\stepm{（3）}{当 $a=0$ 时，$A=\e$，故 $A\sub B$ 恒成立，此时 $b\in\R$.}
\step{当 $a\neq 0$ 时，$A=\left\{\dfrac{1}{a}\right\}$，将 $x=\dfrac{1}{a}$ 代入 $x^2-ax+b=0$，得 $b=1-\dfrac{1}{a^2}$.}
\step{此时方程判别式 $\Delta=a^2-4b=\dfrac{(a^2-2)^2}{a^2}\geqslant 0$，}
\step{若 $\Delta=0$，则 $a=\pm\sqrt2$，$b=\dfrac{1}{2}$；}
\step{若 $\Delta>0$，则 $a\neq 0$ 且 $a\neq\pm\sqrt2$，此时 $b=1-\dfrac{1}{a^2}<1$.}
\step{综上，当 $a=0$ 时，$b\in\R$；当 $a\neq 0$ 时，$b\in(-\infty,1)$.}
\solend

\fi

\ansspace*[12]

\end{enumerate}

\end{document}
"
% 紧凑版：xelatex "\def\iscompact{1}% 高一10_复附浦东_周末作业1.tex —— 高一上数学“2026-2027 年复附浦东高一上周末作业 1”（试卷版/答案版）
% 试卷版：xelatex 高一10_复附浦东_周末作业1.tex
% 答案版：xelatex "\def\isanswer{1}% 高一10_复附浦东_周末作业1.tex —— 高一上数学“2026-2027 年复附浦东高一上周末作业 1”（试卷版/答案版）
% 试卷版：xelatex 高一10_复附浦东_周末作业1.tex
% 答案版：xelatex "\def\isanswer{1}\input{高一10_复附浦东_周末作业1.tex}"
% 紧凑版：xelatex "\def\iscompact{1}\input{高一10_复附浦东_周末作业1.tex}"
%
% 原始材料是微信公众号图文（公众号“魔都数学加油站”），正文为 9 张 PNG 页；原文本身就是一份
% **讲评版**——题目下方直接印出【答案】或【解析】。试题文字、答案与解析均照原卷录入，未改内容。
%
% 与原卷的排版差异（均已在 README“说明”中交代）：
%   ① 原文自**第 3 题**起（第 1、2 题未在该文中发布），故本稿题号亦自 3 起；
%   ② 原卷本就印有“一、填空题”“二、选择题”“三、解答题”三节标题（分别在原图 01/04/06.png
%      上，逐页放大核对过），本稿照原卷分节，只把标题改排成全稿统一的“大节标题”样式；
%   ③ 原卷填空、选择的答案直接印在题目下方，本稿统一改为试卷版隐去、答案版以【答案】给出，
%      使两版彻底分离；
%   ④ 原卷的实数集符号时作正体 $R$、时作粗体 $\mathbf{R}$（第 21(3) 与 (3) 末句均为粗体），
%      本稿按全稿惯例统一写作 $\R$.
%
% 原卷笔误一处（本稿已补，见该处上方的 `%` 注）：
%   ① 第 10 题【解析】首句原卷印作“这个元不是 $1$ 就是 $m$”，“元素”漏了“素”字——
%      本稿按 高一b 第 13 题“原卷漏排、本稿补上”的先例补出“素”字，其余照原卷。
% 除此一处外，本卷 9 页逐页放大核对，未发现别的笔误（题 16 的 $1995=15\times 133$ 与
% “$15^2$ 的倍数有 8 个”两处最易被误判为笔误，实测 $133$、$8$ 均正确）。

% 本篇在公众号“魔都数学加油站”连载里的编号是“27学年高一10”，本地编号与之对齐（高一10）；
% 对应关系见 README“与公众号连载号的对照表”。
\documentclass[a4paper,11pt]{article}
\usepackage{common}

\begin{document}

\examtitlest[3]{2026--2027 年复附浦东高一上周末作业 1}{集合与常用逻辑用语}

\bigsec{一、填空题}

\begin{enumerate}
\setcounter{enumi}{2}

\item 已知集合 $M=\{1,2m+1\}$，$N=\{-1,m^2\}$，且 $M=N$，则 $m=$ \blankline[2.4em].

\ans{$-1$}

\ifanswer
\solbegin
\step{若 $M=N$，必有 $\begin{cases}2m+1=-1\\ m^2=1\end{cases}$，解得 $m=-1$.}
\solend

\fi

\item 集合 $\{(x,y)\mid x+2y=2,\ x,y\in\N\}$ 用列举法表示为 \blankline[4.4em].

\ans{$\{(0,1),(2,0)\}$}

\ifanswer
\solbegin
\step{由题意得 $x+2y=2$ 可化为 $y=1-\dfrac{x}{2}$，}
\step{由 $y\in\N$，有 $y=1-\dfrac{x}{2}\geqslant 0$，解得 $x\leqslant 2$，}
\step{又由 $x\in\N$，得 $x$ 可能取值为 $0$，$1$，$2$，}
\step{当 $x=0$ 时，得 $y=1$，满足条件，}
\step{当 $x=1$ 时，得 $y=\dfrac{1}{2}\notin\N$，不满足条件，}
\step{当 $x=2$ 时，得 $y=0$，满足条件，}
\step{综上，集合用列举法表示为 $\{(0,1),(2,0)\}$.}
\solend

\fi

\item 已知集合 $A=\{1\}$，集合 $B=\{0,1,2\}$，则满足关系 $A\sub P\sub B$ 的所有集合 $P$ 为 \blankline[4.4em].

\ans{$\{1\},\{0,1\},\{1,2\},\{0,1,2\}$}

\ifanswer
\solbegin
\step{由题意，集合 $P$ 可以是 $\{1\},\{0,1\},\{1,2\},\{0,1,2\}$.}
\solend

\fi

\item 已知集合 $A=\{x\mid a-1\leqslant x\leqslant a+2\}$，$B=\{x\mid 3<x<5\}$，则使 $A\supseteq B$ 成立的实数 $a$ 的取值范围是 \blankline[4.4em].

\ans{$[3,4]$}

\ifanswer
\solbegin
\step{因为 $A\supseteq B$，所以 $\begin{cases}a+2\geqslant 5\\ a-1\leqslant 3\end{cases}$，解得 $3\leqslant a\leqslant 4$.}
\solend

\fi

\item 若集合 $\{x\mid mx^2-4x+1=0\}=\{n\}$，则 $m+n=$ \blankline[2.4em].

\ans{$\dfrac{1}{4}$ 或 $\dfrac{9}{2}$}

\ifanswer
\solbegin
\step{若集合 $\{x\mid mx^2-4x+1=0\}=\{n\}$，则 $mx^2-4x+1=0$ 有唯一实数根 $n$，}
\step{当 $m=0$ 时，得 $n=\dfrac{1}{4}$，$m+n=\dfrac{1}{4}$，}
\step{当 $m\neq 0$ 时，$\Delta=16-4m=0$，即 $m=4$，}
\step{此时 $4x^2-4x+1=0$ 的根为 $n=\dfrac{1}{2}$，$m+n=\dfrac{9}{2}$.}
\step{故 $m+n=\dfrac{1}{4}$ 或 $\dfrac{9}{2}$.}
\solend

\fi

\item 集合 $M=\{-1,2,3,5,6\}$ 的全部非空子集的元素和等于 \blankline[4.4em].

\ans{$240$}

\ifanswer
\solbegin
\step{集合 $M=\{-1,2,3,5,6\}$ 中共有 $5$ 个元素，}
\step{在其所有子集中，元素 $-1$ 出现的次数为 $2^{5-1}=2^4=16$，}
\step{同理 $2$、$3$、$5$、$6$ 都出现 $16$ 次，}
\step{则集合 $M$ 的全部非空子集的元素和为 $(-1+2+3+5+6)\times 16=240$.}
\solend

\fi

\item 设 $A=\{x\mid x^2-3x-4=0\}$，$B=\{x\mid ax-1=0\}$，若 $A\cap B=B$，则实数 $a$ 的值可以为 \blankline[4.4em].

\ans{$0$、$-1$、$\dfrac{1}{4}$}

\ifanswer
\solbegin
\step{$A=\{-1,4\}$，因为 $A\cap B=B$，所以 $B\sub A$，}
\step{①$a=0$ 时，$B=\e$，满足 $B\sub A$；}
\step{②$a\neq 0$ 时，$B=\left\{\dfrac{1}{a}\right\}$，则 $\dfrac{1}{a}=-1$ 或 $\dfrac{1}{a}=4$，所以 $a=-1$ 或 $\dfrac{1}{4}$；}
\step{所以 $a$ 的值可以为 $0$、$-1$、$\dfrac{1}{4}$.}
\solend

\fi

\item 设集合 $A=\{1,m\}$，其中 $m$ 为实数. 令集合 $B=\{a+b\mid a\in A,b\in A\}$，若 $A\cap B$ 恰有一个元素，则 $A\cup B$ 的元素之和为 \blankline[2.4em].

\ans{$5$ 或 $10$}

\ifanswer
\solbegin
% 注：原卷本句印作“这个元不是 $1$ 就是 $m$”——“元素”漏了“素”字。按 高一b 第 13 题
%     “原卷漏排、本稿补上”的先例，本稿补出“素”字，其余照原卷，见文件头“原卷笔误”。
\step{由题意得，若 $A\cap B$ 恰有一个元素，这个元素不是 $1$ 就是 $m$，}
\stepm{（1）}{当 $1\in B$ 时，}
\step{①若 $1=1+m$，则 $m=0$，此时 $A=\{1,0\},B=\{2,1,0\},A\cap B=\{0,1\}$，不满足题意，舍去；}
\step{②若 $1=2m$，则 $m=\dfrac{1}{2}$，此时 $A=\left\{1,\dfrac{1}{2}\right\}$，$B=\left\{2,\dfrac{3}{2},1\right\}$，$A\cap B=\{1\}$，}
\step{满足题意，所以 $A\cup B=\left\{1,\dfrac{1}{2},2,\dfrac{3}{2}\right\}$，其元素之和为 $5$，}
\stepm{（2）}{当 $m\in B$ 时，}
\step{①若 $m=1+m$，则 $m$ 无解，}
\step{②若 $m=2m$，则 $m=0$，此时 $A=\{1,0\},B=\{2,1,0\},A\cap B=\{0,1\}$，不满足题意，}
\step{③若 $m=2$，则 $A=\{1,2\},B=\{2,3,4\},A\cap B=\{2\}$，满足题意，}
\step{所以 $A\cup B=\{1,2,3,4\}$，其元素之和为 $10$，}
\step{综上所述，$A\cup B$ 的元素之和为 $5$ 或 $10$.}
\solend

\fi

\item 设非空集合 $Q\sub M$，当 $Q$ 中所有元素和为偶数时（集合为单元素时和为元素本身），称 $Q$ 是 $M$ 的偶子集. 若集合 $M=\{1,2,3,4,5,6,7\}$，则其偶子集 $Q$ 的个数为 \blankline[2.4em].

\ans{$63$}

\ifanswer
\solbegin
\step{偶子集可以分为只有奇数、偶数及偶数个奇数和偶数的集合，}
\step{奇数有 $4$ 个，偶数有 $3$ 个，}
\step{只有奇数的有 $\mathrm{C}_4^2+1=7$ 个，只有偶数的有 $\mathrm{C}_3^1+\mathrm{C}_3^2+1=7$ 个，}
\step{偶数个奇数和偶数的：两奇一偶 $\mathrm{C}_4^2\cdot\mathrm{C}_3^1=18$ 个，两奇二偶 $\mathrm{C}_4^2\cdot\mathrm{C}_3^2=18$ 个，}
\step{两奇三偶 $\mathrm{C}_4^2\cdot\mathrm{C}_3^3=6$ 个，四奇一偶 $\mathrm{C}_4^4\cdot\mathrm{C}_3^1=3$ 个，四奇二偶 $\mathrm{C}_4^4\cdot\mathrm{C}_3^2=3$ 个，}
\step{四奇三偶 $\mathrm{C}_4^4\cdot\mathrm{C}_3^3=1$ 个，所以共有 $7+7+18+18+6+3+3+1=63$ 个.}
\solend

\fi

\item 已知非空集合 $S=\left\{x\mid -\dfrac{1}{2}\leqslant x\leqslant m\right\}$ 满足：当 $x\in S$ 时，有 $x^2\in S$，则实数 $m$ 的取值范围是 \blankline[4.4em].

\ans{$\left[\dfrac{1}{4},1\right]$}

\ifanswer
\solbegin
\step{因为集合 $S=\left\{x\mid -\dfrac{1}{2}\leqslant x\leqslant m\right\}$ 是非空集合，所以 $m\geqslant -\dfrac{1}{2}$，}
\step{又当 $x\in S$ 时，有 $x^2\in S$，所以 $m^2\leqslant m$ 且 $m\geqslant \left(-\dfrac{1}{2}\right)^2=\dfrac{1}{4}$，所以 $\dfrac{1}{4}\leqslant m\leqslant 1$.}
\solend

\fi

\end{enumerate}

\bigsec{二、选择题}

\begin{enumerate}
\setcounter{enumi}{12}

\item 已知集合 $U=\{1+x,1\}$，$x\in U$，则 $x=$\mbox{（\quad）}

\keepwithprev
A. $-1$\qquad B. $0$\qquad C. $1$\qquad D. $2$

\ans{C}

\ifanswer
\solbegin
\step{因为集合 $U=\{1+x,1\}$，$x\in U$，所以 $1+x=x$ 或 $x=1$，}
\step{当 $x=1+x$ 时，即 $0=1$，此时显然不成立，舍去，}
\step{当 $x=1$，此时集合 $U=\{1,2\}$，满足互异性，}
\step{综上所述，$x=1$. 故选 C.}
\solend

\fi

\item 已知集合 $M=\{m\mid m=a+b\sqrt2,\ a,b\in\Q\}$，则下列四个元素中属于 $M$ 的元素的个数是\mbox{（\quad）}

①$1+\sqrt2\pi$；②$\sqrt{11+6\sqrt2}$；③$\dfrac{1}{2+\sqrt2}$；④$\sqrt{2-\sqrt3}+\sqrt{2+\sqrt3}$

\keepwithprev
A. $4$\qquad B. $3$\qquad C. $2$\qquad D. $1$

\ans{C}

\ifanswer
\solbegin
\step{①$1+\sqrt2\pi$，由于 $\sqrt2$、$\pi$ 都为无理数，因此 $1+\sqrt2\pi\notin M$；}
\step{②$\sqrt{11+6\sqrt2}=\sqrt{11+2\sqrt{18}}=3+\sqrt2\in M$；}
\step{③$\dfrac{1}{2+\sqrt2}=1-\dfrac{1}{2}\sqrt2\in M$；}
\step{④$\sqrt{2-\sqrt3}+\sqrt{2+\sqrt3}=\sqrt{\dfrac{8-2\sqrt{12}}{4}}+\sqrt{\dfrac{8+2\sqrt{12}}{4}}=\dfrac{\sqrt6-\sqrt2}{2}+\dfrac{\sqrt6+\sqrt2}{2}=\sqrt6=\sqrt3\cdot\sqrt2\notin M$；}
\step{则上述四个元素中属于 $M$ 的元素的个数是 $2$. 故选 C.}
\solend

\fi

\item 集合 $A=\{2,3,7,8,9,11\}$，集合 $B=\{1,2,4,6,8\}$，下面说法正确的是\mbox{（\quad）}

\keepwithprev
A. 两个集合的交集是 $\{2,8\}$

B. 两个集合没有并集

C. $B$ 集合的元素无法用 $\{x\mid x=f(k),k\in D\}$ 这种形式表示出来

D. 集合 $A$ 有 $6$ 个子集

\ans{A}

\ifanswer
\solbegin
\step{对于 A，$A\cap B=\{2,8\}$，故 A 正确；}
\step{对于 B，$A\cup B=\{1,2,3,4,6,7,8,9,11\}$，故 B 错误；}
\step{对于 C，集合 $B$ 可以表示为 $B=\{x\mid x=f(k),k\in\{1,2,3,4,5\}\}$，}
\step{其中 $f(1)=1,f(2)=2,f(3)=4,f(4)=6,f(5)=8$，}
\step{即集合 $B$ 可描述为 $\{x\mid x=f(k),k\in D\}$ 形式，故 C 错误；}
\step{对于 D，集合 $A$ 有 $6$ 个元素，其子集总数为 $2^6=64\neq 6$，故 D 错误.}
\step{故选 A.}
\solend

\fi

\item 设 $M=\{1,2,\cdots,1995\}$，$A$ 是 $M$ 的子集且满足条件：当 $x\in A$ 时，$15x\notin A$，则 $A$ 中元素的个数的最大值为\mbox{（\quad）}

\keepwithprev
A. $1862$\qquad B. $1866$\qquad C. $1868$\qquad D. $1870$

\ans{D}

\ifanswer
\solbegin
\step{$1995=15\times 133$. 故取出所有不是 $15$ 的倍数的数，共 $1862$ 个，}
\step{这些数均符合要求.}
\step{在所有 $15$ 的倍数的数中，$15^2$ 的倍数有 $8$ 个，这些数又可以取出，}
\step{这样共取出了 $1870$ 个，即 $|A|\geqslant 1870$.}
\step{又 $\{k,15k\}(k=9,10,11,\cdots,133)$ 中的两个元素不能同时取出，}
\step{故 $|A|\leqslant 1995-133+8=1870$，故选 D.}
\solend

\fi

\end{enumerate}

\bigsec{三、解答题}

\begin{enumerate}
\setcounter{enumi}{16}

\item 已知集合 $A=\{x\mid a\leqslant x\leqslant a+4\}$，$B=\{x\mid x\geqslant 6\}$.

\begin{enumerate}[label=(\arabic*)]
\item 若 $a=3$，求 $A\cup B$；
\item 若 $A\cap B=A$，求实数 $a$ 的取值范围.
\end{enumerate}

\ans{（1）$\{x\mid x\geqslant 3\}$；（2）$[6,+\infty)$}

\ifanswer
\solbegin
\stepm{（1）}{当 $a=3$ 时，$A=\{x\mid 3\leqslant x\leqslant 7\}$，所以 $A\cup B=\{x\mid x\geqslant 3\}$；}
\stepm{（2）}{若 $A\cap B=A$，则 $A\sub B$，所以 $a\geqslant 6$，所以实数 $a$ 的取值范围是 $[6,+\infty)$.}
\solend

\fi

\ansspace[7]

\item 已知集合 $M=\{x\mid (x-a)(x^2-ax+a-1)=0\}$.

\begin{enumerate}[label=(\arabic*)]
\item 若 $a=3$，求集合 $M$；
\item 若集合 $M$ 中各元素之和等于 $3$，求实数 $a$ 的值，并用列举法表示集合 $M$.
\end{enumerate}

\ans{（1）$\{1,2,3\}$；（2）$a=2$ 或 $\dfrac{3}{2}$，此时 $M=\{1,2\}$ 或 $M=\left\{\dfrac{1}{2},1,\dfrac{3}{2}\right\}$}

\ifanswer
\solbegin
\stepm{（1）}{$\{1,2,3\}$；}
\stepm{（2）}{当 $(x-a)(x^2-ax+a-1)=0$ 有重根时，重根只能算作集合的一个元素，}
\step{$(x-a)(x^2-ax+a-1)=(x-a)(x-1)[x-(a-1)]=0$，}
\step{当 $a=1$ 时，得 $M=\{0,1\}$，不符合题意；}
\step{当 $a-1=1$ 时，即 $a=2$ 时，得 $M=\{1,2\}$，符合题意；}
\step{当 $a\neq 1$ 且 $a\neq 2$ 时，此时 $M=\{1,a,a-1\}$，得 $1+a+a-1=3$，}
\step{解得 $a=\dfrac{3}{2}$，此时 $M=\left\{\dfrac{1}{2},1,\dfrac{3}{2}\right\}$，符合题意，}
\step{综上，实数 $a$ 的值为 $2$ 或 $\dfrac{3}{2}$，}
\step{当 $a=2$ 时，$M=\{1,2\}$；当 $a=\dfrac{3}{2}$ 时，$M=\left\{\dfrac{1}{2},1,\dfrac{3}{2}\right\}$.}
\solend

\fi

\ansspace[8]

\item 已知集合 $A=\{x\mid x^2-2x-3=0\}$，集合 $B=\{x\mid ax-1=0\}$. 若 $B\sub A$，求实数 $a$ 的值.

\ans{$a=0$ 或 $a=-1$ 或 $a=\dfrac{1}{3}$}

\ifanswer
\solbegin
\step{当 $B=\e$ 时，$a=0$，满足题意；}
\step{当 $B\neq\e$，即 $a\neq 0$ 时，$B=\left\{\dfrac{1}{a}\right\}$，}
\step{又 $A=\{x\mid x^2-2x-3=0\}=\{x\mid x=-1\text{ 或 }x=3\}$，$B\sub A$，}
\step{所以 $\dfrac{1}{a}=-1$ 或 $\dfrac{1}{a}=3$，所以 $a=-1$ 或 $a=\dfrac{1}{3}$.}
\step{综上所述，$a=0$ 或 $a=-1$ 或 $a=\dfrac{1}{3}$.}
\solend

\fi

\ansspace[8]

\item 设集合 $A=\{x\mid -2\leqslant x\leqslant 5\}$，$B=\{x\mid m+1\leqslant x\leqslant 2m-1\}$.

\begin{enumerate}[label=(\arabic*)]
\item 若 $B\sub A$，求实数 $m$ 的取值范围；
\item 当 $x\in\Z$ 时，求 $A$ 的非空真子集个数；
\item 当 $x\in\R$ 时，不存在元素 $x$ 使 $x\in A$ 与 $x\in B$ 同时成立，求实数 $m$ 的取值范围.
\end{enumerate}

\ans{（1）$\{m\mid m\leqslant 3\}$；（2）$254$；（3）$\{m\mid m<2\text{ 或 }m>4\}$}

\ifanswer
\solbegin
\stepm{（1）}{当 $m+1>2m-1$，即 $m<2$ 时，$B=\e$，满足 $B\sub A$.}
\step{当 $m+1\leqslant 2m-1$，即 $m\geqslant 2$ 时，要使 $B\sub A$ 成立，}
\step{需 $\begin{cases}m+1\geqslant -2\\ 2m-1\leqslant 5\end{cases}$，得 $2\leqslant m\leqslant 3$，}
\step{综上，$m$ 的范围为 $\{m\mid m\leqslant 3\}$；}
\stepm{（2）}{$A=\{x\mid -2\leqslant x\leqslant 5,x\in\Z\}=\{-2,-1,0,1,2,3,4,5\}$，}
\step{则 $A$ 的非空真子集个数为 $2^8-1-1=254$；}
\stepm{（3）}{因为没有元素 $x$ 使 $x\in A$ 与 $x\in B$ 同时成立，所以 $A$ 与 $B$ 交集为空集.}
\step{①若 $B=\e$，即 $m+1>2m-1$，得 $m<2$ 时满足条件；}
\step{②若 $B\neq\e$，则要满足的条件是 $\begin{cases}m+1\leqslant 2m-1\\ m+1>5\end{cases}$ 或 $\begin{cases}m+1\leqslant 2m-1\\ 2m-1<-2\end{cases}$，}
\step{解得 $m>4$.}
\step{综上，有 $m<2$ 或 $m>4$，所以 $m$ 的范围为 $\{m\mid m<2\text{ 或 }m>4\}$.}
\solend

\fi

\ansspace[10]

% 压轴题：其后已无内容，故作答区全用可丢弃胶水（\ansspace*），并在题目前先量一下版面。
\ansneed{7cm}

\item 已知集合 $A=\{x\mid ax=1\}$，$B=\{x\mid x^2-ax+b=0\}$.

\begin{enumerate}[label=(\arabic*)]
\item 若 $a=2$，且 $A\sub B$，求 $b$ 的值；
\item 当 $a^2\geqslant 4b$ 时，若 $B\sub A$，求 $a$、$b$ 的值；
\item 若 $A\sub B$，讨论 $a$、$b$ 的取值范围.
\end{enumerate}

\ans{（1）$\dfrac{3}{4}$；（2）$a=\pm\sqrt2$，$b=\dfrac{1}{2}$；（3）$a=0$ 时 $b\in\R$，$a\neq 0$ 时 $b\in(-\infty,1)$}

\ifanswer
\solbegin
\stepm{（1）}{当 $a=2$ 时，$A=\left\{\dfrac{1}{2}\right\}$，}
\step{由 $A\sub B$ 得 $\dfrac{1}{2}$ 是方程 $x^2-2x+b=0$ 的解，代入得 $\left(\dfrac{1}{2}\right)^2-2\times\dfrac{1}{2}+b=0$，}
\step{解得 $b=\dfrac{3}{4}$.}
\stepm{（2）}{当 $a^2\geqslant 4b$ 时，方程 $x^2-ax+b=0$ 有实根.}
\step{若 $a=0$，则 $A=\e$，要 $B\sub A$ 需 $B=\e$，}
\step{但此时方程 $x^2+b=0$ 在 $a^2\geqslant 4b$（即 $b\leqslant 0$）时有实根，矛盾，舍去.}
\step{若 $a\neq 0$，则 $A=\left\{\dfrac{1}{a}\right\}$，$B\sub A$ 且方程有实根，故方程有且仅有一个根 $\dfrac{1}{a}$.}
\step{由韦达定理，根的和为 $a=\dfrac{1}{a}+\dfrac{1}{a}$，根的积为 $b=\dfrac{1}{a}\times\dfrac{1}{a}$.}
\step{由根的和得 $a=\dfrac{2}{a}$，即 $a^2=2$，$a=\pm\sqrt2$，$b=\dfrac{1}{2}$.}
\stepm{（3）}{当 $a=0$ 时，$A=\e$，故 $A\sub B$ 恒成立，此时 $b\in\R$.}
\step{当 $a\neq 0$ 时，$A=\left\{\dfrac{1}{a}\right\}$，将 $x=\dfrac{1}{a}$ 代入 $x^2-ax+b=0$，得 $b=1-\dfrac{1}{a^2}$.}
\step{此时方程判别式 $\Delta=a^2-4b=\dfrac{(a^2-2)^2}{a^2}\geqslant 0$，}
\step{若 $\Delta=0$，则 $a=\pm\sqrt2$，$b=\dfrac{1}{2}$；}
\step{若 $\Delta>0$，则 $a\neq 0$ 且 $a\neq\pm\sqrt2$，此时 $b=1-\dfrac{1}{a^2}<1$.}
\step{综上，当 $a=0$ 时，$b\in\R$；当 $a\neq 0$ 时，$b\in(-\infty,1)$.}
\solend

\fi

\ansspace*[12]

\end{enumerate}

\end{document}
"
% 紧凑版：xelatex "\def\iscompact{1}% 高一10_复附浦东_周末作业1.tex —— 高一上数学“2026-2027 年复附浦东高一上周末作业 1”（试卷版/答案版）
% 试卷版：xelatex 高一10_复附浦东_周末作业1.tex
% 答案版：xelatex "\def\isanswer{1}\input{高一10_复附浦东_周末作业1.tex}"
% 紧凑版：xelatex "\def\iscompact{1}\input{高一10_复附浦东_周末作业1.tex}"
%
% 原始材料是微信公众号图文（公众号“魔都数学加油站”），正文为 9 张 PNG 页；原文本身就是一份
% **讲评版**——题目下方直接印出【答案】或【解析】。试题文字、答案与解析均照原卷录入，未改内容。
%
% 与原卷的排版差异（均已在 README“说明”中交代）：
%   ① 原文自**第 3 题**起（第 1、2 题未在该文中发布），故本稿题号亦自 3 起；
%   ② 原卷本就印有“一、填空题”“二、选择题”“三、解答题”三节标题（分别在原图 01/04/06.png
%      上，逐页放大核对过），本稿照原卷分节，只把标题改排成全稿统一的“大节标题”样式；
%   ③ 原卷填空、选择的答案直接印在题目下方，本稿统一改为试卷版隐去、答案版以【答案】给出，
%      使两版彻底分离；
%   ④ 原卷的实数集符号时作正体 $R$、时作粗体 $\mathbf{R}$（第 21(3) 与 (3) 末句均为粗体），
%      本稿按全稿惯例统一写作 $\R$.
%
% 原卷笔误一处（本稿已补，见该处上方的 `%` 注）：
%   ① 第 10 题【解析】首句原卷印作“这个元不是 $1$ 就是 $m$”，“元素”漏了“素”字——
%      本稿按 高一b 第 13 题“原卷漏排、本稿补上”的先例补出“素”字，其余照原卷。
% 除此一处外，本卷 9 页逐页放大核对，未发现别的笔误（题 16 的 $1995=15\times 133$ 与
% “$15^2$ 的倍数有 8 个”两处最易被误判为笔误，实测 $133$、$8$ 均正确）。

% 本篇在公众号“魔都数学加油站”连载里的编号是“27学年高一10”，本地编号与之对齐（高一10）；
% 对应关系见 README“与公众号连载号的对照表”。
\documentclass[a4paper,11pt]{article}
\usepackage{common}

\begin{document}

\examtitlest[3]{2026--2027 年复附浦东高一上周末作业 1}{集合与常用逻辑用语}

\bigsec{一、填空题}

\begin{enumerate}
\setcounter{enumi}{2}

\item 已知集合 $M=\{1,2m+1\}$，$N=\{-1,m^2\}$，且 $M=N$，则 $m=$ \blankline[2.4em].

\ans{$-1$}

\ifanswer
\solbegin
\step{若 $M=N$，必有 $\begin{cases}2m+1=-1\\ m^2=1\end{cases}$，解得 $m=-1$.}
\solend

\fi

\item 集合 $\{(x,y)\mid x+2y=2,\ x,y\in\N\}$ 用列举法表示为 \blankline[4.4em].

\ans{$\{(0,1),(2,0)\}$}

\ifanswer
\solbegin
\step{由题意得 $x+2y=2$ 可化为 $y=1-\dfrac{x}{2}$，}
\step{由 $y\in\N$，有 $y=1-\dfrac{x}{2}\geqslant 0$，解得 $x\leqslant 2$，}
\step{又由 $x\in\N$，得 $x$ 可能取值为 $0$，$1$，$2$，}
\step{当 $x=0$ 时，得 $y=1$，满足条件，}
\step{当 $x=1$ 时，得 $y=\dfrac{1}{2}\notin\N$，不满足条件，}
\step{当 $x=2$ 时，得 $y=0$，满足条件，}
\step{综上，集合用列举法表示为 $\{(0,1),(2,0)\}$.}
\solend

\fi

\item 已知集合 $A=\{1\}$，集合 $B=\{0,1,2\}$，则满足关系 $A\sub P\sub B$ 的所有集合 $P$ 为 \blankline[4.4em].

\ans{$\{1\},\{0,1\},\{1,2\},\{0,1,2\}$}

\ifanswer
\solbegin
\step{由题意，集合 $P$ 可以是 $\{1\},\{0,1\},\{1,2\},\{0,1,2\}$.}
\solend

\fi

\item 已知集合 $A=\{x\mid a-1\leqslant x\leqslant a+2\}$，$B=\{x\mid 3<x<5\}$，则使 $A\supseteq B$ 成立的实数 $a$ 的取值范围是 \blankline[4.4em].

\ans{$[3,4]$}

\ifanswer
\solbegin
\step{因为 $A\supseteq B$，所以 $\begin{cases}a+2\geqslant 5\\ a-1\leqslant 3\end{cases}$，解得 $3\leqslant a\leqslant 4$.}
\solend

\fi

\item 若集合 $\{x\mid mx^2-4x+1=0\}=\{n\}$，则 $m+n=$ \blankline[2.4em].

\ans{$\dfrac{1}{4}$ 或 $\dfrac{9}{2}$}

\ifanswer
\solbegin
\step{若集合 $\{x\mid mx^2-4x+1=0\}=\{n\}$，则 $mx^2-4x+1=0$ 有唯一实数根 $n$，}
\step{当 $m=0$ 时，得 $n=\dfrac{1}{4}$，$m+n=\dfrac{1}{4}$，}
\step{当 $m\neq 0$ 时，$\Delta=16-4m=0$，即 $m=4$，}
\step{此时 $4x^2-4x+1=0$ 的根为 $n=\dfrac{1}{2}$，$m+n=\dfrac{9}{2}$.}
\step{故 $m+n=\dfrac{1}{4}$ 或 $\dfrac{9}{2}$.}
\solend

\fi

\item 集合 $M=\{-1,2,3,5,6\}$ 的全部非空子集的元素和等于 \blankline[4.4em].

\ans{$240$}

\ifanswer
\solbegin
\step{集合 $M=\{-1,2,3,5,6\}$ 中共有 $5$ 个元素，}
\step{在其所有子集中，元素 $-1$ 出现的次数为 $2^{5-1}=2^4=16$，}
\step{同理 $2$、$3$、$5$、$6$ 都出现 $16$ 次，}
\step{则集合 $M$ 的全部非空子集的元素和为 $(-1+2+3+5+6)\times 16=240$.}
\solend

\fi

\item 设 $A=\{x\mid x^2-3x-4=0\}$，$B=\{x\mid ax-1=0\}$，若 $A\cap B=B$，则实数 $a$ 的值可以为 \blankline[4.4em].

\ans{$0$、$-1$、$\dfrac{1}{4}$}

\ifanswer
\solbegin
\step{$A=\{-1,4\}$，因为 $A\cap B=B$，所以 $B\sub A$，}
\step{①$a=0$ 时，$B=\e$，满足 $B\sub A$；}
\step{②$a\neq 0$ 时，$B=\left\{\dfrac{1}{a}\right\}$，则 $\dfrac{1}{a}=-1$ 或 $\dfrac{1}{a}=4$，所以 $a=-1$ 或 $\dfrac{1}{4}$；}
\step{所以 $a$ 的值可以为 $0$、$-1$、$\dfrac{1}{4}$.}
\solend

\fi

\item 设集合 $A=\{1,m\}$，其中 $m$ 为实数. 令集合 $B=\{a+b\mid a\in A,b\in A\}$，若 $A\cap B$ 恰有一个元素，则 $A\cup B$ 的元素之和为 \blankline[2.4em].

\ans{$5$ 或 $10$}

\ifanswer
\solbegin
% 注：原卷本句印作“这个元不是 $1$ 就是 $m$”——“元素”漏了“素”字。按 高一b 第 13 题
%     “原卷漏排、本稿补上”的先例，本稿补出“素”字，其余照原卷，见文件头“原卷笔误”。
\step{由题意得，若 $A\cap B$ 恰有一个元素，这个元素不是 $1$ 就是 $m$，}
\stepm{（1）}{当 $1\in B$ 时，}
\step{①若 $1=1+m$，则 $m=0$，此时 $A=\{1,0\},B=\{2,1,0\},A\cap B=\{0,1\}$，不满足题意，舍去；}
\step{②若 $1=2m$，则 $m=\dfrac{1}{2}$，此时 $A=\left\{1,\dfrac{1}{2}\right\}$，$B=\left\{2,\dfrac{3}{2},1\right\}$，$A\cap B=\{1\}$，}
\step{满足题意，所以 $A\cup B=\left\{1,\dfrac{1}{2},2,\dfrac{3}{2}\right\}$，其元素之和为 $5$，}
\stepm{（2）}{当 $m\in B$ 时，}
\step{①若 $m=1+m$，则 $m$ 无解，}
\step{②若 $m=2m$，则 $m=0$，此时 $A=\{1,0\},B=\{2,1,0\},A\cap B=\{0,1\}$，不满足题意，}
\step{③若 $m=2$，则 $A=\{1,2\},B=\{2,3,4\},A\cap B=\{2\}$，满足题意，}
\step{所以 $A\cup B=\{1,2,3,4\}$，其元素之和为 $10$，}
\step{综上所述，$A\cup B$ 的元素之和为 $5$ 或 $10$.}
\solend

\fi

\item 设非空集合 $Q\sub M$，当 $Q$ 中所有元素和为偶数时（集合为单元素时和为元素本身），称 $Q$ 是 $M$ 的偶子集. 若集合 $M=\{1,2,3,4,5,6,7\}$，则其偶子集 $Q$ 的个数为 \blankline[2.4em].

\ans{$63$}

\ifanswer
\solbegin
\step{偶子集可以分为只有奇数、偶数及偶数个奇数和偶数的集合，}
\step{奇数有 $4$ 个，偶数有 $3$ 个，}
\step{只有奇数的有 $\mathrm{C}_4^2+1=7$ 个，只有偶数的有 $\mathrm{C}_3^1+\mathrm{C}_3^2+1=7$ 个，}
\step{偶数个奇数和偶数的：两奇一偶 $\mathrm{C}_4^2\cdot\mathrm{C}_3^1=18$ 个，两奇二偶 $\mathrm{C}_4^2\cdot\mathrm{C}_3^2=18$ 个，}
\step{两奇三偶 $\mathrm{C}_4^2\cdot\mathrm{C}_3^3=6$ 个，四奇一偶 $\mathrm{C}_4^4\cdot\mathrm{C}_3^1=3$ 个，四奇二偶 $\mathrm{C}_4^4\cdot\mathrm{C}_3^2=3$ 个，}
\step{四奇三偶 $\mathrm{C}_4^4\cdot\mathrm{C}_3^3=1$ 个，所以共有 $7+7+18+18+6+3+3+1=63$ 个.}
\solend

\fi

\item 已知非空集合 $S=\left\{x\mid -\dfrac{1}{2}\leqslant x\leqslant m\right\}$ 满足：当 $x\in S$ 时，有 $x^2\in S$，则实数 $m$ 的取值范围是 \blankline[4.4em].

\ans{$\left[\dfrac{1}{4},1\right]$}

\ifanswer
\solbegin
\step{因为集合 $S=\left\{x\mid -\dfrac{1}{2}\leqslant x\leqslant m\right\}$ 是非空集合，所以 $m\geqslant -\dfrac{1}{2}$，}
\step{又当 $x\in S$ 时，有 $x^2\in S$，所以 $m^2\leqslant m$ 且 $m\geqslant \left(-\dfrac{1}{2}\right)^2=\dfrac{1}{4}$，所以 $\dfrac{1}{4}\leqslant m\leqslant 1$.}
\solend

\fi

\end{enumerate}

\bigsec{二、选择题}

\begin{enumerate}
\setcounter{enumi}{12}

\item 已知集合 $U=\{1+x,1\}$，$x\in U$，则 $x=$\mbox{（\quad）}

\keepwithprev
A. $-1$\qquad B. $0$\qquad C. $1$\qquad D. $2$

\ans{C}

\ifanswer
\solbegin
\step{因为集合 $U=\{1+x,1\}$，$x\in U$，所以 $1+x=x$ 或 $x=1$，}
\step{当 $x=1+x$ 时，即 $0=1$，此时显然不成立，舍去，}
\step{当 $x=1$，此时集合 $U=\{1,2\}$，满足互异性，}
\step{综上所述，$x=1$. 故选 C.}
\solend

\fi

\item 已知集合 $M=\{m\mid m=a+b\sqrt2,\ a,b\in\Q\}$，则下列四个元素中属于 $M$ 的元素的个数是\mbox{（\quad）}

①$1+\sqrt2\pi$；②$\sqrt{11+6\sqrt2}$；③$\dfrac{1}{2+\sqrt2}$；④$\sqrt{2-\sqrt3}+\sqrt{2+\sqrt3}$

\keepwithprev
A. $4$\qquad B. $3$\qquad C. $2$\qquad D. $1$

\ans{C}

\ifanswer
\solbegin
\step{①$1+\sqrt2\pi$，由于 $\sqrt2$、$\pi$ 都为无理数，因此 $1+\sqrt2\pi\notin M$；}
\step{②$\sqrt{11+6\sqrt2}=\sqrt{11+2\sqrt{18}}=3+\sqrt2\in M$；}
\step{③$\dfrac{1}{2+\sqrt2}=1-\dfrac{1}{2}\sqrt2\in M$；}
\step{④$\sqrt{2-\sqrt3}+\sqrt{2+\sqrt3}=\sqrt{\dfrac{8-2\sqrt{12}}{4}}+\sqrt{\dfrac{8+2\sqrt{12}}{4}}=\dfrac{\sqrt6-\sqrt2}{2}+\dfrac{\sqrt6+\sqrt2}{2}=\sqrt6=\sqrt3\cdot\sqrt2\notin M$；}
\step{则上述四个元素中属于 $M$ 的元素的个数是 $2$. 故选 C.}
\solend

\fi

\item 集合 $A=\{2,3,7,8,9,11\}$，集合 $B=\{1,2,4,6,8\}$，下面说法正确的是\mbox{（\quad）}

\keepwithprev
A. 两个集合的交集是 $\{2,8\}$

B. 两个集合没有并集

C. $B$ 集合的元素无法用 $\{x\mid x=f(k),k\in D\}$ 这种形式表示出来

D. 集合 $A$ 有 $6$ 个子集

\ans{A}

\ifanswer
\solbegin
\step{对于 A，$A\cap B=\{2,8\}$，故 A 正确；}
\step{对于 B，$A\cup B=\{1,2,3,4,6,7,8,9,11\}$，故 B 错误；}
\step{对于 C，集合 $B$ 可以表示为 $B=\{x\mid x=f(k),k\in\{1,2,3,4,5\}\}$，}
\step{其中 $f(1)=1,f(2)=2,f(3)=4,f(4)=6,f(5)=8$，}
\step{即集合 $B$ 可描述为 $\{x\mid x=f(k),k\in D\}$ 形式，故 C 错误；}
\step{对于 D，集合 $A$ 有 $6$ 个元素，其子集总数为 $2^6=64\neq 6$，故 D 错误.}
\step{故选 A.}
\solend

\fi

\item 设 $M=\{1,2,\cdots,1995\}$，$A$ 是 $M$ 的子集且满足条件：当 $x\in A$ 时，$15x\notin A$，则 $A$ 中元素的个数的最大值为\mbox{（\quad）}

\keepwithprev
A. $1862$\qquad B. $1866$\qquad C. $1868$\qquad D. $1870$

\ans{D}

\ifanswer
\solbegin
\step{$1995=15\times 133$. 故取出所有不是 $15$ 的倍数的数，共 $1862$ 个，}
\step{这些数均符合要求.}
\step{在所有 $15$ 的倍数的数中，$15^2$ 的倍数有 $8$ 个，这些数又可以取出，}
\step{这样共取出了 $1870$ 个，即 $|A|\geqslant 1870$.}
\step{又 $\{k,15k\}(k=9,10,11,\cdots,133)$ 中的两个元素不能同时取出，}
\step{故 $|A|\leqslant 1995-133+8=1870$，故选 D.}
\solend

\fi

\end{enumerate}

\bigsec{三、解答题}

\begin{enumerate}
\setcounter{enumi}{16}

\item 已知集合 $A=\{x\mid a\leqslant x\leqslant a+4\}$，$B=\{x\mid x\geqslant 6\}$.

\begin{enumerate}[label=(\arabic*)]
\item 若 $a=3$，求 $A\cup B$；
\item 若 $A\cap B=A$，求实数 $a$ 的取值范围.
\end{enumerate}

\ans{（1）$\{x\mid x\geqslant 3\}$；（2）$[6,+\infty)$}

\ifanswer
\solbegin
\stepm{（1）}{当 $a=3$ 时，$A=\{x\mid 3\leqslant x\leqslant 7\}$，所以 $A\cup B=\{x\mid x\geqslant 3\}$；}
\stepm{（2）}{若 $A\cap B=A$，则 $A\sub B$，所以 $a\geqslant 6$，所以实数 $a$ 的取值范围是 $[6,+\infty)$.}
\solend

\fi

\ansspace[7]

\item 已知集合 $M=\{x\mid (x-a)(x^2-ax+a-1)=0\}$.

\begin{enumerate}[label=(\arabic*)]
\item 若 $a=3$，求集合 $M$；
\item 若集合 $M$ 中各元素之和等于 $3$，求实数 $a$ 的值，并用列举法表示集合 $M$.
\end{enumerate}

\ans{（1）$\{1,2,3\}$；（2）$a=2$ 或 $\dfrac{3}{2}$，此时 $M=\{1,2\}$ 或 $M=\left\{\dfrac{1}{2},1,\dfrac{3}{2}\right\}$}

\ifanswer
\solbegin
\stepm{（1）}{$\{1,2,3\}$；}
\stepm{（2）}{当 $(x-a)(x^2-ax+a-1)=0$ 有重根时，重根只能算作集合的一个元素，}
\step{$(x-a)(x^2-ax+a-1)=(x-a)(x-1)[x-(a-1)]=0$，}
\step{当 $a=1$ 时，得 $M=\{0,1\}$，不符合题意；}
\step{当 $a-1=1$ 时，即 $a=2$ 时，得 $M=\{1,2\}$，符合题意；}
\step{当 $a\neq 1$ 且 $a\neq 2$ 时，此时 $M=\{1,a,a-1\}$，得 $1+a+a-1=3$，}
\step{解得 $a=\dfrac{3}{2}$，此时 $M=\left\{\dfrac{1}{2},1,\dfrac{3}{2}\right\}$，符合题意，}
\step{综上，实数 $a$ 的值为 $2$ 或 $\dfrac{3}{2}$，}
\step{当 $a=2$ 时，$M=\{1,2\}$；当 $a=\dfrac{3}{2}$ 时，$M=\left\{\dfrac{1}{2},1,\dfrac{3}{2}\right\}$.}
\solend

\fi

\ansspace[8]

\item 已知集合 $A=\{x\mid x^2-2x-3=0\}$，集合 $B=\{x\mid ax-1=0\}$. 若 $B\sub A$，求实数 $a$ 的值.

\ans{$a=0$ 或 $a=-1$ 或 $a=\dfrac{1}{3}$}

\ifanswer
\solbegin
\step{当 $B=\e$ 时，$a=0$，满足题意；}
\step{当 $B\neq\e$，即 $a\neq 0$ 时，$B=\left\{\dfrac{1}{a}\right\}$，}
\step{又 $A=\{x\mid x^2-2x-3=0\}=\{x\mid x=-1\text{ 或 }x=3\}$，$B\sub A$，}
\step{所以 $\dfrac{1}{a}=-1$ 或 $\dfrac{1}{a}=3$，所以 $a=-1$ 或 $a=\dfrac{1}{3}$.}
\step{综上所述，$a=0$ 或 $a=-1$ 或 $a=\dfrac{1}{3}$.}
\solend

\fi

\ansspace[8]

\item 设集合 $A=\{x\mid -2\leqslant x\leqslant 5\}$，$B=\{x\mid m+1\leqslant x\leqslant 2m-1\}$.

\begin{enumerate}[label=(\arabic*)]
\item 若 $B\sub A$，求实数 $m$ 的取值范围；
\item 当 $x\in\Z$ 时，求 $A$ 的非空真子集个数；
\item 当 $x\in\R$ 时，不存在元素 $x$ 使 $x\in A$ 与 $x\in B$ 同时成立，求实数 $m$ 的取值范围.
\end{enumerate}

\ans{（1）$\{m\mid m\leqslant 3\}$；（2）$254$；（3）$\{m\mid m<2\text{ 或 }m>4\}$}

\ifanswer
\solbegin
\stepm{（1）}{当 $m+1>2m-1$，即 $m<2$ 时，$B=\e$，满足 $B\sub A$.}
\step{当 $m+1\leqslant 2m-1$，即 $m\geqslant 2$ 时，要使 $B\sub A$ 成立，}
\step{需 $\begin{cases}m+1\geqslant -2\\ 2m-1\leqslant 5\end{cases}$，得 $2\leqslant m\leqslant 3$，}
\step{综上，$m$ 的范围为 $\{m\mid m\leqslant 3\}$；}
\stepm{（2）}{$A=\{x\mid -2\leqslant x\leqslant 5,x\in\Z\}=\{-2,-1,0,1,2,3,4,5\}$，}
\step{则 $A$ 的非空真子集个数为 $2^8-1-1=254$；}
\stepm{（3）}{因为没有元素 $x$ 使 $x\in A$ 与 $x\in B$ 同时成立，所以 $A$ 与 $B$ 交集为空集.}
\step{①若 $B=\e$，即 $m+1>2m-1$，得 $m<2$ 时满足条件；}
\step{②若 $B\neq\e$，则要满足的条件是 $\begin{cases}m+1\leqslant 2m-1\\ m+1>5\end{cases}$ 或 $\begin{cases}m+1\leqslant 2m-1\\ 2m-1<-2\end{cases}$，}
\step{解得 $m>4$.}
\step{综上，有 $m<2$ 或 $m>4$，所以 $m$ 的范围为 $\{m\mid m<2\text{ 或 }m>4\}$.}
\solend

\fi

\ansspace[10]

% 压轴题：其后已无内容，故作答区全用可丢弃胶水（\ansspace*），并在题目前先量一下版面。
\ansneed{7cm}

\item 已知集合 $A=\{x\mid ax=1\}$，$B=\{x\mid x^2-ax+b=0\}$.

\begin{enumerate}[label=(\arabic*)]
\item 若 $a=2$，且 $A\sub B$，求 $b$ 的值；
\item 当 $a^2\geqslant 4b$ 时，若 $B\sub A$，求 $a$、$b$ 的值；
\item 若 $A\sub B$，讨论 $a$、$b$ 的取值范围.
\end{enumerate}

\ans{（1）$\dfrac{3}{4}$；（2）$a=\pm\sqrt2$，$b=\dfrac{1}{2}$；（3）$a=0$ 时 $b\in\R$，$a\neq 0$ 时 $b\in(-\infty,1)$}

\ifanswer
\solbegin
\stepm{（1）}{当 $a=2$ 时，$A=\left\{\dfrac{1}{2}\right\}$，}
\step{由 $A\sub B$ 得 $\dfrac{1}{2}$ 是方程 $x^2-2x+b=0$ 的解，代入得 $\left(\dfrac{1}{2}\right)^2-2\times\dfrac{1}{2}+b=0$，}
\step{解得 $b=\dfrac{3}{4}$.}
\stepm{（2）}{当 $a^2\geqslant 4b$ 时，方程 $x^2-ax+b=0$ 有实根.}
\step{若 $a=0$，则 $A=\e$，要 $B\sub A$ 需 $B=\e$，}
\step{但此时方程 $x^2+b=0$ 在 $a^2\geqslant 4b$（即 $b\leqslant 0$）时有实根，矛盾，舍去.}
\step{若 $a\neq 0$，则 $A=\left\{\dfrac{1}{a}\right\}$，$B\sub A$ 且方程有实根，故方程有且仅有一个根 $\dfrac{1}{a}$.}
\step{由韦达定理，根的和为 $a=\dfrac{1}{a}+\dfrac{1}{a}$，根的积为 $b=\dfrac{1}{a}\times\dfrac{1}{a}$.}
\step{由根的和得 $a=\dfrac{2}{a}$，即 $a^2=2$，$a=\pm\sqrt2$，$b=\dfrac{1}{2}$.}
\stepm{（3）}{当 $a=0$ 时，$A=\e$，故 $A\sub B$ 恒成立，此时 $b\in\R$.}
\step{当 $a\neq 0$ 时，$A=\left\{\dfrac{1}{a}\right\}$，将 $x=\dfrac{1}{a}$ 代入 $x^2-ax+b=0$，得 $b=1-\dfrac{1}{a^2}$.}
\step{此时方程判别式 $\Delta=a^2-4b=\dfrac{(a^2-2)^2}{a^2}\geqslant 0$，}
\step{若 $\Delta=0$，则 $a=\pm\sqrt2$，$b=\dfrac{1}{2}$；}
\step{若 $\Delta>0$，则 $a\neq 0$ 且 $a\neq\pm\sqrt2$，此时 $b=1-\dfrac{1}{a^2}<1$.}
\step{综上，当 $a=0$ 时，$b\in\R$；当 $a\neq 0$ 时，$b\in(-\infty,1)$.}
\solend

\fi

\ansspace*[12]

\end{enumerate}

\end{document}
"
%
% 原始材料是微信公众号图文（公众号“魔都数学加油站”），正文为 9 张 PNG 页；原文本身就是一份
% **讲评版**——题目下方直接印出【答案】或【解析】。试题文字、答案与解析均照原卷录入，未改内容。
%
% 与原卷的排版差异（均已在 README“说明”中交代）：
%   ① 原文自**第 3 题**起（第 1、2 题未在该文中发布），故本稿题号亦自 3 起；
%   ② 原卷本就印有“一、填空题”“二、选择题”“三、解答题”三节标题（分别在原图 01/04/06.png
%      上，逐页放大核对过），本稿照原卷分节，只把标题改排成全稿统一的“大节标题”样式；
%   ③ 原卷填空、选择的答案直接印在题目下方，本稿统一改为试卷版隐去、答案版以【答案】给出，
%      使两版彻底分离；
%   ④ 原卷的实数集符号时作正体 $R$、时作粗体 $\mathbf{R}$（第 21(3) 与 (3) 末句均为粗体），
%      本稿按全稿惯例统一写作 $\R$.
%
% 原卷笔误一处（本稿已补，见该处上方的 `%` 注）：
%   ① 第 10 题【解析】首句原卷印作“这个元不是 $1$ 就是 $m$”，“元素”漏了“素”字——
%      本稿按 高一b 第 13 题“原卷漏排、本稿补上”的先例补出“素”字，其余照原卷。
% 除此一处外，本卷 9 页逐页放大核对，未发现别的笔误（题 16 的 $1995=15\times 133$ 与
% “$15^2$ 的倍数有 8 个”两处最易被误判为笔误，实测 $133$、$8$ 均正确）。

% 本篇在公众号“魔都数学加油站”连载里的编号是“27学年高一10”，本地编号与之对齐（高一10）；
% 对应关系见 README“与公众号连载号的对照表”。
\documentclass[a4paper,11pt]{article}
\usepackage{common}

\begin{document}

\examtitlest[3]{2026--2027 年复附浦东高一上周末作业 1}{集合与常用逻辑用语}

\bigsec{一、填空题}

\begin{enumerate}
\setcounter{enumi}{2}

\item 已知集合 $M=\{1,2m+1\}$，$N=\{-1,m^2\}$，且 $M=N$，则 $m=$ \blankline[2.4em].

\ans{$-1$}

\ifanswer
\solbegin
\step{若 $M=N$，必有 $\begin{cases}2m+1=-1\\ m^2=1\end{cases}$，解得 $m=-1$.}
\solend

\fi

\item 集合 $\{(x,y)\mid x+2y=2,\ x,y\in\N\}$ 用列举法表示为 \blankline[4.4em].

\ans{$\{(0,1),(2,0)\}$}

\ifanswer
\solbegin
\step{由题意得 $x+2y=2$ 可化为 $y=1-\dfrac{x}{2}$，}
\step{由 $y\in\N$，有 $y=1-\dfrac{x}{2}\geqslant 0$，解得 $x\leqslant 2$，}
\step{又由 $x\in\N$，得 $x$ 可能取值为 $0$，$1$，$2$，}
\step{当 $x=0$ 时，得 $y=1$，满足条件，}
\step{当 $x=1$ 时，得 $y=\dfrac{1}{2}\notin\N$，不满足条件，}
\step{当 $x=2$ 时，得 $y=0$，满足条件，}
\step{综上，集合用列举法表示为 $\{(0,1),(2,0)\}$.}
\solend

\fi

\item 已知集合 $A=\{1\}$，集合 $B=\{0,1,2\}$，则满足关系 $A\sub P\sub B$ 的所有集合 $P$ 为 \blankline[4.4em].

\ans{$\{1\},\{0,1\},\{1,2\},\{0,1,2\}$}

\ifanswer
\solbegin
\step{由题意，集合 $P$ 可以是 $\{1\},\{0,1\},\{1,2\},\{0,1,2\}$.}
\solend

\fi

\item 已知集合 $A=\{x\mid a-1\leqslant x\leqslant a+2\}$，$B=\{x\mid 3<x<5\}$，则使 $A\supseteq B$ 成立的实数 $a$ 的取值范围是 \blankline[4.4em].

\ans{$[3,4]$}

\ifanswer
\solbegin
\step{因为 $A\supseteq B$，所以 $\begin{cases}a+2\geqslant 5\\ a-1\leqslant 3\end{cases}$，解得 $3\leqslant a\leqslant 4$.}
\solend

\fi

\item 若集合 $\{x\mid mx^2-4x+1=0\}=\{n\}$，则 $m+n=$ \blankline[2.4em].

\ans{$\dfrac{1}{4}$ 或 $\dfrac{9}{2}$}

\ifanswer
\solbegin
\step{若集合 $\{x\mid mx^2-4x+1=0\}=\{n\}$，则 $mx^2-4x+1=0$ 有唯一实数根 $n$，}
\step{当 $m=0$ 时，得 $n=\dfrac{1}{4}$，$m+n=\dfrac{1}{4}$，}
\step{当 $m\neq 0$ 时，$\Delta=16-4m=0$，即 $m=4$，}
\step{此时 $4x^2-4x+1=0$ 的根为 $n=\dfrac{1}{2}$，$m+n=\dfrac{9}{2}$.}
\step{故 $m+n=\dfrac{1}{4}$ 或 $\dfrac{9}{2}$.}
\solend

\fi

\item 集合 $M=\{-1,2,3,5,6\}$ 的全部非空子集的元素和等于 \blankline[4.4em].

\ans{$240$}

\ifanswer
\solbegin
\step{集合 $M=\{-1,2,3,5,6\}$ 中共有 $5$ 个元素，}
\step{在其所有子集中，元素 $-1$ 出现的次数为 $2^{5-1}=2^4=16$，}
\step{同理 $2$、$3$、$5$、$6$ 都出现 $16$ 次，}
\step{则集合 $M$ 的全部非空子集的元素和为 $(-1+2+3+5+6)\times 16=240$.}
\solend

\fi

\item 设 $A=\{x\mid x^2-3x-4=0\}$，$B=\{x\mid ax-1=0\}$，若 $A\cap B=B$，则实数 $a$ 的值可以为 \blankline[4.4em].

\ans{$0$、$-1$、$\dfrac{1}{4}$}

\ifanswer
\solbegin
\step{$A=\{-1,4\}$，因为 $A\cap B=B$，所以 $B\sub A$，}
\step{①$a=0$ 时，$B=\e$，满足 $B\sub A$；}
\step{②$a\neq 0$ 时，$B=\left\{\dfrac{1}{a}\right\}$，则 $\dfrac{1}{a}=-1$ 或 $\dfrac{1}{a}=4$，所以 $a=-1$ 或 $\dfrac{1}{4}$；}
\step{所以 $a$ 的值可以为 $0$、$-1$、$\dfrac{1}{4}$.}
\solend

\fi

\item 设集合 $A=\{1,m\}$，其中 $m$ 为实数. 令集合 $B=\{a+b\mid a\in A,b\in A\}$，若 $A\cap B$ 恰有一个元素，则 $A\cup B$ 的元素之和为 \blankline[2.4em].

\ans{$5$ 或 $10$}

\ifanswer
\solbegin
% 注：原卷本句印作“这个元不是 $1$ 就是 $m$”——“元素”漏了“素”字。按 高一b 第 13 题
%     “原卷漏排、本稿补上”的先例，本稿补出“素”字，其余照原卷，见文件头“原卷笔误”。
\step{由题意得，若 $A\cap B$ 恰有一个元素，这个元素不是 $1$ 就是 $m$，}
\stepm{（1）}{当 $1\in B$ 时，}
\step{①若 $1=1+m$，则 $m=0$，此时 $A=\{1,0\},B=\{2,1,0\},A\cap B=\{0,1\}$，不满足题意，舍去；}
\step{②若 $1=2m$，则 $m=\dfrac{1}{2}$，此时 $A=\left\{1,\dfrac{1}{2}\right\}$，$B=\left\{2,\dfrac{3}{2},1\right\}$，$A\cap B=\{1\}$，}
\step{满足题意，所以 $A\cup B=\left\{1,\dfrac{1}{2},2,\dfrac{3}{2}\right\}$，其元素之和为 $5$，}
\stepm{（2）}{当 $m\in B$ 时，}
\step{①若 $m=1+m$，则 $m$ 无解，}
\step{②若 $m=2m$，则 $m=0$，此时 $A=\{1,0\},B=\{2,1,0\},A\cap B=\{0,1\}$，不满足题意，}
\step{③若 $m=2$，则 $A=\{1,2\},B=\{2,3,4\},A\cap B=\{2\}$，满足题意，}
\step{所以 $A\cup B=\{1,2,3,4\}$，其元素之和为 $10$，}
\step{综上所述，$A\cup B$ 的元素之和为 $5$ 或 $10$.}
\solend

\fi

\item 设非空集合 $Q\sub M$，当 $Q$ 中所有元素和为偶数时（集合为单元素时和为元素本身），称 $Q$ 是 $M$ 的偶子集. 若集合 $M=\{1,2,3,4,5,6,7\}$，则其偶子集 $Q$ 的个数为 \blankline[2.4em].

\ans{$63$}

\ifanswer
\solbegin
\step{偶子集可以分为只有奇数、偶数及偶数个奇数和偶数的集合，}
\step{奇数有 $4$ 个，偶数有 $3$ 个，}
\step{只有奇数的有 $\mathrm{C}_4^2+1=7$ 个，只有偶数的有 $\mathrm{C}_3^1+\mathrm{C}_3^2+1=7$ 个，}
\step{偶数个奇数和偶数的：两奇一偶 $\mathrm{C}_4^2\cdot\mathrm{C}_3^1=18$ 个，两奇二偶 $\mathrm{C}_4^2\cdot\mathrm{C}_3^2=18$ 个，}
\step{两奇三偶 $\mathrm{C}_4^2\cdot\mathrm{C}_3^3=6$ 个，四奇一偶 $\mathrm{C}_4^4\cdot\mathrm{C}_3^1=3$ 个，四奇二偶 $\mathrm{C}_4^4\cdot\mathrm{C}_3^2=3$ 个，}
\step{四奇三偶 $\mathrm{C}_4^4\cdot\mathrm{C}_3^3=1$ 个，所以共有 $7+7+18+18+6+3+3+1=63$ 个.}
\solend

\fi

\item 已知非空集合 $S=\left\{x\mid -\dfrac{1}{2}\leqslant x\leqslant m\right\}$ 满足：当 $x\in S$ 时，有 $x^2\in S$，则实数 $m$ 的取值范围是 \blankline[4.4em].

\ans{$\left[\dfrac{1}{4},1\right]$}

\ifanswer
\solbegin
\step{因为集合 $S=\left\{x\mid -\dfrac{1}{2}\leqslant x\leqslant m\right\}$ 是非空集合，所以 $m\geqslant -\dfrac{1}{2}$，}
\step{又当 $x\in S$ 时，有 $x^2\in S$，所以 $m^2\leqslant m$ 且 $m\geqslant \left(-\dfrac{1}{2}\right)^2=\dfrac{1}{4}$，所以 $\dfrac{1}{4}\leqslant m\leqslant 1$.}
\solend

\fi

\end{enumerate}

\bigsec{二、选择题}

\begin{enumerate}
\setcounter{enumi}{12}

\item 已知集合 $U=\{1+x,1\}$，$x\in U$，则 $x=$\mbox{（\quad）}

\keepwithprev
A. $-1$\qquad B. $0$\qquad C. $1$\qquad D. $2$

\ans{C}

\ifanswer
\solbegin
\step{因为集合 $U=\{1+x,1\}$，$x\in U$，所以 $1+x=x$ 或 $x=1$，}
\step{当 $x=1+x$ 时，即 $0=1$，此时显然不成立，舍去，}
\step{当 $x=1$，此时集合 $U=\{1,2\}$，满足互异性，}
\step{综上所述，$x=1$. 故选 C.}
\solend

\fi

\item 已知集合 $M=\{m\mid m=a+b\sqrt2,\ a,b\in\Q\}$，则下列四个元素中属于 $M$ 的元素的个数是\mbox{（\quad）}

①$1+\sqrt2\pi$；②$\sqrt{11+6\sqrt2}$；③$\dfrac{1}{2+\sqrt2}$；④$\sqrt{2-\sqrt3}+\sqrt{2+\sqrt3}$

\keepwithprev
A. $4$\qquad B. $3$\qquad C. $2$\qquad D. $1$

\ans{C}

\ifanswer
\solbegin
\step{①$1+\sqrt2\pi$，由于 $\sqrt2$、$\pi$ 都为无理数，因此 $1+\sqrt2\pi\notin M$；}
\step{②$\sqrt{11+6\sqrt2}=\sqrt{11+2\sqrt{18}}=3+\sqrt2\in M$；}
\step{③$\dfrac{1}{2+\sqrt2}=1-\dfrac{1}{2}\sqrt2\in M$；}
\step{④$\sqrt{2-\sqrt3}+\sqrt{2+\sqrt3}=\sqrt{\dfrac{8-2\sqrt{12}}{4}}+\sqrt{\dfrac{8+2\sqrt{12}}{4}}=\dfrac{\sqrt6-\sqrt2}{2}+\dfrac{\sqrt6+\sqrt2}{2}=\sqrt6=\sqrt3\cdot\sqrt2\notin M$；}
\step{则上述四个元素中属于 $M$ 的元素的个数是 $2$. 故选 C.}
\solend

\fi

\item 集合 $A=\{2,3,7,8,9,11\}$，集合 $B=\{1,2,4,6,8\}$，下面说法正确的是\mbox{（\quad）}

\keepwithprev
A. 两个集合的交集是 $\{2,8\}$

B. 两个集合没有并集

C. $B$ 集合的元素无法用 $\{x\mid x=f(k),k\in D\}$ 这种形式表示出来

D. 集合 $A$ 有 $6$ 个子集

\ans{A}

\ifanswer
\solbegin
\step{对于 A，$A\cap B=\{2,8\}$，故 A 正确；}
\step{对于 B，$A\cup B=\{1,2,3,4,6,7,8,9,11\}$，故 B 错误；}
\step{对于 C，集合 $B$ 可以表示为 $B=\{x\mid x=f(k),k\in\{1,2,3,4,5\}\}$，}
\step{其中 $f(1)=1,f(2)=2,f(3)=4,f(4)=6,f(5)=8$，}
\step{即集合 $B$ 可描述为 $\{x\mid x=f(k),k\in D\}$ 形式，故 C 错误；}
\step{对于 D，集合 $A$ 有 $6$ 个元素，其子集总数为 $2^6=64\neq 6$，故 D 错误.}
\step{故选 A.}
\solend

\fi

\item 设 $M=\{1,2,\cdots,1995\}$，$A$ 是 $M$ 的子集且满足条件：当 $x\in A$ 时，$15x\notin A$，则 $A$ 中元素的个数的最大值为\mbox{（\quad）}

\keepwithprev
A. $1862$\qquad B. $1866$\qquad C. $1868$\qquad D. $1870$

\ans{D}

\ifanswer
\solbegin
\step{$1995=15\times 133$. 故取出所有不是 $15$ 的倍数的数，共 $1862$ 个，}
\step{这些数均符合要求.}
\step{在所有 $15$ 的倍数的数中，$15^2$ 的倍数有 $8$ 个，这些数又可以取出，}
\step{这样共取出了 $1870$ 个，即 $|A|\geqslant 1870$.}
\step{又 $\{k,15k\}(k=9,10,11,\cdots,133)$ 中的两个元素不能同时取出，}
\step{故 $|A|\leqslant 1995-133+8=1870$，故选 D.}
\solend

\fi

\end{enumerate}

\bigsec{三、解答题}

\begin{enumerate}
\setcounter{enumi}{16}

\item 已知集合 $A=\{x\mid a\leqslant x\leqslant a+4\}$，$B=\{x\mid x\geqslant 6\}$.

\begin{enumerate}[label=(\arabic*)]
\item 若 $a=3$，求 $A\cup B$；
\item 若 $A\cap B=A$，求实数 $a$ 的取值范围.
\end{enumerate}

\ans{（1）$\{x\mid x\geqslant 3\}$；（2）$[6,+\infty)$}

\ifanswer
\solbegin
\stepm{（1）}{当 $a=3$ 时，$A=\{x\mid 3\leqslant x\leqslant 7\}$，所以 $A\cup B=\{x\mid x\geqslant 3\}$；}
\stepm{（2）}{若 $A\cap B=A$，则 $A\sub B$，所以 $a\geqslant 6$，所以实数 $a$ 的取值范围是 $[6,+\infty)$.}
\solend

\fi

\ansspace[7]

\item 已知集合 $M=\{x\mid (x-a)(x^2-ax+a-1)=0\}$.

\begin{enumerate}[label=(\arabic*)]
\item 若 $a=3$，求集合 $M$；
\item 若集合 $M$ 中各元素之和等于 $3$，求实数 $a$ 的值，并用列举法表示集合 $M$.
\end{enumerate}

\ans{（1）$\{1,2,3\}$；（2）$a=2$ 或 $\dfrac{3}{2}$，此时 $M=\{1,2\}$ 或 $M=\left\{\dfrac{1}{2},1,\dfrac{3}{2}\right\}$}

\ifanswer
\solbegin
\stepm{（1）}{$\{1,2,3\}$；}
\stepm{（2）}{当 $(x-a)(x^2-ax+a-1)=0$ 有重根时，重根只能算作集合的一个元素，}
\step{$(x-a)(x^2-ax+a-1)=(x-a)(x-1)[x-(a-1)]=0$，}
\step{当 $a=1$ 时，得 $M=\{0,1\}$，不符合题意；}
\step{当 $a-1=1$ 时，即 $a=2$ 时，得 $M=\{1,2\}$，符合题意；}
\step{当 $a\neq 1$ 且 $a\neq 2$ 时，此时 $M=\{1,a,a-1\}$，得 $1+a+a-1=3$，}
\step{解得 $a=\dfrac{3}{2}$，此时 $M=\left\{\dfrac{1}{2},1,\dfrac{3}{2}\right\}$，符合题意，}
\step{综上，实数 $a$ 的值为 $2$ 或 $\dfrac{3}{2}$，}
\step{当 $a=2$ 时，$M=\{1,2\}$；当 $a=\dfrac{3}{2}$ 时，$M=\left\{\dfrac{1}{2},1,\dfrac{3}{2}\right\}$.}
\solend

\fi

\ansspace[8]

\item 已知集合 $A=\{x\mid x^2-2x-3=0\}$，集合 $B=\{x\mid ax-1=0\}$. 若 $B\sub A$，求实数 $a$ 的值.

\ans{$a=0$ 或 $a=-1$ 或 $a=\dfrac{1}{3}$}

\ifanswer
\solbegin
\step{当 $B=\e$ 时，$a=0$，满足题意；}
\step{当 $B\neq\e$，即 $a\neq 0$ 时，$B=\left\{\dfrac{1}{a}\right\}$，}
\step{又 $A=\{x\mid x^2-2x-3=0\}=\{x\mid x=-1\text{ 或 }x=3\}$，$B\sub A$，}
\step{所以 $\dfrac{1}{a}=-1$ 或 $\dfrac{1}{a}=3$，所以 $a=-1$ 或 $a=\dfrac{1}{3}$.}
\step{综上所述，$a=0$ 或 $a=-1$ 或 $a=\dfrac{1}{3}$.}
\solend

\fi

\ansspace[8]

\item 设集合 $A=\{x\mid -2\leqslant x\leqslant 5\}$，$B=\{x\mid m+1\leqslant x\leqslant 2m-1\}$.

\begin{enumerate}[label=(\arabic*)]
\item 若 $B\sub A$，求实数 $m$ 的取值范围；
\item 当 $x\in\Z$ 时，求 $A$ 的非空真子集个数；
\item 当 $x\in\R$ 时，不存在元素 $x$ 使 $x\in A$ 与 $x\in B$ 同时成立，求实数 $m$ 的取值范围.
\end{enumerate}

\ans{（1）$\{m\mid m\leqslant 3\}$；（2）$254$；（3）$\{m\mid m<2\text{ 或 }m>4\}$}

\ifanswer
\solbegin
\stepm{（1）}{当 $m+1>2m-1$，即 $m<2$ 时，$B=\e$，满足 $B\sub A$.}
\step{当 $m+1\leqslant 2m-1$，即 $m\geqslant 2$ 时，要使 $B\sub A$ 成立，}
\step{需 $\begin{cases}m+1\geqslant -2\\ 2m-1\leqslant 5\end{cases}$，得 $2\leqslant m\leqslant 3$，}
\step{综上，$m$ 的范围为 $\{m\mid m\leqslant 3\}$；}
\stepm{（2）}{$A=\{x\mid -2\leqslant x\leqslant 5,x\in\Z\}=\{-2,-1,0,1,2,3,4,5\}$，}
\step{则 $A$ 的非空真子集个数为 $2^8-1-1=254$；}
\stepm{（3）}{因为没有元素 $x$ 使 $x\in A$ 与 $x\in B$ 同时成立，所以 $A$ 与 $B$ 交集为空集.}
\step{①若 $B=\e$，即 $m+1>2m-1$，得 $m<2$ 时满足条件；}
\step{②若 $B\neq\e$，则要满足的条件是 $\begin{cases}m+1\leqslant 2m-1\\ m+1>5\end{cases}$ 或 $\begin{cases}m+1\leqslant 2m-1\\ 2m-1<-2\end{cases}$，}
\step{解得 $m>4$.}
\step{综上，有 $m<2$ 或 $m>4$，所以 $m$ 的范围为 $\{m\mid m<2\text{ 或 }m>4\}$.}
\solend

\fi

\ansspace[10]

% 压轴题：其后已无内容，故作答区全用可丢弃胶水（\ansspace*），并在题目前先量一下版面。
\ansneed{7cm}

\item 已知集合 $A=\{x\mid ax=1\}$，$B=\{x\mid x^2-ax+b=0\}$.

\begin{enumerate}[label=(\arabic*)]
\item 若 $a=2$，且 $A\sub B$，求 $b$ 的值；
\item 当 $a^2\geqslant 4b$ 时，若 $B\sub A$，求 $a$、$b$ 的值；
\item 若 $A\sub B$，讨论 $a$、$b$ 的取值范围.
\end{enumerate}

\ans{（1）$\dfrac{3}{4}$；（2）$a=\pm\sqrt2$，$b=\dfrac{1}{2}$；（3）$a=0$ 时 $b\in\R$，$a\neq 0$ 时 $b\in(-\infty,1)$}

\ifanswer
\solbegin
\stepm{（1）}{当 $a=2$ 时，$A=\left\{\dfrac{1}{2}\right\}$，}
\step{由 $A\sub B$ 得 $\dfrac{1}{2}$ 是方程 $x^2-2x+b=0$ 的解，代入得 $\left(\dfrac{1}{2}\right)^2-2\times\dfrac{1}{2}+b=0$，}
\step{解得 $b=\dfrac{3}{4}$.}
\stepm{（2）}{当 $a^2\geqslant 4b$ 时，方程 $x^2-ax+b=0$ 有实根.}
\step{若 $a=0$，则 $A=\e$，要 $B\sub A$ 需 $B=\e$，}
\step{但此时方程 $x^2+b=0$ 在 $a^2\geqslant 4b$（即 $b\leqslant 0$）时有实根，矛盾，舍去.}
\step{若 $a\neq 0$，则 $A=\left\{\dfrac{1}{a}\right\}$，$B\sub A$ 且方程有实根，故方程有且仅有一个根 $\dfrac{1}{a}$.}
\step{由韦达定理，根的和为 $a=\dfrac{1}{a}+\dfrac{1}{a}$，根的积为 $b=\dfrac{1}{a}\times\dfrac{1}{a}$.}
\step{由根的和得 $a=\dfrac{2}{a}$，即 $a^2=2$，$a=\pm\sqrt2$，$b=\dfrac{1}{2}$.}
\stepm{（3）}{当 $a=0$ 时，$A=\e$，故 $A\sub B$ 恒成立，此时 $b\in\R$.}
\step{当 $a\neq 0$ 时，$A=\left\{\dfrac{1}{a}\right\}$，将 $x=\dfrac{1}{a}$ 代入 $x^2-ax+b=0$，得 $b=1-\dfrac{1}{a^2}$.}
\step{此时方程判别式 $\Delta=a^2-4b=\dfrac{(a^2-2)^2}{a^2}\geqslant 0$，}
\step{若 $\Delta=0$，则 $a=\pm\sqrt2$，$b=\dfrac{1}{2}$；}
\step{若 $\Delta>0$，则 $a\neq 0$ 且 $a\neq\pm\sqrt2$，此时 $b=1-\dfrac{1}{a^2}<1$.}
\step{综上，当 $a=0$ 时，$b\in\R$；当 $a\neq 0$ 时，$b\in(-\infty,1)$.}
\solend

\fi

\ansspace*[12]

\end{enumerate}

\end{document}
"
%
% 原始材料是微信公众号图文（公众号“魔都数学加油站”），正文为 9 张 PNG 页；原文本身就是一份
% **讲评版**——题目下方直接印出【答案】或【解析】。试题文字、答案与解析均照原卷录入，未改内容。
%
% 与原卷的排版差异（均已在 README“说明”中交代）：
%   ① 原文自**第 3 题**起（第 1、2 题未在该文中发布），故本稿题号亦自 3 起；
%   ② 原卷本就印有“一、填空题”“二、选择题”“三、解答题”三节标题（分别在原图 01/04/06.png
%      上，逐页放大核对过），本稿照原卷分节，只把标题改排成全稿统一的“大节标题”样式；
%   ③ 原卷填空、选择的答案直接印在题目下方，本稿统一改为试卷版隐去、答案版以【答案】给出，
%      使两版彻底分离；
%   ④ 原卷的实数集符号时作正体 $R$、时作粗体 $\mathbf{R}$（第 21(3) 与 (3) 末句均为粗体），
%      本稿按全稿惯例统一写作 $\R$.
%
% 原卷笔误一处（本稿已补，见该处上方的 `%` 注）：
%   ① 第 10 题【解析】首句原卷印作“这个元不是 $1$ 就是 $m$”，“元素”漏了“素”字——
%      本稿按 高一b 第 13 题“原卷漏排、本稿补上”的先例补出“素”字，其余照原卷。
% 除此一处外，本卷 9 页逐页放大核对，未发现别的笔误（题 16 的 $1995=15\times 133$ 与
% “$15^2$ 的倍数有 8 个”两处最易被误判为笔误，实测 $133$、$8$ 均正确）。

% 本篇在公众号“魔都数学加油站”连载里的编号是“27学年高一10”，本地编号与之对齐（高一10）；
% 对应关系见 README“与公众号连载号的对照表”。
\documentclass[a4paper,11pt]{article}
\usepackage{common}

\begin{document}

\examtitlest[3]{2026--2027 年复附浦东高一上周末作业 1}{集合与常用逻辑用语}

\bigsec{一、填空题}

\begin{enumerate}
\setcounter{enumi}{2}

\item 已知集合 $M=\{1,2m+1\}$，$N=\{-1,m^2\}$，且 $M=N$，则 $m=$ \blankline[2.4em].

\ans{$-1$}

\ifanswer
\solbegin
\step{若 $M=N$，必有 $\begin{cases}2m+1=-1\\ m^2=1\end{cases}$，解得 $m=-1$.}
\solend

\fi

\item 集合 $\{(x,y)\mid x+2y=2,\ x,y\in\N\}$ 用列举法表示为 \blankline[4.4em].

\ans{$\{(0,1),(2,0)\}$}

\ifanswer
\solbegin
\step{由题意得 $x+2y=2$ 可化为 $y=1-\dfrac{x}{2}$，}
\step{由 $y\in\N$，有 $y=1-\dfrac{x}{2}\geqslant 0$，解得 $x\leqslant 2$，}
\step{又由 $x\in\N$，得 $x$ 可能取值为 $0$，$1$，$2$，}
\step{当 $x=0$ 时，得 $y=1$，满足条件，}
\step{当 $x=1$ 时，得 $y=\dfrac{1}{2}\notin\N$，不满足条件，}
\step{当 $x=2$ 时，得 $y=0$，满足条件，}
\step{综上，集合用列举法表示为 $\{(0,1),(2,0)\}$.}
\solend

\fi

\item 已知集合 $A=\{1\}$，集合 $B=\{0,1,2\}$，则满足关系 $A\sub P\sub B$ 的所有集合 $P$ 为 \blankline[4.4em].

\ans{$\{1\},\{0,1\},\{1,2\},\{0,1,2\}$}

\ifanswer
\solbegin
\step{由题意，集合 $P$ 可以是 $\{1\},\{0,1\},\{1,2\},\{0,1,2\}$.}
\solend

\fi

\item 已知集合 $A=\{x\mid a-1\leqslant x\leqslant a+2\}$，$B=\{x\mid 3<x<5\}$，则使 $A\supseteq B$ 成立的实数 $a$ 的取值范围是 \blankline[4.4em].

\ans{$[3,4]$}

\ifanswer
\solbegin
\step{因为 $A\supseteq B$，所以 $\begin{cases}a+2\geqslant 5\\ a-1\leqslant 3\end{cases}$，解得 $3\leqslant a\leqslant 4$.}
\solend

\fi

\item 若集合 $\{x\mid mx^2-4x+1=0\}=\{n\}$，则 $m+n=$ \blankline[2.4em].

\ans{$\dfrac{1}{4}$ 或 $\dfrac{9}{2}$}

\ifanswer
\solbegin
\step{若集合 $\{x\mid mx^2-4x+1=0\}=\{n\}$，则 $mx^2-4x+1=0$ 有唯一实数根 $n$，}
\step{当 $m=0$ 时，得 $n=\dfrac{1}{4}$，$m+n=\dfrac{1}{4}$，}
\step{当 $m\neq 0$ 时，$\Delta=16-4m=0$，即 $m=4$，}
\step{此时 $4x^2-4x+1=0$ 的根为 $n=\dfrac{1}{2}$，$m+n=\dfrac{9}{2}$.}
\step{故 $m+n=\dfrac{1}{4}$ 或 $\dfrac{9}{2}$.}
\solend

\fi

\item 集合 $M=\{-1,2,3,5,6\}$ 的全部非空子集的元素和等于 \blankline[4.4em].

\ans{$240$}

\ifanswer
\solbegin
\step{集合 $M=\{-1,2,3,5,6\}$ 中共有 $5$ 个元素，}
\step{在其所有子集中，元素 $-1$ 出现的次数为 $2^{5-1}=2^4=16$，}
\step{同理 $2$、$3$、$5$、$6$ 都出现 $16$ 次，}
\step{则集合 $M$ 的全部非空子集的元素和为 $(-1+2+3+5+6)\times 16=240$.}
\solend

\fi

\item 设 $A=\{x\mid x^2-3x-4=0\}$，$B=\{x\mid ax-1=0\}$，若 $A\cap B=B$，则实数 $a$ 的值可以为 \blankline[4.4em].

\ans{$0$、$-1$、$\dfrac{1}{4}$}

\ifanswer
\solbegin
\step{$A=\{-1,4\}$，因为 $A\cap B=B$，所以 $B\sub A$，}
\step{①$a=0$ 时，$B=\e$，满足 $B\sub A$；}
\step{②$a\neq 0$ 时，$B=\left\{\dfrac{1}{a}\right\}$，则 $\dfrac{1}{a}=-1$ 或 $\dfrac{1}{a}=4$，所以 $a=-1$ 或 $\dfrac{1}{4}$；}
\step{所以 $a$ 的值可以为 $0$、$-1$、$\dfrac{1}{4}$.}
\solend

\fi

\item 设集合 $A=\{1,m\}$，其中 $m$ 为实数. 令集合 $B=\{a+b\mid a\in A,b\in A\}$，若 $A\cap B$ 恰有一个元素，则 $A\cup B$ 的元素之和为 \blankline[2.4em].

\ans{$5$ 或 $10$}

\ifanswer
\solbegin
% 注：原卷本句印作“这个元不是 $1$ 就是 $m$”——“元素”漏了“素”字。按 高一b 第 13 题
%     “原卷漏排、本稿补上”的先例，本稿补出“素”字，其余照原卷，见文件头“原卷笔误”。
\step{由题意得，若 $A\cap B$ 恰有一个元素，这个元素不是 $1$ 就是 $m$，}
\stepm{（1）}{当 $1\in B$ 时，}
\step{①若 $1=1+m$，则 $m=0$，此时 $A=\{1,0\},B=\{2,1,0\},A\cap B=\{0,1\}$，不满足题意，舍去；}
\step{②若 $1=2m$，则 $m=\dfrac{1}{2}$，此时 $A=\left\{1,\dfrac{1}{2}\right\}$，$B=\left\{2,\dfrac{3}{2},1\right\}$，$A\cap B=\{1\}$，}
\step{满足题意，所以 $A\cup B=\left\{1,\dfrac{1}{2},2,\dfrac{3}{2}\right\}$，其元素之和为 $5$，}
\stepm{（2）}{当 $m\in B$ 时，}
\step{①若 $m=1+m$，则 $m$ 无解，}
\step{②若 $m=2m$，则 $m=0$，此时 $A=\{1,0\},B=\{2,1,0\},A\cap B=\{0,1\}$，不满足题意，}
\step{③若 $m=2$，则 $A=\{1,2\},B=\{2,3,4\},A\cap B=\{2\}$，满足题意，}
\step{所以 $A\cup B=\{1,2,3,4\}$，其元素之和为 $10$，}
\step{综上所述，$A\cup B$ 的元素之和为 $5$ 或 $10$.}
\solend

\fi

\item 设非空集合 $Q\sub M$，当 $Q$ 中所有元素和为偶数时（集合为单元素时和为元素本身），称 $Q$ 是 $M$ 的偶子集. 若集合 $M=\{1,2,3,4,5,6,7\}$，则其偶子集 $Q$ 的个数为 \blankline[2.4em].

\ans{$63$}

\ifanswer
\solbegin
\step{偶子集可以分为只有奇数、偶数及偶数个奇数和偶数的集合，}
\step{奇数有 $4$ 个，偶数有 $3$ 个，}
\step{只有奇数的有 $\mathrm{C}_4^2+1=7$ 个，只有偶数的有 $\mathrm{C}_3^1+\mathrm{C}_3^2+1=7$ 个，}
\step{偶数个奇数和偶数的：两奇一偶 $\mathrm{C}_4^2\cdot\mathrm{C}_3^1=18$ 个，两奇二偶 $\mathrm{C}_4^2\cdot\mathrm{C}_3^2=18$ 个，}
\step{两奇三偶 $\mathrm{C}_4^2\cdot\mathrm{C}_3^3=6$ 个，四奇一偶 $\mathrm{C}_4^4\cdot\mathrm{C}_3^1=3$ 个，四奇二偶 $\mathrm{C}_4^4\cdot\mathrm{C}_3^2=3$ 个，}
\step{四奇三偶 $\mathrm{C}_4^4\cdot\mathrm{C}_3^3=1$ 个，所以共有 $7+7+18+18+6+3+3+1=63$ 个.}
\solend

\fi

\item 已知非空集合 $S=\left\{x\mid -\dfrac{1}{2}\leqslant x\leqslant m\right\}$ 满足：当 $x\in S$ 时，有 $x^2\in S$，则实数 $m$ 的取值范围是 \blankline[4.4em].

\ans{$\left[\dfrac{1}{4},1\right]$}

\ifanswer
\solbegin
\step{因为集合 $S=\left\{x\mid -\dfrac{1}{2}\leqslant x\leqslant m\right\}$ 是非空集合，所以 $m\geqslant -\dfrac{1}{2}$，}
\step{又当 $x\in S$ 时，有 $x^2\in S$，所以 $m^2\leqslant m$ 且 $m\geqslant \left(-\dfrac{1}{2}\right)^2=\dfrac{1}{4}$，所以 $\dfrac{1}{4}\leqslant m\leqslant 1$.}
\solend

\fi

\end{enumerate}

\bigsec{二、选择题}

\begin{enumerate}
\setcounter{enumi}{12}

\item 已知集合 $U=\{1+x,1\}$，$x\in U$，则 $x=$\mbox{（\quad）}

\keepwithprev
A. $-1$\qquad B. $0$\qquad C. $1$\qquad D. $2$

\ans{C}

\ifanswer
\solbegin
\step{因为集合 $U=\{1+x,1\}$，$x\in U$，所以 $1+x=x$ 或 $x=1$，}
\step{当 $x=1+x$ 时，即 $0=1$，此时显然不成立，舍去，}
\step{当 $x=1$，此时集合 $U=\{1,2\}$，满足互异性，}
\step{综上所述，$x=1$. 故选 C.}
\solend

\fi

\item 已知集合 $M=\{m\mid m=a+b\sqrt2,\ a,b\in\Q\}$，则下列四个元素中属于 $M$ 的元素的个数是\mbox{（\quad）}

①$1+\sqrt2\pi$；②$\sqrt{11+6\sqrt2}$；③$\dfrac{1}{2+\sqrt2}$；④$\sqrt{2-\sqrt3}+\sqrt{2+\sqrt3}$

\keepwithprev
A. $4$\qquad B. $3$\qquad C. $2$\qquad D. $1$

\ans{C}

\ifanswer
\solbegin
\step{①$1+\sqrt2\pi$，由于 $\sqrt2$、$\pi$ 都为无理数，因此 $1+\sqrt2\pi\notin M$；}
\step{②$\sqrt{11+6\sqrt2}=\sqrt{11+2\sqrt{18}}=3+\sqrt2\in M$；}
\step{③$\dfrac{1}{2+\sqrt2}=1-\dfrac{1}{2}\sqrt2\in M$；}
\step{④$\sqrt{2-\sqrt3}+\sqrt{2+\sqrt3}=\sqrt{\dfrac{8-2\sqrt{12}}{4}}+\sqrt{\dfrac{8+2\sqrt{12}}{4}}=\dfrac{\sqrt6-\sqrt2}{2}+\dfrac{\sqrt6+\sqrt2}{2}=\sqrt6=\sqrt3\cdot\sqrt2\notin M$；}
\step{则上述四个元素中属于 $M$ 的元素的个数是 $2$. 故选 C.}
\solend

\fi

\item 集合 $A=\{2,3,7,8,9,11\}$，集合 $B=\{1,2,4,6,8\}$，下面说法正确的是\mbox{（\quad）}

\keepwithprev
A. 两个集合的交集是 $\{2,8\}$

B. 两个集合没有并集

C. $B$ 集合的元素无法用 $\{x\mid x=f(k),k\in D\}$ 这种形式表示出来

D. 集合 $A$ 有 $6$ 个子集

\ans{A}

\ifanswer
\solbegin
\step{对于 A，$A\cap B=\{2,8\}$，故 A 正确；}
\step{对于 B，$A\cup B=\{1,2,3,4,6,7,8,9,11\}$，故 B 错误；}
\step{对于 C，集合 $B$ 可以表示为 $B=\{x\mid x=f(k),k\in\{1,2,3,4,5\}\}$，}
\step{其中 $f(1)=1,f(2)=2,f(3)=4,f(4)=6,f(5)=8$，}
\step{即集合 $B$ 可描述为 $\{x\mid x=f(k),k\in D\}$ 形式，故 C 错误；}
\step{对于 D，集合 $A$ 有 $6$ 个元素，其子集总数为 $2^6=64\neq 6$，故 D 错误.}
\step{故选 A.}
\solend

\fi

\item 设 $M=\{1,2,\cdots,1995\}$，$A$ 是 $M$ 的子集且满足条件：当 $x\in A$ 时，$15x\notin A$，则 $A$ 中元素的个数的最大值为\mbox{（\quad）}

\keepwithprev
A. $1862$\qquad B. $1866$\qquad C. $1868$\qquad D. $1870$

\ans{D}

\ifanswer
\solbegin
\step{$1995=15\times 133$. 故取出所有不是 $15$ 的倍数的数，共 $1862$ 个，}
\step{这些数均符合要求.}
\step{在所有 $15$ 的倍数的数中，$15^2$ 的倍数有 $8$ 个，这些数又可以取出，}
\step{这样共取出了 $1870$ 个，即 $|A|\geqslant 1870$.}
\step{又 $\{k,15k\}(k=9,10,11,\cdots,133)$ 中的两个元素不能同时取出，}
\step{故 $|A|\leqslant 1995-133+8=1870$，故选 D.}
\solend

\fi

\end{enumerate}

\bigsec{三、解答题}

\begin{enumerate}
\setcounter{enumi}{16}

\item 已知集合 $A=\{x\mid a\leqslant x\leqslant a+4\}$，$B=\{x\mid x\geqslant 6\}$.

\begin{enumerate}[label=(\arabic*)]
\item 若 $a=3$，求 $A\cup B$；
\item 若 $A\cap B=A$，求实数 $a$ 的取值范围.
\end{enumerate}

\ans{（1）$\{x\mid x\geqslant 3\}$；（2）$[6,+\infty)$}

\ifanswer
\solbegin
\stepm{（1）}{当 $a=3$ 时，$A=\{x\mid 3\leqslant x\leqslant 7\}$，所以 $A\cup B=\{x\mid x\geqslant 3\}$；}
\stepm{（2）}{若 $A\cap B=A$，则 $A\sub B$，所以 $a\geqslant 6$，所以实数 $a$ 的取值范围是 $[6,+\infty)$.}
\solend

\fi

\ansspace[7]

\item 已知集合 $M=\{x\mid (x-a)(x^2-ax+a-1)=0\}$.

\begin{enumerate}[label=(\arabic*)]
\item 若 $a=3$，求集合 $M$；
\item 若集合 $M$ 中各元素之和等于 $3$，求实数 $a$ 的值，并用列举法表示集合 $M$.
\end{enumerate}

\ans{（1）$\{1,2,3\}$；（2）$a=2$ 或 $\dfrac{3}{2}$，此时 $M=\{1,2\}$ 或 $M=\left\{\dfrac{1}{2},1,\dfrac{3}{2}\right\}$}

\ifanswer
\solbegin
\stepm{（1）}{$\{1,2,3\}$；}
\stepm{（2）}{当 $(x-a)(x^2-ax+a-1)=0$ 有重根时，重根只能算作集合的一个元素，}
\step{$(x-a)(x^2-ax+a-1)=(x-a)(x-1)[x-(a-1)]=0$，}
\step{当 $a=1$ 时，得 $M=\{0,1\}$，不符合题意；}
\step{当 $a-1=1$ 时，即 $a=2$ 时，得 $M=\{1,2\}$，符合题意；}
\step{当 $a\neq 1$ 且 $a\neq 2$ 时，此时 $M=\{1,a,a-1\}$，得 $1+a+a-1=3$，}
\step{解得 $a=\dfrac{3}{2}$，此时 $M=\left\{\dfrac{1}{2},1,\dfrac{3}{2}\right\}$，符合题意，}
\step{综上，实数 $a$ 的值为 $2$ 或 $\dfrac{3}{2}$，}
\step{当 $a=2$ 时，$M=\{1,2\}$；当 $a=\dfrac{3}{2}$ 时，$M=\left\{\dfrac{1}{2},1,\dfrac{3}{2}\right\}$.}
\solend

\fi

\ansspace[8]

\item 已知集合 $A=\{x\mid x^2-2x-3=0\}$，集合 $B=\{x\mid ax-1=0\}$. 若 $B\sub A$，求实数 $a$ 的值.

\ans{$a=0$ 或 $a=-1$ 或 $a=\dfrac{1}{3}$}

\ifanswer
\solbegin
\step{当 $B=\e$ 时，$a=0$，满足题意；}
\step{当 $B\neq\e$，即 $a\neq 0$ 时，$B=\left\{\dfrac{1}{a}\right\}$，}
\step{又 $A=\{x\mid x^2-2x-3=0\}=\{x\mid x=-1\text{ 或 }x=3\}$，$B\sub A$，}
\step{所以 $\dfrac{1}{a}=-1$ 或 $\dfrac{1}{a}=3$，所以 $a=-1$ 或 $a=\dfrac{1}{3}$.}
\step{综上所述，$a=0$ 或 $a=-1$ 或 $a=\dfrac{1}{3}$.}
\solend

\fi

\ansspace[8]

\item 设集合 $A=\{x\mid -2\leqslant x\leqslant 5\}$，$B=\{x\mid m+1\leqslant x\leqslant 2m-1\}$.

\begin{enumerate}[label=(\arabic*)]
\item 若 $B\sub A$，求实数 $m$ 的取值范围；
\item 当 $x\in\Z$ 时，求 $A$ 的非空真子集个数；
\item 当 $x\in\R$ 时，不存在元素 $x$ 使 $x\in A$ 与 $x\in B$ 同时成立，求实数 $m$ 的取值范围.
\end{enumerate}

\ans{（1）$\{m\mid m\leqslant 3\}$；（2）$254$；（3）$\{m\mid m<2\text{ 或 }m>4\}$}

\ifanswer
\solbegin
\stepm{（1）}{当 $m+1>2m-1$，即 $m<2$ 时，$B=\e$，满足 $B\sub A$.}
\step{当 $m+1\leqslant 2m-1$，即 $m\geqslant 2$ 时，要使 $B\sub A$ 成立，}
\step{需 $\begin{cases}m+1\geqslant -2\\ 2m-1\leqslant 5\end{cases}$，得 $2\leqslant m\leqslant 3$，}
\step{综上，$m$ 的范围为 $\{m\mid m\leqslant 3\}$；}
\stepm{（2）}{$A=\{x\mid -2\leqslant x\leqslant 5,x\in\Z\}=\{-2,-1,0,1,2,3,4,5\}$，}
\step{则 $A$ 的非空真子集个数为 $2^8-1-1=254$；}
\stepm{（3）}{因为没有元素 $x$ 使 $x\in A$ 与 $x\in B$ 同时成立，所以 $A$ 与 $B$ 交集为空集.}
\step{①若 $B=\e$，即 $m+1>2m-1$，得 $m<2$ 时满足条件；}
\step{②若 $B\neq\e$，则要满足的条件是 $\begin{cases}m+1\leqslant 2m-1\\ m+1>5\end{cases}$ 或 $\begin{cases}m+1\leqslant 2m-1\\ 2m-1<-2\end{cases}$，}
\step{解得 $m>4$.}
\step{综上，有 $m<2$ 或 $m>4$，所以 $m$ 的范围为 $\{m\mid m<2\text{ 或 }m>4\}$.}
\solend

\fi

\ansspace[10]

% 压轴题：其后已无内容，故作答区全用可丢弃胶水（\ansspace*），并在题目前先量一下版面。
\ansneed{7cm}

\item 已知集合 $A=\{x\mid ax=1\}$，$B=\{x\mid x^2-ax+b=0\}$.

\begin{enumerate}[label=(\arabic*)]
\item 若 $a=2$，且 $A\sub B$，求 $b$ 的值；
\item 当 $a^2\geqslant 4b$ 时，若 $B\sub A$，求 $a$、$b$ 的值；
\item 若 $A\sub B$，讨论 $a$、$b$ 的取值范围.
\end{enumerate}

\ans{（1）$\dfrac{3}{4}$；（2）$a=\pm\sqrt2$，$b=\dfrac{1}{2}$；（3）$a=0$ 时 $b\in\R$，$a\neq 0$ 时 $b\in(-\infty,1)$}

\ifanswer
\solbegin
\stepm{（1）}{当 $a=2$ 时，$A=\left\{\dfrac{1}{2}\right\}$，}
\step{由 $A\sub B$ 得 $\dfrac{1}{2}$ 是方程 $x^2-2x+b=0$ 的解，代入得 $\left(\dfrac{1}{2}\right)^2-2\times\dfrac{1}{2}+b=0$，}
\step{解得 $b=\dfrac{3}{4}$.}
\stepm{（2）}{当 $a^2\geqslant 4b$ 时，方程 $x^2-ax+b=0$ 有实根.}
\step{若 $a=0$，则 $A=\e$，要 $B\sub A$ 需 $B=\e$，}
\step{但此时方程 $x^2+b=0$ 在 $a^2\geqslant 4b$（即 $b\leqslant 0$）时有实根，矛盾，舍去.}
\step{若 $a\neq 0$，则 $A=\left\{\dfrac{1}{a}\right\}$，$B\sub A$ 且方程有实根，故方程有且仅有一个根 $\dfrac{1}{a}$.}
\step{由韦达定理，根的和为 $a=\dfrac{1}{a}+\dfrac{1}{a}$，根的积为 $b=\dfrac{1}{a}\times\dfrac{1}{a}$.}
\step{由根的和得 $a=\dfrac{2}{a}$，即 $a^2=2$，$a=\pm\sqrt2$，$b=\dfrac{1}{2}$.}
\stepm{（3）}{当 $a=0$ 时，$A=\e$，故 $A\sub B$ 恒成立，此时 $b\in\R$.}
\step{当 $a\neq 0$ 时，$A=\left\{\dfrac{1}{a}\right\}$，将 $x=\dfrac{1}{a}$ 代入 $x^2-ax+b=0$，得 $b=1-\dfrac{1}{a^2}$.}
\step{此时方程判别式 $\Delta=a^2-4b=\dfrac{(a^2-2)^2}{a^2}\geqslant 0$，}
\step{若 $\Delta=0$，则 $a=\pm\sqrt2$，$b=\dfrac{1}{2}$；}
\step{若 $\Delta>0$，则 $a\neq 0$ 且 $a\neq\pm\sqrt2$，此时 $b=1-\dfrac{1}{a^2}<1$.}
\step{综上，当 $a=0$ 时，$b\in\R$；当 $a\neq 0$ 时，$b\in(-\infty,1)$.}
\solend

\fi

\ansspace*[12]

\end{enumerate}

\end{document}
"
% 紧凑版：xelatex "\def\iscompact{1}% 高一10_复附浦东_周末作业1.tex —— 高一上数学“2026-2027 年复附浦东高一上周末作业 1”（试卷版/答案版）
% 试卷版：xelatex 高一10_复附浦东_周末作业1.tex
% 答案版：xelatex "\def\isanswer{1}% 高一10_复附浦东_周末作业1.tex —— 高一上数学“2026-2027 年复附浦东高一上周末作业 1”（试卷版/答案版）
% 试卷版：xelatex 高一10_复附浦东_周末作业1.tex
% 答案版：xelatex "\def\isanswer{1}% 高一10_复附浦东_周末作业1.tex —— 高一上数学“2026-2027 年复附浦东高一上周末作业 1”（试卷版/答案版）
% 试卷版：xelatex 高一10_复附浦东_周末作业1.tex
% 答案版：xelatex "\def\isanswer{1}\input{高一10_复附浦东_周末作业1.tex}"
% 紧凑版：xelatex "\def\iscompact{1}\input{高一10_复附浦东_周末作业1.tex}"
%
% 原始材料是微信公众号图文（公众号“魔都数学加油站”），正文为 9 张 PNG 页；原文本身就是一份
% **讲评版**——题目下方直接印出【答案】或【解析】。试题文字、答案与解析均照原卷录入，未改内容。
%
% 与原卷的排版差异（均已在 README“说明”中交代）：
%   ① 原文自**第 3 题**起（第 1、2 题未在该文中发布），故本稿题号亦自 3 起；
%   ② 原卷本就印有“一、填空题”“二、选择题”“三、解答题”三节标题（分别在原图 01/04/06.png
%      上，逐页放大核对过），本稿照原卷分节，只把标题改排成全稿统一的“大节标题”样式；
%   ③ 原卷填空、选择的答案直接印在题目下方，本稿统一改为试卷版隐去、答案版以【答案】给出，
%      使两版彻底分离；
%   ④ 原卷的实数集符号时作正体 $R$、时作粗体 $\mathbf{R}$（第 21(3) 与 (3) 末句均为粗体），
%      本稿按全稿惯例统一写作 $\R$.
%
% 原卷笔误一处（本稿已补，见该处上方的 `%` 注）：
%   ① 第 10 题【解析】首句原卷印作“这个元不是 $1$ 就是 $m$”，“元素”漏了“素”字——
%      本稿按 高一b 第 13 题“原卷漏排、本稿补上”的先例补出“素”字，其余照原卷。
% 除此一处外，本卷 9 页逐页放大核对，未发现别的笔误（题 16 的 $1995=15\times 133$ 与
% “$15^2$ 的倍数有 8 个”两处最易被误判为笔误，实测 $133$、$8$ 均正确）。

% 本篇在公众号“魔都数学加油站”连载里的编号是“27学年高一10”，本地编号与之对齐（高一10）；
% 对应关系见 README“与公众号连载号的对照表”。
\documentclass[a4paper,11pt]{article}
\usepackage{common}

\begin{document}

\examtitlest[3]{2026--2027 年复附浦东高一上周末作业 1}{集合与常用逻辑用语}

\bigsec{一、填空题}

\begin{enumerate}
\setcounter{enumi}{2}

\item 已知集合 $M=\{1,2m+1\}$，$N=\{-1,m^2\}$，且 $M=N$，则 $m=$ \blankline[2.4em].

\ans{$-1$}

\ifanswer
\solbegin
\step{若 $M=N$，必有 $\begin{cases}2m+1=-1\\ m^2=1\end{cases}$，解得 $m=-1$.}
\solend

\fi

\item 集合 $\{(x,y)\mid x+2y=2,\ x,y\in\N\}$ 用列举法表示为 \blankline[4.4em].

\ans{$\{(0,1),(2,0)\}$}

\ifanswer
\solbegin
\step{由题意得 $x+2y=2$ 可化为 $y=1-\dfrac{x}{2}$，}
\step{由 $y\in\N$，有 $y=1-\dfrac{x}{2}\geqslant 0$，解得 $x\leqslant 2$，}
\step{又由 $x\in\N$，得 $x$ 可能取值为 $0$，$1$，$2$，}
\step{当 $x=0$ 时，得 $y=1$，满足条件，}
\step{当 $x=1$ 时，得 $y=\dfrac{1}{2}\notin\N$，不满足条件，}
\step{当 $x=2$ 时，得 $y=0$，满足条件，}
\step{综上，集合用列举法表示为 $\{(0,1),(2,0)\}$.}
\solend

\fi

\item 已知集合 $A=\{1\}$，集合 $B=\{0,1,2\}$，则满足关系 $A\sub P\sub B$ 的所有集合 $P$ 为 \blankline[4.4em].

\ans{$\{1\},\{0,1\},\{1,2\},\{0,1,2\}$}

\ifanswer
\solbegin
\step{由题意，集合 $P$ 可以是 $\{1\},\{0,1\},\{1,2\},\{0,1,2\}$.}
\solend

\fi

\item 已知集合 $A=\{x\mid a-1\leqslant x\leqslant a+2\}$，$B=\{x\mid 3<x<5\}$，则使 $A\supseteq B$ 成立的实数 $a$ 的取值范围是 \blankline[4.4em].

\ans{$[3,4]$}

\ifanswer
\solbegin
\step{因为 $A\supseteq B$，所以 $\begin{cases}a+2\geqslant 5\\ a-1\leqslant 3\end{cases}$，解得 $3\leqslant a\leqslant 4$.}
\solend

\fi

\item 若集合 $\{x\mid mx^2-4x+1=0\}=\{n\}$，则 $m+n=$ \blankline[2.4em].

\ans{$\dfrac{1}{4}$ 或 $\dfrac{9}{2}$}

\ifanswer
\solbegin
\step{若集合 $\{x\mid mx^2-4x+1=0\}=\{n\}$，则 $mx^2-4x+1=0$ 有唯一实数根 $n$，}
\step{当 $m=0$ 时，得 $n=\dfrac{1}{4}$，$m+n=\dfrac{1}{4}$，}
\step{当 $m\neq 0$ 时，$\Delta=16-4m=0$，即 $m=4$，}
\step{此时 $4x^2-4x+1=0$ 的根为 $n=\dfrac{1}{2}$，$m+n=\dfrac{9}{2}$.}
\step{故 $m+n=\dfrac{1}{4}$ 或 $\dfrac{9}{2}$.}
\solend

\fi

\item 集合 $M=\{-1,2,3,5,6\}$ 的全部非空子集的元素和等于 \blankline[4.4em].

\ans{$240$}

\ifanswer
\solbegin
\step{集合 $M=\{-1,2,3,5,6\}$ 中共有 $5$ 个元素，}
\step{在其所有子集中，元素 $-1$ 出现的次数为 $2^{5-1}=2^4=16$，}
\step{同理 $2$、$3$、$5$、$6$ 都出现 $16$ 次，}
\step{则集合 $M$ 的全部非空子集的元素和为 $(-1+2+3+5+6)\times 16=240$.}
\solend

\fi

\item 设 $A=\{x\mid x^2-3x-4=0\}$，$B=\{x\mid ax-1=0\}$，若 $A\cap B=B$，则实数 $a$ 的值可以为 \blankline[4.4em].

\ans{$0$、$-1$、$\dfrac{1}{4}$}

\ifanswer
\solbegin
\step{$A=\{-1,4\}$，因为 $A\cap B=B$，所以 $B\sub A$，}
\step{①$a=0$ 时，$B=\e$，满足 $B\sub A$；}
\step{②$a\neq 0$ 时，$B=\left\{\dfrac{1}{a}\right\}$，则 $\dfrac{1}{a}=-1$ 或 $\dfrac{1}{a}=4$，所以 $a=-1$ 或 $\dfrac{1}{4}$；}
\step{所以 $a$ 的值可以为 $0$、$-1$、$\dfrac{1}{4}$.}
\solend

\fi

\item 设集合 $A=\{1,m\}$，其中 $m$ 为实数. 令集合 $B=\{a+b\mid a\in A,b\in A\}$，若 $A\cap B$ 恰有一个元素，则 $A\cup B$ 的元素之和为 \blankline[2.4em].

\ans{$5$ 或 $10$}

\ifanswer
\solbegin
% 注：原卷本句印作“这个元不是 $1$ 就是 $m$”——“元素”漏了“素”字。按 高一b 第 13 题
%     “原卷漏排、本稿补上”的先例，本稿补出“素”字，其余照原卷，见文件头“原卷笔误”。
\step{由题意得，若 $A\cap B$ 恰有一个元素，这个元素不是 $1$ 就是 $m$，}
\stepm{（1）}{当 $1\in B$ 时，}
\step{①若 $1=1+m$，则 $m=0$，此时 $A=\{1,0\},B=\{2,1,0\},A\cap B=\{0,1\}$，不满足题意，舍去；}
\step{②若 $1=2m$，则 $m=\dfrac{1}{2}$，此时 $A=\left\{1,\dfrac{1}{2}\right\}$，$B=\left\{2,\dfrac{3}{2},1\right\}$，$A\cap B=\{1\}$，}
\step{满足题意，所以 $A\cup B=\left\{1,\dfrac{1}{2},2,\dfrac{3}{2}\right\}$，其元素之和为 $5$，}
\stepm{（2）}{当 $m\in B$ 时，}
\step{①若 $m=1+m$，则 $m$ 无解，}
\step{②若 $m=2m$，则 $m=0$，此时 $A=\{1,0\},B=\{2,1,0\},A\cap B=\{0,1\}$，不满足题意，}
\step{③若 $m=2$，则 $A=\{1,2\},B=\{2,3,4\},A\cap B=\{2\}$，满足题意，}
\step{所以 $A\cup B=\{1,2,3,4\}$，其元素之和为 $10$，}
\step{综上所述，$A\cup B$ 的元素之和为 $5$ 或 $10$.}
\solend

\fi

\item 设非空集合 $Q\sub M$，当 $Q$ 中所有元素和为偶数时（集合为单元素时和为元素本身），称 $Q$ 是 $M$ 的偶子集. 若集合 $M=\{1,2,3,4,5,6,7\}$，则其偶子集 $Q$ 的个数为 \blankline[2.4em].

\ans{$63$}

\ifanswer
\solbegin
\step{偶子集可以分为只有奇数、偶数及偶数个奇数和偶数的集合，}
\step{奇数有 $4$ 个，偶数有 $3$ 个，}
\step{只有奇数的有 $\mathrm{C}_4^2+1=7$ 个，只有偶数的有 $\mathrm{C}_3^1+\mathrm{C}_3^2+1=7$ 个，}
\step{偶数个奇数和偶数的：两奇一偶 $\mathrm{C}_4^2\cdot\mathrm{C}_3^1=18$ 个，两奇二偶 $\mathrm{C}_4^2\cdot\mathrm{C}_3^2=18$ 个，}
\step{两奇三偶 $\mathrm{C}_4^2\cdot\mathrm{C}_3^3=6$ 个，四奇一偶 $\mathrm{C}_4^4\cdot\mathrm{C}_3^1=3$ 个，四奇二偶 $\mathrm{C}_4^4\cdot\mathrm{C}_3^2=3$ 个，}
\step{四奇三偶 $\mathrm{C}_4^4\cdot\mathrm{C}_3^3=1$ 个，所以共有 $7+7+18+18+6+3+3+1=63$ 个.}
\solend

\fi

\item 已知非空集合 $S=\left\{x\mid -\dfrac{1}{2}\leqslant x\leqslant m\right\}$ 满足：当 $x\in S$ 时，有 $x^2\in S$，则实数 $m$ 的取值范围是 \blankline[4.4em].

\ans{$\left[\dfrac{1}{4},1\right]$}

\ifanswer
\solbegin
\step{因为集合 $S=\left\{x\mid -\dfrac{1}{2}\leqslant x\leqslant m\right\}$ 是非空集合，所以 $m\geqslant -\dfrac{1}{2}$，}
\step{又当 $x\in S$ 时，有 $x^2\in S$，所以 $m^2\leqslant m$ 且 $m\geqslant \left(-\dfrac{1}{2}\right)^2=\dfrac{1}{4}$，所以 $\dfrac{1}{4}\leqslant m\leqslant 1$.}
\solend

\fi

\end{enumerate}

\bigsec{二、选择题}

\begin{enumerate}
\setcounter{enumi}{12}

\item 已知集合 $U=\{1+x,1\}$，$x\in U$，则 $x=$\mbox{（\quad）}

\keepwithprev
A. $-1$\qquad B. $0$\qquad C. $1$\qquad D. $2$

\ans{C}

\ifanswer
\solbegin
\step{因为集合 $U=\{1+x,1\}$，$x\in U$，所以 $1+x=x$ 或 $x=1$，}
\step{当 $x=1+x$ 时，即 $0=1$，此时显然不成立，舍去，}
\step{当 $x=1$，此时集合 $U=\{1,2\}$，满足互异性，}
\step{综上所述，$x=1$. 故选 C.}
\solend

\fi

\item 已知集合 $M=\{m\mid m=a+b\sqrt2,\ a,b\in\Q\}$，则下列四个元素中属于 $M$ 的元素的个数是\mbox{（\quad）}

①$1+\sqrt2\pi$；②$\sqrt{11+6\sqrt2}$；③$\dfrac{1}{2+\sqrt2}$；④$\sqrt{2-\sqrt3}+\sqrt{2+\sqrt3}$

\keepwithprev
A. $4$\qquad B. $3$\qquad C. $2$\qquad D. $1$

\ans{C}

\ifanswer
\solbegin
\step{①$1+\sqrt2\pi$，由于 $\sqrt2$、$\pi$ 都为无理数，因此 $1+\sqrt2\pi\notin M$；}
\step{②$\sqrt{11+6\sqrt2}=\sqrt{11+2\sqrt{18}}=3+\sqrt2\in M$；}
\step{③$\dfrac{1}{2+\sqrt2}=1-\dfrac{1}{2}\sqrt2\in M$；}
\step{④$\sqrt{2-\sqrt3}+\sqrt{2+\sqrt3}=\sqrt{\dfrac{8-2\sqrt{12}}{4}}+\sqrt{\dfrac{8+2\sqrt{12}}{4}}=\dfrac{\sqrt6-\sqrt2}{2}+\dfrac{\sqrt6+\sqrt2}{2}=\sqrt6=\sqrt3\cdot\sqrt2\notin M$；}
\step{则上述四个元素中属于 $M$ 的元素的个数是 $2$. 故选 C.}
\solend

\fi

\item 集合 $A=\{2,3,7,8,9,11\}$，集合 $B=\{1,2,4,6,8\}$，下面说法正确的是\mbox{（\quad）}

\keepwithprev
A. 两个集合的交集是 $\{2,8\}$

B. 两个集合没有并集

C. $B$ 集合的元素无法用 $\{x\mid x=f(k),k\in D\}$ 这种形式表示出来

D. 集合 $A$ 有 $6$ 个子集

\ans{A}

\ifanswer
\solbegin
\step{对于 A，$A\cap B=\{2,8\}$，故 A 正确；}
\step{对于 B，$A\cup B=\{1,2,3,4,6,7,8,9,11\}$，故 B 错误；}
\step{对于 C，集合 $B$ 可以表示为 $B=\{x\mid x=f(k),k\in\{1,2,3,4,5\}\}$，}
\step{其中 $f(1)=1,f(2)=2,f(3)=4,f(4)=6,f(5)=8$，}
\step{即集合 $B$ 可描述为 $\{x\mid x=f(k),k\in D\}$ 形式，故 C 错误；}
\step{对于 D，集合 $A$ 有 $6$ 个元素，其子集总数为 $2^6=64\neq 6$，故 D 错误.}
\step{故选 A.}
\solend

\fi

\item 设 $M=\{1,2,\cdots,1995\}$，$A$ 是 $M$ 的子集且满足条件：当 $x\in A$ 时，$15x\notin A$，则 $A$ 中元素的个数的最大值为\mbox{（\quad）}

\keepwithprev
A. $1862$\qquad B. $1866$\qquad C. $1868$\qquad D. $1870$

\ans{D}

\ifanswer
\solbegin
\step{$1995=15\times 133$. 故取出所有不是 $15$ 的倍数的数，共 $1862$ 个，}
\step{这些数均符合要求.}
\step{在所有 $15$ 的倍数的数中，$15^2$ 的倍数有 $8$ 个，这些数又可以取出，}
\step{这样共取出了 $1870$ 个，即 $|A|\geqslant 1870$.}
\step{又 $\{k,15k\}(k=9,10,11,\cdots,133)$ 中的两个元素不能同时取出，}
\step{故 $|A|\leqslant 1995-133+8=1870$，故选 D.}
\solend

\fi

\end{enumerate}

\bigsec{三、解答题}

\begin{enumerate}
\setcounter{enumi}{16}

\item 已知集合 $A=\{x\mid a\leqslant x\leqslant a+4\}$，$B=\{x\mid x\geqslant 6\}$.

\begin{enumerate}[label=(\arabic*)]
\item 若 $a=3$，求 $A\cup B$；
\item 若 $A\cap B=A$，求实数 $a$ 的取值范围.
\end{enumerate}

\ans{（1）$\{x\mid x\geqslant 3\}$；（2）$[6,+\infty)$}

\ifanswer
\solbegin
\stepm{（1）}{当 $a=3$ 时，$A=\{x\mid 3\leqslant x\leqslant 7\}$，所以 $A\cup B=\{x\mid x\geqslant 3\}$；}
\stepm{（2）}{若 $A\cap B=A$，则 $A\sub B$，所以 $a\geqslant 6$，所以实数 $a$ 的取值范围是 $[6,+\infty)$.}
\solend

\fi

\ansspace[7]

\item 已知集合 $M=\{x\mid (x-a)(x^2-ax+a-1)=0\}$.

\begin{enumerate}[label=(\arabic*)]
\item 若 $a=3$，求集合 $M$；
\item 若集合 $M$ 中各元素之和等于 $3$，求实数 $a$ 的值，并用列举法表示集合 $M$.
\end{enumerate}

\ans{（1）$\{1,2,3\}$；（2）$a=2$ 或 $\dfrac{3}{2}$，此时 $M=\{1,2\}$ 或 $M=\left\{\dfrac{1}{2},1,\dfrac{3}{2}\right\}$}

\ifanswer
\solbegin
\stepm{（1）}{$\{1,2,3\}$；}
\stepm{（2）}{当 $(x-a)(x^2-ax+a-1)=0$ 有重根时，重根只能算作集合的一个元素，}
\step{$(x-a)(x^2-ax+a-1)=(x-a)(x-1)[x-(a-1)]=0$，}
\step{当 $a=1$ 时，得 $M=\{0,1\}$，不符合题意；}
\step{当 $a-1=1$ 时，即 $a=2$ 时，得 $M=\{1,2\}$，符合题意；}
\step{当 $a\neq 1$ 且 $a\neq 2$ 时，此时 $M=\{1,a,a-1\}$，得 $1+a+a-1=3$，}
\step{解得 $a=\dfrac{3}{2}$，此时 $M=\left\{\dfrac{1}{2},1,\dfrac{3}{2}\right\}$，符合题意，}
\step{综上，实数 $a$ 的值为 $2$ 或 $\dfrac{3}{2}$，}
\step{当 $a=2$ 时，$M=\{1,2\}$；当 $a=\dfrac{3}{2}$ 时，$M=\left\{\dfrac{1}{2},1,\dfrac{3}{2}\right\}$.}
\solend

\fi

\ansspace[8]

\item 已知集合 $A=\{x\mid x^2-2x-3=0\}$，集合 $B=\{x\mid ax-1=0\}$. 若 $B\sub A$，求实数 $a$ 的值.

\ans{$a=0$ 或 $a=-1$ 或 $a=\dfrac{1}{3}$}

\ifanswer
\solbegin
\step{当 $B=\e$ 时，$a=0$，满足题意；}
\step{当 $B\neq\e$，即 $a\neq 0$ 时，$B=\left\{\dfrac{1}{a}\right\}$，}
\step{又 $A=\{x\mid x^2-2x-3=0\}=\{x\mid x=-1\text{ 或 }x=3\}$，$B\sub A$，}
\step{所以 $\dfrac{1}{a}=-1$ 或 $\dfrac{1}{a}=3$，所以 $a=-1$ 或 $a=\dfrac{1}{3}$.}
\step{综上所述，$a=0$ 或 $a=-1$ 或 $a=\dfrac{1}{3}$.}
\solend

\fi

\ansspace[8]

\item 设集合 $A=\{x\mid -2\leqslant x\leqslant 5\}$，$B=\{x\mid m+1\leqslant x\leqslant 2m-1\}$.

\begin{enumerate}[label=(\arabic*)]
\item 若 $B\sub A$，求实数 $m$ 的取值范围；
\item 当 $x\in\Z$ 时，求 $A$ 的非空真子集个数；
\item 当 $x\in\R$ 时，不存在元素 $x$ 使 $x\in A$ 与 $x\in B$ 同时成立，求实数 $m$ 的取值范围.
\end{enumerate}

\ans{（1）$\{m\mid m\leqslant 3\}$；（2）$254$；（3）$\{m\mid m<2\text{ 或 }m>4\}$}

\ifanswer
\solbegin
\stepm{（1）}{当 $m+1>2m-1$，即 $m<2$ 时，$B=\e$，满足 $B\sub A$.}
\step{当 $m+1\leqslant 2m-1$，即 $m\geqslant 2$ 时，要使 $B\sub A$ 成立，}
\step{需 $\begin{cases}m+1\geqslant -2\\ 2m-1\leqslant 5\end{cases}$，得 $2\leqslant m\leqslant 3$，}
\step{综上，$m$ 的范围为 $\{m\mid m\leqslant 3\}$；}
\stepm{（2）}{$A=\{x\mid -2\leqslant x\leqslant 5,x\in\Z\}=\{-2,-1,0,1,2,3,4,5\}$，}
\step{则 $A$ 的非空真子集个数为 $2^8-1-1=254$；}
\stepm{（3）}{因为没有元素 $x$ 使 $x\in A$ 与 $x\in B$ 同时成立，所以 $A$ 与 $B$ 交集为空集.}
\step{①若 $B=\e$，即 $m+1>2m-1$，得 $m<2$ 时满足条件；}
\step{②若 $B\neq\e$，则要满足的条件是 $\begin{cases}m+1\leqslant 2m-1\\ m+1>5\end{cases}$ 或 $\begin{cases}m+1\leqslant 2m-1\\ 2m-1<-2\end{cases}$，}
\step{解得 $m>4$.}
\step{综上，有 $m<2$ 或 $m>4$，所以 $m$ 的范围为 $\{m\mid m<2\text{ 或 }m>4\}$.}
\solend

\fi

\ansspace[10]

% 压轴题：其后已无内容，故作答区全用可丢弃胶水（\ansspace*），并在题目前先量一下版面。
\ansneed{7cm}

\item 已知集合 $A=\{x\mid ax=1\}$，$B=\{x\mid x^2-ax+b=0\}$.

\begin{enumerate}[label=(\arabic*)]
\item 若 $a=2$，且 $A\sub B$，求 $b$ 的值；
\item 当 $a^2\geqslant 4b$ 时，若 $B\sub A$，求 $a$、$b$ 的值；
\item 若 $A\sub B$，讨论 $a$、$b$ 的取值范围.
\end{enumerate}

\ans{（1）$\dfrac{3}{4}$；（2）$a=\pm\sqrt2$，$b=\dfrac{1}{2}$；（3）$a=0$ 时 $b\in\R$，$a\neq 0$ 时 $b\in(-\infty,1)$}

\ifanswer
\solbegin
\stepm{（1）}{当 $a=2$ 时，$A=\left\{\dfrac{1}{2}\right\}$，}
\step{由 $A\sub B$ 得 $\dfrac{1}{2}$ 是方程 $x^2-2x+b=0$ 的解，代入得 $\left(\dfrac{1}{2}\right)^2-2\times\dfrac{1}{2}+b=0$，}
\step{解得 $b=\dfrac{3}{4}$.}
\stepm{（2）}{当 $a^2\geqslant 4b$ 时，方程 $x^2-ax+b=0$ 有实根.}
\step{若 $a=0$，则 $A=\e$，要 $B\sub A$ 需 $B=\e$，}
\step{但此时方程 $x^2+b=0$ 在 $a^2\geqslant 4b$（即 $b\leqslant 0$）时有实根，矛盾，舍去.}
\step{若 $a\neq 0$，则 $A=\left\{\dfrac{1}{a}\right\}$，$B\sub A$ 且方程有实根，故方程有且仅有一个根 $\dfrac{1}{a}$.}
\step{由韦达定理，根的和为 $a=\dfrac{1}{a}+\dfrac{1}{a}$，根的积为 $b=\dfrac{1}{a}\times\dfrac{1}{a}$.}
\step{由根的和得 $a=\dfrac{2}{a}$，即 $a^2=2$，$a=\pm\sqrt2$，$b=\dfrac{1}{2}$.}
\stepm{（3）}{当 $a=0$ 时，$A=\e$，故 $A\sub B$ 恒成立，此时 $b\in\R$.}
\step{当 $a\neq 0$ 时，$A=\left\{\dfrac{1}{a}\right\}$，将 $x=\dfrac{1}{a}$ 代入 $x^2-ax+b=0$，得 $b=1-\dfrac{1}{a^2}$.}
\step{此时方程判别式 $\Delta=a^2-4b=\dfrac{(a^2-2)^2}{a^2}\geqslant 0$，}
\step{若 $\Delta=0$，则 $a=\pm\sqrt2$，$b=\dfrac{1}{2}$；}
\step{若 $\Delta>0$，则 $a\neq 0$ 且 $a\neq\pm\sqrt2$，此时 $b=1-\dfrac{1}{a^2}<1$.}
\step{综上，当 $a=0$ 时，$b\in\R$；当 $a\neq 0$ 时，$b\in(-\infty,1)$.}
\solend

\fi

\ansspace*[12]

\end{enumerate}

\end{document}
"
% 紧凑版：xelatex "\def\iscompact{1}% 高一10_复附浦东_周末作业1.tex —— 高一上数学“2026-2027 年复附浦东高一上周末作业 1”（试卷版/答案版）
% 试卷版：xelatex 高一10_复附浦东_周末作业1.tex
% 答案版：xelatex "\def\isanswer{1}\input{高一10_复附浦东_周末作业1.tex}"
% 紧凑版：xelatex "\def\iscompact{1}\input{高一10_复附浦东_周末作业1.tex}"
%
% 原始材料是微信公众号图文（公众号“魔都数学加油站”），正文为 9 张 PNG 页；原文本身就是一份
% **讲评版**——题目下方直接印出【答案】或【解析】。试题文字、答案与解析均照原卷录入，未改内容。
%
% 与原卷的排版差异（均已在 README“说明”中交代）：
%   ① 原文自**第 3 题**起（第 1、2 题未在该文中发布），故本稿题号亦自 3 起；
%   ② 原卷本就印有“一、填空题”“二、选择题”“三、解答题”三节标题（分别在原图 01/04/06.png
%      上，逐页放大核对过），本稿照原卷分节，只把标题改排成全稿统一的“大节标题”样式；
%   ③ 原卷填空、选择的答案直接印在题目下方，本稿统一改为试卷版隐去、答案版以【答案】给出，
%      使两版彻底分离；
%   ④ 原卷的实数集符号时作正体 $R$、时作粗体 $\mathbf{R}$（第 21(3) 与 (3) 末句均为粗体），
%      本稿按全稿惯例统一写作 $\R$.
%
% 原卷笔误一处（本稿已补，见该处上方的 `%` 注）：
%   ① 第 10 题【解析】首句原卷印作“这个元不是 $1$ 就是 $m$”，“元素”漏了“素”字——
%      本稿按 高一b 第 13 题“原卷漏排、本稿补上”的先例补出“素”字，其余照原卷。
% 除此一处外，本卷 9 页逐页放大核对，未发现别的笔误（题 16 的 $1995=15\times 133$ 与
% “$15^2$ 的倍数有 8 个”两处最易被误判为笔误，实测 $133$、$8$ 均正确）。

% 本篇在公众号“魔都数学加油站”连载里的编号是“27学年高一10”，本地编号与之对齐（高一10）；
% 对应关系见 README“与公众号连载号的对照表”。
\documentclass[a4paper,11pt]{article}
\usepackage{common}

\begin{document}

\examtitlest[3]{2026--2027 年复附浦东高一上周末作业 1}{集合与常用逻辑用语}

\bigsec{一、填空题}

\begin{enumerate}
\setcounter{enumi}{2}

\item 已知集合 $M=\{1,2m+1\}$，$N=\{-1,m^2\}$，且 $M=N$，则 $m=$ \blankline[2.4em].

\ans{$-1$}

\ifanswer
\solbegin
\step{若 $M=N$，必有 $\begin{cases}2m+1=-1\\ m^2=1\end{cases}$，解得 $m=-1$.}
\solend

\fi

\item 集合 $\{(x,y)\mid x+2y=2,\ x,y\in\N\}$ 用列举法表示为 \blankline[4.4em].

\ans{$\{(0,1),(2,0)\}$}

\ifanswer
\solbegin
\step{由题意得 $x+2y=2$ 可化为 $y=1-\dfrac{x}{2}$，}
\step{由 $y\in\N$，有 $y=1-\dfrac{x}{2}\geqslant 0$，解得 $x\leqslant 2$，}
\step{又由 $x\in\N$，得 $x$ 可能取值为 $0$，$1$，$2$，}
\step{当 $x=0$ 时，得 $y=1$，满足条件，}
\step{当 $x=1$ 时，得 $y=\dfrac{1}{2}\notin\N$，不满足条件，}
\step{当 $x=2$ 时，得 $y=0$，满足条件，}
\step{综上，集合用列举法表示为 $\{(0,1),(2,0)\}$.}
\solend

\fi

\item 已知集合 $A=\{1\}$，集合 $B=\{0,1,2\}$，则满足关系 $A\sub P\sub B$ 的所有集合 $P$ 为 \blankline[4.4em].

\ans{$\{1\},\{0,1\},\{1,2\},\{0,1,2\}$}

\ifanswer
\solbegin
\step{由题意，集合 $P$ 可以是 $\{1\},\{0,1\},\{1,2\},\{0,1,2\}$.}
\solend

\fi

\item 已知集合 $A=\{x\mid a-1\leqslant x\leqslant a+2\}$，$B=\{x\mid 3<x<5\}$，则使 $A\supseteq B$ 成立的实数 $a$ 的取值范围是 \blankline[4.4em].

\ans{$[3,4]$}

\ifanswer
\solbegin
\step{因为 $A\supseteq B$，所以 $\begin{cases}a+2\geqslant 5\\ a-1\leqslant 3\end{cases}$，解得 $3\leqslant a\leqslant 4$.}
\solend

\fi

\item 若集合 $\{x\mid mx^2-4x+1=0\}=\{n\}$，则 $m+n=$ \blankline[2.4em].

\ans{$\dfrac{1}{4}$ 或 $\dfrac{9}{2}$}

\ifanswer
\solbegin
\step{若集合 $\{x\mid mx^2-4x+1=0\}=\{n\}$，则 $mx^2-4x+1=0$ 有唯一实数根 $n$，}
\step{当 $m=0$ 时，得 $n=\dfrac{1}{4}$，$m+n=\dfrac{1}{4}$，}
\step{当 $m\neq 0$ 时，$\Delta=16-4m=0$，即 $m=4$，}
\step{此时 $4x^2-4x+1=0$ 的根为 $n=\dfrac{1}{2}$，$m+n=\dfrac{9}{2}$.}
\step{故 $m+n=\dfrac{1}{4}$ 或 $\dfrac{9}{2}$.}
\solend

\fi

\item 集合 $M=\{-1,2,3,5,6\}$ 的全部非空子集的元素和等于 \blankline[4.4em].

\ans{$240$}

\ifanswer
\solbegin
\step{集合 $M=\{-1,2,3,5,6\}$ 中共有 $5$ 个元素，}
\step{在其所有子集中，元素 $-1$ 出现的次数为 $2^{5-1}=2^4=16$，}
\step{同理 $2$、$3$、$5$、$6$ 都出现 $16$ 次，}
\step{则集合 $M$ 的全部非空子集的元素和为 $(-1+2+3+5+6)\times 16=240$.}
\solend

\fi

\item 设 $A=\{x\mid x^2-3x-4=0\}$，$B=\{x\mid ax-1=0\}$，若 $A\cap B=B$，则实数 $a$ 的值可以为 \blankline[4.4em].

\ans{$0$、$-1$、$\dfrac{1}{4}$}

\ifanswer
\solbegin
\step{$A=\{-1,4\}$，因为 $A\cap B=B$，所以 $B\sub A$，}
\step{①$a=0$ 时，$B=\e$，满足 $B\sub A$；}
\step{②$a\neq 0$ 时，$B=\left\{\dfrac{1}{a}\right\}$，则 $\dfrac{1}{a}=-1$ 或 $\dfrac{1}{a}=4$，所以 $a=-1$ 或 $\dfrac{1}{4}$；}
\step{所以 $a$ 的值可以为 $0$、$-1$、$\dfrac{1}{4}$.}
\solend

\fi

\item 设集合 $A=\{1,m\}$，其中 $m$ 为实数. 令集合 $B=\{a+b\mid a\in A,b\in A\}$，若 $A\cap B$ 恰有一个元素，则 $A\cup B$ 的元素之和为 \blankline[2.4em].

\ans{$5$ 或 $10$}

\ifanswer
\solbegin
% 注：原卷本句印作“这个元不是 $1$ 就是 $m$”——“元素”漏了“素”字。按 高一b 第 13 题
%     “原卷漏排、本稿补上”的先例，本稿补出“素”字，其余照原卷，见文件头“原卷笔误”。
\step{由题意得，若 $A\cap B$ 恰有一个元素，这个元素不是 $1$ 就是 $m$，}
\stepm{（1）}{当 $1\in B$ 时，}
\step{①若 $1=1+m$，则 $m=0$，此时 $A=\{1,0\},B=\{2,1,0\},A\cap B=\{0,1\}$，不满足题意，舍去；}
\step{②若 $1=2m$，则 $m=\dfrac{1}{2}$，此时 $A=\left\{1,\dfrac{1}{2}\right\}$，$B=\left\{2,\dfrac{3}{2},1\right\}$，$A\cap B=\{1\}$，}
\step{满足题意，所以 $A\cup B=\left\{1,\dfrac{1}{2},2,\dfrac{3}{2}\right\}$，其元素之和为 $5$，}
\stepm{（2）}{当 $m\in B$ 时，}
\step{①若 $m=1+m$，则 $m$ 无解，}
\step{②若 $m=2m$，则 $m=0$，此时 $A=\{1,0\},B=\{2,1,0\},A\cap B=\{0,1\}$，不满足题意，}
\step{③若 $m=2$，则 $A=\{1,2\},B=\{2,3,4\},A\cap B=\{2\}$，满足题意，}
\step{所以 $A\cup B=\{1,2,3,4\}$，其元素之和为 $10$，}
\step{综上所述，$A\cup B$ 的元素之和为 $5$ 或 $10$.}
\solend

\fi

\item 设非空集合 $Q\sub M$，当 $Q$ 中所有元素和为偶数时（集合为单元素时和为元素本身），称 $Q$ 是 $M$ 的偶子集. 若集合 $M=\{1,2,3,4,5,6,7\}$，则其偶子集 $Q$ 的个数为 \blankline[2.4em].

\ans{$63$}

\ifanswer
\solbegin
\step{偶子集可以分为只有奇数、偶数及偶数个奇数和偶数的集合，}
\step{奇数有 $4$ 个，偶数有 $3$ 个，}
\step{只有奇数的有 $\mathrm{C}_4^2+1=7$ 个，只有偶数的有 $\mathrm{C}_3^1+\mathrm{C}_3^2+1=7$ 个，}
\step{偶数个奇数和偶数的：两奇一偶 $\mathrm{C}_4^2\cdot\mathrm{C}_3^1=18$ 个，两奇二偶 $\mathrm{C}_4^2\cdot\mathrm{C}_3^2=18$ 个，}
\step{两奇三偶 $\mathrm{C}_4^2\cdot\mathrm{C}_3^3=6$ 个，四奇一偶 $\mathrm{C}_4^4\cdot\mathrm{C}_3^1=3$ 个，四奇二偶 $\mathrm{C}_4^4\cdot\mathrm{C}_3^2=3$ 个，}
\step{四奇三偶 $\mathrm{C}_4^4\cdot\mathrm{C}_3^3=1$ 个，所以共有 $7+7+18+18+6+3+3+1=63$ 个.}
\solend

\fi

\item 已知非空集合 $S=\left\{x\mid -\dfrac{1}{2}\leqslant x\leqslant m\right\}$ 满足：当 $x\in S$ 时，有 $x^2\in S$，则实数 $m$ 的取值范围是 \blankline[4.4em].

\ans{$\left[\dfrac{1}{4},1\right]$}

\ifanswer
\solbegin
\step{因为集合 $S=\left\{x\mid -\dfrac{1}{2}\leqslant x\leqslant m\right\}$ 是非空集合，所以 $m\geqslant -\dfrac{1}{2}$，}
\step{又当 $x\in S$ 时，有 $x^2\in S$，所以 $m^2\leqslant m$ 且 $m\geqslant \left(-\dfrac{1}{2}\right)^2=\dfrac{1}{4}$，所以 $\dfrac{1}{4}\leqslant m\leqslant 1$.}
\solend

\fi

\end{enumerate}

\bigsec{二、选择题}

\begin{enumerate}
\setcounter{enumi}{12}

\item 已知集合 $U=\{1+x,1\}$，$x\in U$，则 $x=$\mbox{（\quad）}

\keepwithprev
A. $-1$\qquad B. $0$\qquad C. $1$\qquad D. $2$

\ans{C}

\ifanswer
\solbegin
\step{因为集合 $U=\{1+x,1\}$，$x\in U$，所以 $1+x=x$ 或 $x=1$，}
\step{当 $x=1+x$ 时，即 $0=1$，此时显然不成立，舍去，}
\step{当 $x=1$，此时集合 $U=\{1,2\}$，满足互异性，}
\step{综上所述，$x=1$. 故选 C.}
\solend

\fi

\item 已知集合 $M=\{m\mid m=a+b\sqrt2,\ a,b\in\Q\}$，则下列四个元素中属于 $M$ 的元素的个数是\mbox{（\quad）}

①$1+\sqrt2\pi$；②$\sqrt{11+6\sqrt2}$；③$\dfrac{1}{2+\sqrt2}$；④$\sqrt{2-\sqrt3}+\sqrt{2+\sqrt3}$

\keepwithprev
A. $4$\qquad B. $3$\qquad C. $2$\qquad D. $1$

\ans{C}

\ifanswer
\solbegin
\step{①$1+\sqrt2\pi$，由于 $\sqrt2$、$\pi$ 都为无理数，因此 $1+\sqrt2\pi\notin M$；}
\step{②$\sqrt{11+6\sqrt2}=\sqrt{11+2\sqrt{18}}=3+\sqrt2\in M$；}
\step{③$\dfrac{1}{2+\sqrt2}=1-\dfrac{1}{2}\sqrt2\in M$；}
\step{④$\sqrt{2-\sqrt3}+\sqrt{2+\sqrt3}=\sqrt{\dfrac{8-2\sqrt{12}}{4}}+\sqrt{\dfrac{8+2\sqrt{12}}{4}}=\dfrac{\sqrt6-\sqrt2}{2}+\dfrac{\sqrt6+\sqrt2}{2}=\sqrt6=\sqrt3\cdot\sqrt2\notin M$；}
\step{则上述四个元素中属于 $M$ 的元素的个数是 $2$. 故选 C.}
\solend

\fi

\item 集合 $A=\{2,3,7,8,9,11\}$，集合 $B=\{1,2,4,6,8\}$，下面说法正确的是\mbox{（\quad）}

\keepwithprev
A. 两个集合的交集是 $\{2,8\}$

B. 两个集合没有并集

C. $B$ 集合的元素无法用 $\{x\mid x=f(k),k\in D\}$ 这种形式表示出来

D. 集合 $A$ 有 $6$ 个子集

\ans{A}

\ifanswer
\solbegin
\step{对于 A，$A\cap B=\{2,8\}$，故 A 正确；}
\step{对于 B，$A\cup B=\{1,2,3,4,6,7,8,9,11\}$，故 B 错误；}
\step{对于 C，集合 $B$ 可以表示为 $B=\{x\mid x=f(k),k\in\{1,2,3,4,5\}\}$，}
\step{其中 $f(1)=1,f(2)=2,f(3)=4,f(4)=6,f(5)=8$，}
\step{即集合 $B$ 可描述为 $\{x\mid x=f(k),k\in D\}$ 形式，故 C 错误；}
\step{对于 D，集合 $A$ 有 $6$ 个元素，其子集总数为 $2^6=64\neq 6$，故 D 错误.}
\step{故选 A.}
\solend

\fi

\item 设 $M=\{1,2,\cdots,1995\}$，$A$ 是 $M$ 的子集且满足条件：当 $x\in A$ 时，$15x\notin A$，则 $A$ 中元素的个数的最大值为\mbox{（\quad）}

\keepwithprev
A. $1862$\qquad B. $1866$\qquad C. $1868$\qquad D. $1870$

\ans{D}

\ifanswer
\solbegin
\step{$1995=15\times 133$. 故取出所有不是 $15$ 的倍数的数，共 $1862$ 个，}
\step{这些数均符合要求.}
\step{在所有 $15$ 的倍数的数中，$15^2$ 的倍数有 $8$ 个，这些数又可以取出，}
\step{这样共取出了 $1870$ 个，即 $|A|\geqslant 1870$.}
\step{又 $\{k,15k\}(k=9,10,11,\cdots,133)$ 中的两个元素不能同时取出，}
\step{故 $|A|\leqslant 1995-133+8=1870$，故选 D.}
\solend

\fi

\end{enumerate}

\bigsec{三、解答题}

\begin{enumerate}
\setcounter{enumi}{16}

\item 已知集合 $A=\{x\mid a\leqslant x\leqslant a+4\}$，$B=\{x\mid x\geqslant 6\}$.

\begin{enumerate}[label=(\arabic*)]
\item 若 $a=3$，求 $A\cup B$；
\item 若 $A\cap B=A$，求实数 $a$ 的取值范围.
\end{enumerate}

\ans{（1）$\{x\mid x\geqslant 3\}$；（2）$[6,+\infty)$}

\ifanswer
\solbegin
\stepm{（1）}{当 $a=3$ 时，$A=\{x\mid 3\leqslant x\leqslant 7\}$，所以 $A\cup B=\{x\mid x\geqslant 3\}$；}
\stepm{（2）}{若 $A\cap B=A$，则 $A\sub B$，所以 $a\geqslant 6$，所以实数 $a$ 的取值范围是 $[6,+\infty)$.}
\solend

\fi

\ansspace[7]

\item 已知集合 $M=\{x\mid (x-a)(x^2-ax+a-1)=0\}$.

\begin{enumerate}[label=(\arabic*)]
\item 若 $a=3$，求集合 $M$；
\item 若集合 $M$ 中各元素之和等于 $3$，求实数 $a$ 的值，并用列举法表示集合 $M$.
\end{enumerate}

\ans{（1）$\{1,2,3\}$；（2）$a=2$ 或 $\dfrac{3}{2}$，此时 $M=\{1,2\}$ 或 $M=\left\{\dfrac{1}{2},1,\dfrac{3}{2}\right\}$}

\ifanswer
\solbegin
\stepm{（1）}{$\{1,2,3\}$；}
\stepm{（2）}{当 $(x-a)(x^2-ax+a-1)=0$ 有重根时，重根只能算作集合的一个元素，}
\step{$(x-a)(x^2-ax+a-1)=(x-a)(x-1)[x-(a-1)]=0$，}
\step{当 $a=1$ 时，得 $M=\{0,1\}$，不符合题意；}
\step{当 $a-1=1$ 时，即 $a=2$ 时，得 $M=\{1,2\}$，符合题意；}
\step{当 $a\neq 1$ 且 $a\neq 2$ 时，此时 $M=\{1,a,a-1\}$，得 $1+a+a-1=3$，}
\step{解得 $a=\dfrac{3}{2}$，此时 $M=\left\{\dfrac{1}{2},1,\dfrac{3}{2}\right\}$，符合题意，}
\step{综上，实数 $a$ 的值为 $2$ 或 $\dfrac{3}{2}$，}
\step{当 $a=2$ 时，$M=\{1,2\}$；当 $a=\dfrac{3}{2}$ 时，$M=\left\{\dfrac{1}{2},1,\dfrac{3}{2}\right\}$.}
\solend

\fi

\ansspace[8]

\item 已知集合 $A=\{x\mid x^2-2x-3=0\}$，集合 $B=\{x\mid ax-1=0\}$. 若 $B\sub A$，求实数 $a$ 的值.

\ans{$a=0$ 或 $a=-1$ 或 $a=\dfrac{1}{3}$}

\ifanswer
\solbegin
\step{当 $B=\e$ 时，$a=0$，满足题意；}
\step{当 $B\neq\e$，即 $a\neq 0$ 时，$B=\left\{\dfrac{1}{a}\right\}$，}
\step{又 $A=\{x\mid x^2-2x-3=0\}=\{x\mid x=-1\text{ 或 }x=3\}$，$B\sub A$，}
\step{所以 $\dfrac{1}{a}=-1$ 或 $\dfrac{1}{a}=3$，所以 $a=-1$ 或 $a=\dfrac{1}{3}$.}
\step{综上所述，$a=0$ 或 $a=-1$ 或 $a=\dfrac{1}{3}$.}
\solend

\fi

\ansspace[8]

\item 设集合 $A=\{x\mid -2\leqslant x\leqslant 5\}$，$B=\{x\mid m+1\leqslant x\leqslant 2m-1\}$.

\begin{enumerate}[label=(\arabic*)]
\item 若 $B\sub A$，求实数 $m$ 的取值范围；
\item 当 $x\in\Z$ 时，求 $A$ 的非空真子集个数；
\item 当 $x\in\R$ 时，不存在元素 $x$ 使 $x\in A$ 与 $x\in B$ 同时成立，求实数 $m$ 的取值范围.
\end{enumerate}

\ans{（1）$\{m\mid m\leqslant 3\}$；（2）$254$；（3）$\{m\mid m<2\text{ 或 }m>4\}$}

\ifanswer
\solbegin
\stepm{（1）}{当 $m+1>2m-1$，即 $m<2$ 时，$B=\e$，满足 $B\sub A$.}
\step{当 $m+1\leqslant 2m-1$，即 $m\geqslant 2$ 时，要使 $B\sub A$ 成立，}
\step{需 $\begin{cases}m+1\geqslant -2\\ 2m-1\leqslant 5\end{cases}$，得 $2\leqslant m\leqslant 3$，}
\step{综上，$m$ 的范围为 $\{m\mid m\leqslant 3\}$；}
\stepm{（2）}{$A=\{x\mid -2\leqslant x\leqslant 5,x\in\Z\}=\{-2,-1,0,1,2,3,4,5\}$，}
\step{则 $A$ 的非空真子集个数为 $2^8-1-1=254$；}
\stepm{（3）}{因为没有元素 $x$ 使 $x\in A$ 与 $x\in B$ 同时成立，所以 $A$ 与 $B$ 交集为空集.}
\step{①若 $B=\e$，即 $m+1>2m-1$，得 $m<2$ 时满足条件；}
\step{②若 $B\neq\e$，则要满足的条件是 $\begin{cases}m+1\leqslant 2m-1\\ m+1>5\end{cases}$ 或 $\begin{cases}m+1\leqslant 2m-1\\ 2m-1<-2\end{cases}$，}
\step{解得 $m>4$.}
\step{综上，有 $m<2$ 或 $m>4$，所以 $m$ 的范围为 $\{m\mid m<2\text{ 或 }m>4\}$.}
\solend

\fi

\ansspace[10]

% 压轴题：其后已无内容，故作答区全用可丢弃胶水（\ansspace*），并在题目前先量一下版面。
\ansneed{7cm}

\item 已知集合 $A=\{x\mid ax=1\}$，$B=\{x\mid x^2-ax+b=0\}$.

\begin{enumerate}[label=(\arabic*)]
\item 若 $a=2$，且 $A\sub B$，求 $b$ 的值；
\item 当 $a^2\geqslant 4b$ 时，若 $B\sub A$，求 $a$、$b$ 的值；
\item 若 $A\sub B$，讨论 $a$、$b$ 的取值范围.
\end{enumerate}

\ans{（1）$\dfrac{3}{4}$；（2）$a=\pm\sqrt2$，$b=\dfrac{1}{2}$；（3）$a=0$ 时 $b\in\R$，$a\neq 0$ 时 $b\in(-\infty,1)$}

\ifanswer
\solbegin
\stepm{（1）}{当 $a=2$ 时，$A=\left\{\dfrac{1}{2}\right\}$，}
\step{由 $A\sub B$ 得 $\dfrac{1}{2}$ 是方程 $x^2-2x+b=0$ 的解，代入得 $\left(\dfrac{1}{2}\right)^2-2\times\dfrac{1}{2}+b=0$，}
\step{解得 $b=\dfrac{3}{4}$.}
\stepm{（2）}{当 $a^2\geqslant 4b$ 时，方程 $x^2-ax+b=0$ 有实根.}
\step{若 $a=0$，则 $A=\e$，要 $B\sub A$ 需 $B=\e$，}
\step{但此时方程 $x^2+b=0$ 在 $a^2\geqslant 4b$（即 $b\leqslant 0$）时有实根，矛盾，舍去.}
\step{若 $a\neq 0$，则 $A=\left\{\dfrac{1}{a}\right\}$，$B\sub A$ 且方程有实根，故方程有且仅有一个根 $\dfrac{1}{a}$.}
\step{由韦达定理，根的和为 $a=\dfrac{1}{a}+\dfrac{1}{a}$，根的积为 $b=\dfrac{1}{a}\times\dfrac{1}{a}$.}
\step{由根的和得 $a=\dfrac{2}{a}$，即 $a^2=2$，$a=\pm\sqrt2$，$b=\dfrac{1}{2}$.}
\stepm{（3）}{当 $a=0$ 时，$A=\e$，故 $A\sub B$ 恒成立，此时 $b\in\R$.}
\step{当 $a\neq 0$ 时，$A=\left\{\dfrac{1}{a}\right\}$，将 $x=\dfrac{1}{a}$ 代入 $x^2-ax+b=0$，得 $b=1-\dfrac{1}{a^2}$.}
\step{此时方程判别式 $\Delta=a^2-4b=\dfrac{(a^2-2)^2}{a^2}\geqslant 0$，}
\step{若 $\Delta=0$，则 $a=\pm\sqrt2$，$b=\dfrac{1}{2}$；}
\step{若 $\Delta>0$，则 $a\neq 0$ 且 $a\neq\pm\sqrt2$，此时 $b=1-\dfrac{1}{a^2}<1$.}
\step{综上，当 $a=0$ 时，$b\in\R$；当 $a\neq 0$ 时，$b\in(-\infty,1)$.}
\solend

\fi

\ansspace*[12]

\end{enumerate}

\end{document}
"
%
% 原始材料是微信公众号图文（公众号“魔都数学加油站”），正文为 9 张 PNG 页；原文本身就是一份
% **讲评版**——题目下方直接印出【答案】或【解析】。试题文字、答案与解析均照原卷录入，未改内容。
%
% 与原卷的排版差异（均已在 README“说明”中交代）：
%   ① 原文自**第 3 题**起（第 1、2 题未在该文中发布），故本稿题号亦自 3 起；
%   ② 原卷本就印有“一、填空题”“二、选择题”“三、解答题”三节标题（分别在原图 01/04/06.png
%      上，逐页放大核对过），本稿照原卷分节，只把标题改排成全稿统一的“大节标题”样式；
%   ③ 原卷填空、选择的答案直接印在题目下方，本稿统一改为试卷版隐去、答案版以【答案】给出，
%      使两版彻底分离；
%   ④ 原卷的实数集符号时作正体 $R$、时作粗体 $\mathbf{R}$（第 21(3) 与 (3) 末句均为粗体），
%      本稿按全稿惯例统一写作 $\R$.
%
% 原卷笔误一处（本稿已补，见该处上方的 `%` 注）：
%   ① 第 10 题【解析】首句原卷印作“这个元不是 $1$ 就是 $m$”，“元素”漏了“素”字——
%      本稿按 高一b 第 13 题“原卷漏排、本稿补上”的先例补出“素”字，其余照原卷。
% 除此一处外，本卷 9 页逐页放大核对，未发现别的笔误（题 16 的 $1995=15\times 133$ 与
% “$15^2$ 的倍数有 8 个”两处最易被误判为笔误，实测 $133$、$8$ 均正确）。

% 本篇在公众号“魔都数学加油站”连载里的编号是“27学年高一10”，本地编号与之对齐（高一10）；
% 对应关系见 README“与公众号连载号的对照表”。
\documentclass[a4paper,11pt]{article}
\usepackage{common}

\begin{document}

\examtitlest[3]{2026--2027 年复附浦东高一上周末作业 1}{集合与常用逻辑用语}

\bigsec{一、填空题}

\begin{enumerate}
\setcounter{enumi}{2}

\item 已知集合 $M=\{1,2m+1\}$，$N=\{-1,m^2\}$，且 $M=N$，则 $m=$ \blankline[2.4em].

\ans{$-1$}

\ifanswer
\solbegin
\step{若 $M=N$，必有 $\begin{cases}2m+1=-1\\ m^2=1\end{cases}$，解得 $m=-1$.}
\solend

\fi

\item 集合 $\{(x,y)\mid x+2y=2,\ x,y\in\N\}$ 用列举法表示为 \blankline[4.4em].

\ans{$\{(0,1),(2,0)\}$}

\ifanswer
\solbegin
\step{由题意得 $x+2y=2$ 可化为 $y=1-\dfrac{x}{2}$，}
\step{由 $y\in\N$，有 $y=1-\dfrac{x}{2}\geqslant 0$，解得 $x\leqslant 2$，}
\step{又由 $x\in\N$，得 $x$ 可能取值为 $0$，$1$，$2$，}
\step{当 $x=0$ 时，得 $y=1$，满足条件，}
\step{当 $x=1$ 时，得 $y=\dfrac{1}{2}\notin\N$，不满足条件，}
\step{当 $x=2$ 时，得 $y=0$，满足条件，}
\step{综上，集合用列举法表示为 $\{(0,1),(2,0)\}$.}
\solend

\fi

\item 已知集合 $A=\{1\}$，集合 $B=\{0,1,2\}$，则满足关系 $A\sub P\sub B$ 的所有集合 $P$ 为 \blankline[4.4em].

\ans{$\{1\},\{0,1\},\{1,2\},\{0,1,2\}$}

\ifanswer
\solbegin
\step{由题意，集合 $P$ 可以是 $\{1\},\{0,1\},\{1,2\},\{0,1,2\}$.}
\solend

\fi

\item 已知集合 $A=\{x\mid a-1\leqslant x\leqslant a+2\}$，$B=\{x\mid 3<x<5\}$，则使 $A\supseteq B$ 成立的实数 $a$ 的取值范围是 \blankline[4.4em].

\ans{$[3,4]$}

\ifanswer
\solbegin
\step{因为 $A\supseteq B$，所以 $\begin{cases}a+2\geqslant 5\\ a-1\leqslant 3\end{cases}$，解得 $3\leqslant a\leqslant 4$.}
\solend

\fi

\item 若集合 $\{x\mid mx^2-4x+1=0\}=\{n\}$，则 $m+n=$ \blankline[2.4em].

\ans{$\dfrac{1}{4}$ 或 $\dfrac{9}{2}$}

\ifanswer
\solbegin
\step{若集合 $\{x\mid mx^2-4x+1=0\}=\{n\}$，则 $mx^2-4x+1=0$ 有唯一实数根 $n$，}
\step{当 $m=0$ 时，得 $n=\dfrac{1}{4}$，$m+n=\dfrac{1}{4}$，}
\step{当 $m\neq 0$ 时，$\Delta=16-4m=0$，即 $m=4$，}
\step{此时 $4x^2-4x+1=0$ 的根为 $n=\dfrac{1}{2}$，$m+n=\dfrac{9}{2}$.}
\step{故 $m+n=\dfrac{1}{4}$ 或 $\dfrac{9}{2}$.}
\solend

\fi

\item 集合 $M=\{-1,2,3,5,6\}$ 的全部非空子集的元素和等于 \blankline[4.4em].

\ans{$240$}

\ifanswer
\solbegin
\step{集合 $M=\{-1,2,3,5,6\}$ 中共有 $5$ 个元素，}
\step{在其所有子集中，元素 $-1$ 出现的次数为 $2^{5-1}=2^4=16$，}
\step{同理 $2$、$3$、$5$、$6$ 都出现 $16$ 次，}
\step{则集合 $M$ 的全部非空子集的元素和为 $(-1+2+3+5+6)\times 16=240$.}
\solend

\fi

\item 设 $A=\{x\mid x^2-3x-4=0\}$，$B=\{x\mid ax-1=0\}$，若 $A\cap B=B$，则实数 $a$ 的值可以为 \blankline[4.4em].

\ans{$0$、$-1$、$\dfrac{1}{4}$}

\ifanswer
\solbegin
\step{$A=\{-1,4\}$，因为 $A\cap B=B$，所以 $B\sub A$，}
\step{①$a=0$ 时，$B=\e$，满足 $B\sub A$；}
\step{②$a\neq 0$ 时，$B=\left\{\dfrac{1}{a}\right\}$，则 $\dfrac{1}{a}=-1$ 或 $\dfrac{1}{a}=4$，所以 $a=-1$ 或 $\dfrac{1}{4}$；}
\step{所以 $a$ 的值可以为 $0$、$-1$、$\dfrac{1}{4}$.}
\solend

\fi

\item 设集合 $A=\{1,m\}$，其中 $m$ 为实数. 令集合 $B=\{a+b\mid a\in A,b\in A\}$，若 $A\cap B$ 恰有一个元素，则 $A\cup B$ 的元素之和为 \blankline[2.4em].

\ans{$5$ 或 $10$}

\ifanswer
\solbegin
% 注：原卷本句印作“这个元不是 $1$ 就是 $m$”——“元素”漏了“素”字。按 高一b 第 13 题
%     “原卷漏排、本稿补上”的先例，本稿补出“素”字，其余照原卷，见文件头“原卷笔误”。
\step{由题意得，若 $A\cap B$ 恰有一个元素，这个元素不是 $1$ 就是 $m$，}
\stepm{（1）}{当 $1\in B$ 时，}
\step{①若 $1=1+m$，则 $m=0$，此时 $A=\{1,0\},B=\{2,1,0\},A\cap B=\{0,1\}$，不满足题意，舍去；}
\step{②若 $1=2m$，则 $m=\dfrac{1}{2}$，此时 $A=\left\{1,\dfrac{1}{2}\right\}$，$B=\left\{2,\dfrac{3}{2},1\right\}$，$A\cap B=\{1\}$，}
\step{满足题意，所以 $A\cup B=\left\{1,\dfrac{1}{2},2,\dfrac{3}{2}\right\}$，其元素之和为 $5$，}
\stepm{（2）}{当 $m\in B$ 时，}
\step{①若 $m=1+m$，则 $m$ 无解，}
\step{②若 $m=2m$，则 $m=0$，此时 $A=\{1,0\},B=\{2,1,0\},A\cap B=\{0,1\}$，不满足题意，}
\step{③若 $m=2$，则 $A=\{1,2\},B=\{2,3,4\},A\cap B=\{2\}$，满足题意，}
\step{所以 $A\cup B=\{1,2,3,4\}$，其元素之和为 $10$，}
\step{综上所述，$A\cup B$ 的元素之和为 $5$ 或 $10$.}
\solend

\fi

\item 设非空集合 $Q\sub M$，当 $Q$ 中所有元素和为偶数时（集合为单元素时和为元素本身），称 $Q$ 是 $M$ 的偶子集. 若集合 $M=\{1,2,3,4,5,6,7\}$，则其偶子集 $Q$ 的个数为 \blankline[2.4em].

\ans{$63$}

\ifanswer
\solbegin
\step{偶子集可以分为只有奇数、偶数及偶数个奇数和偶数的集合，}
\step{奇数有 $4$ 个，偶数有 $3$ 个，}
\step{只有奇数的有 $\mathrm{C}_4^2+1=7$ 个，只有偶数的有 $\mathrm{C}_3^1+\mathrm{C}_3^2+1=7$ 个，}
\step{偶数个奇数和偶数的：两奇一偶 $\mathrm{C}_4^2\cdot\mathrm{C}_3^1=18$ 个，两奇二偶 $\mathrm{C}_4^2\cdot\mathrm{C}_3^2=18$ 个，}
\step{两奇三偶 $\mathrm{C}_4^2\cdot\mathrm{C}_3^3=6$ 个，四奇一偶 $\mathrm{C}_4^4\cdot\mathrm{C}_3^1=3$ 个，四奇二偶 $\mathrm{C}_4^4\cdot\mathrm{C}_3^2=3$ 个，}
\step{四奇三偶 $\mathrm{C}_4^4\cdot\mathrm{C}_3^3=1$ 个，所以共有 $7+7+18+18+6+3+3+1=63$ 个.}
\solend

\fi

\item 已知非空集合 $S=\left\{x\mid -\dfrac{1}{2}\leqslant x\leqslant m\right\}$ 满足：当 $x\in S$ 时，有 $x^2\in S$，则实数 $m$ 的取值范围是 \blankline[4.4em].

\ans{$\left[\dfrac{1}{4},1\right]$}

\ifanswer
\solbegin
\step{因为集合 $S=\left\{x\mid -\dfrac{1}{2}\leqslant x\leqslant m\right\}$ 是非空集合，所以 $m\geqslant -\dfrac{1}{2}$，}
\step{又当 $x\in S$ 时，有 $x^2\in S$，所以 $m^2\leqslant m$ 且 $m\geqslant \left(-\dfrac{1}{2}\right)^2=\dfrac{1}{4}$，所以 $\dfrac{1}{4}\leqslant m\leqslant 1$.}
\solend

\fi

\end{enumerate}

\bigsec{二、选择题}

\begin{enumerate}
\setcounter{enumi}{12}

\item 已知集合 $U=\{1+x,1\}$，$x\in U$，则 $x=$\mbox{（\quad）}

\keepwithprev
A. $-1$\qquad B. $0$\qquad C. $1$\qquad D. $2$

\ans{C}

\ifanswer
\solbegin
\step{因为集合 $U=\{1+x,1\}$，$x\in U$，所以 $1+x=x$ 或 $x=1$，}
\step{当 $x=1+x$ 时，即 $0=1$，此时显然不成立，舍去，}
\step{当 $x=1$，此时集合 $U=\{1,2\}$，满足互异性，}
\step{综上所述，$x=1$. 故选 C.}
\solend

\fi

\item 已知集合 $M=\{m\mid m=a+b\sqrt2,\ a,b\in\Q\}$，则下列四个元素中属于 $M$ 的元素的个数是\mbox{（\quad）}

①$1+\sqrt2\pi$；②$\sqrt{11+6\sqrt2}$；③$\dfrac{1}{2+\sqrt2}$；④$\sqrt{2-\sqrt3}+\sqrt{2+\sqrt3}$

\keepwithprev
A. $4$\qquad B. $3$\qquad C. $2$\qquad D. $1$

\ans{C}

\ifanswer
\solbegin
\step{①$1+\sqrt2\pi$，由于 $\sqrt2$、$\pi$ 都为无理数，因此 $1+\sqrt2\pi\notin M$；}
\step{②$\sqrt{11+6\sqrt2}=\sqrt{11+2\sqrt{18}}=3+\sqrt2\in M$；}
\step{③$\dfrac{1}{2+\sqrt2}=1-\dfrac{1}{2}\sqrt2\in M$；}
\step{④$\sqrt{2-\sqrt3}+\sqrt{2+\sqrt3}=\sqrt{\dfrac{8-2\sqrt{12}}{4}}+\sqrt{\dfrac{8+2\sqrt{12}}{4}}=\dfrac{\sqrt6-\sqrt2}{2}+\dfrac{\sqrt6+\sqrt2}{2}=\sqrt6=\sqrt3\cdot\sqrt2\notin M$；}
\step{则上述四个元素中属于 $M$ 的元素的个数是 $2$. 故选 C.}
\solend

\fi

\item 集合 $A=\{2,3,7,8,9,11\}$，集合 $B=\{1,2,4,6,8\}$，下面说法正确的是\mbox{（\quad）}

\keepwithprev
A. 两个集合的交集是 $\{2,8\}$

B. 两个集合没有并集

C. $B$ 集合的元素无法用 $\{x\mid x=f(k),k\in D\}$ 这种形式表示出来

D. 集合 $A$ 有 $6$ 个子集

\ans{A}

\ifanswer
\solbegin
\step{对于 A，$A\cap B=\{2,8\}$，故 A 正确；}
\step{对于 B，$A\cup B=\{1,2,3,4,6,7,8,9,11\}$，故 B 错误；}
\step{对于 C，集合 $B$ 可以表示为 $B=\{x\mid x=f(k),k\in\{1,2,3,4,5\}\}$，}
\step{其中 $f(1)=1,f(2)=2,f(3)=4,f(4)=6,f(5)=8$，}
\step{即集合 $B$ 可描述为 $\{x\mid x=f(k),k\in D\}$ 形式，故 C 错误；}
\step{对于 D，集合 $A$ 有 $6$ 个元素，其子集总数为 $2^6=64\neq 6$，故 D 错误.}
\step{故选 A.}
\solend

\fi

\item 设 $M=\{1,2,\cdots,1995\}$，$A$ 是 $M$ 的子集且满足条件：当 $x\in A$ 时，$15x\notin A$，则 $A$ 中元素的个数的最大值为\mbox{（\quad）}

\keepwithprev
A. $1862$\qquad B. $1866$\qquad C. $1868$\qquad D. $1870$

\ans{D}

\ifanswer
\solbegin
\step{$1995=15\times 133$. 故取出所有不是 $15$ 的倍数的数，共 $1862$ 个，}
\step{这些数均符合要求.}
\step{在所有 $15$ 的倍数的数中，$15^2$ 的倍数有 $8$ 个，这些数又可以取出，}
\step{这样共取出了 $1870$ 个，即 $|A|\geqslant 1870$.}
\step{又 $\{k,15k\}(k=9,10,11,\cdots,133)$ 中的两个元素不能同时取出，}
\step{故 $|A|\leqslant 1995-133+8=1870$，故选 D.}
\solend

\fi

\end{enumerate}

\bigsec{三、解答题}

\begin{enumerate}
\setcounter{enumi}{16}

\item 已知集合 $A=\{x\mid a\leqslant x\leqslant a+4\}$，$B=\{x\mid x\geqslant 6\}$.

\begin{enumerate}[label=(\arabic*)]
\item 若 $a=3$，求 $A\cup B$；
\item 若 $A\cap B=A$，求实数 $a$ 的取值范围.
\end{enumerate}

\ans{（1）$\{x\mid x\geqslant 3\}$；（2）$[6,+\infty)$}

\ifanswer
\solbegin
\stepm{（1）}{当 $a=3$ 时，$A=\{x\mid 3\leqslant x\leqslant 7\}$，所以 $A\cup B=\{x\mid x\geqslant 3\}$；}
\stepm{（2）}{若 $A\cap B=A$，则 $A\sub B$，所以 $a\geqslant 6$，所以实数 $a$ 的取值范围是 $[6,+\infty)$.}
\solend

\fi

\ansspace[7]

\item 已知集合 $M=\{x\mid (x-a)(x^2-ax+a-1)=0\}$.

\begin{enumerate}[label=(\arabic*)]
\item 若 $a=3$，求集合 $M$；
\item 若集合 $M$ 中各元素之和等于 $3$，求实数 $a$ 的值，并用列举法表示集合 $M$.
\end{enumerate}

\ans{（1）$\{1,2,3\}$；（2）$a=2$ 或 $\dfrac{3}{2}$，此时 $M=\{1,2\}$ 或 $M=\left\{\dfrac{1}{2},1,\dfrac{3}{2}\right\}$}

\ifanswer
\solbegin
\stepm{（1）}{$\{1,2,3\}$；}
\stepm{（2）}{当 $(x-a)(x^2-ax+a-1)=0$ 有重根时，重根只能算作集合的一个元素，}
\step{$(x-a)(x^2-ax+a-1)=(x-a)(x-1)[x-(a-1)]=0$，}
\step{当 $a=1$ 时，得 $M=\{0,1\}$，不符合题意；}
\step{当 $a-1=1$ 时，即 $a=2$ 时，得 $M=\{1,2\}$，符合题意；}
\step{当 $a\neq 1$ 且 $a\neq 2$ 时，此时 $M=\{1,a,a-1\}$，得 $1+a+a-1=3$，}
\step{解得 $a=\dfrac{3}{2}$，此时 $M=\left\{\dfrac{1}{2},1,\dfrac{3}{2}\right\}$，符合题意，}
\step{综上，实数 $a$ 的值为 $2$ 或 $\dfrac{3}{2}$，}
\step{当 $a=2$ 时，$M=\{1,2\}$；当 $a=\dfrac{3}{2}$ 时，$M=\left\{\dfrac{1}{2},1,\dfrac{3}{2}\right\}$.}
\solend

\fi

\ansspace[8]

\item 已知集合 $A=\{x\mid x^2-2x-3=0\}$，集合 $B=\{x\mid ax-1=0\}$. 若 $B\sub A$，求实数 $a$ 的值.

\ans{$a=0$ 或 $a=-1$ 或 $a=\dfrac{1}{3}$}

\ifanswer
\solbegin
\step{当 $B=\e$ 时，$a=0$，满足题意；}
\step{当 $B\neq\e$，即 $a\neq 0$ 时，$B=\left\{\dfrac{1}{a}\right\}$，}
\step{又 $A=\{x\mid x^2-2x-3=0\}=\{x\mid x=-1\text{ 或 }x=3\}$，$B\sub A$，}
\step{所以 $\dfrac{1}{a}=-1$ 或 $\dfrac{1}{a}=3$，所以 $a=-1$ 或 $a=\dfrac{1}{3}$.}
\step{综上所述，$a=0$ 或 $a=-1$ 或 $a=\dfrac{1}{3}$.}
\solend

\fi

\ansspace[8]

\item 设集合 $A=\{x\mid -2\leqslant x\leqslant 5\}$，$B=\{x\mid m+1\leqslant x\leqslant 2m-1\}$.

\begin{enumerate}[label=(\arabic*)]
\item 若 $B\sub A$，求实数 $m$ 的取值范围；
\item 当 $x\in\Z$ 时，求 $A$ 的非空真子集个数；
\item 当 $x\in\R$ 时，不存在元素 $x$ 使 $x\in A$ 与 $x\in B$ 同时成立，求实数 $m$ 的取值范围.
\end{enumerate}

\ans{（1）$\{m\mid m\leqslant 3\}$；（2）$254$；（3）$\{m\mid m<2\text{ 或 }m>4\}$}

\ifanswer
\solbegin
\stepm{（1）}{当 $m+1>2m-1$，即 $m<2$ 时，$B=\e$，满足 $B\sub A$.}
\step{当 $m+1\leqslant 2m-1$，即 $m\geqslant 2$ 时，要使 $B\sub A$ 成立，}
\step{需 $\begin{cases}m+1\geqslant -2\\ 2m-1\leqslant 5\end{cases}$，得 $2\leqslant m\leqslant 3$，}
\step{综上，$m$ 的范围为 $\{m\mid m\leqslant 3\}$；}
\stepm{（2）}{$A=\{x\mid -2\leqslant x\leqslant 5,x\in\Z\}=\{-2,-1,0,1,2,3,4,5\}$，}
\step{则 $A$ 的非空真子集个数为 $2^8-1-1=254$；}
\stepm{（3）}{因为没有元素 $x$ 使 $x\in A$ 与 $x\in B$ 同时成立，所以 $A$ 与 $B$ 交集为空集.}
\step{①若 $B=\e$，即 $m+1>2m-1$，得 $m<2$ 时满足条件；}
\step{②若 $B\neq\e$，则要满足的条件是 $\begin{cases}m+1\leqslant 2m-1\\ m+1>5\end{cases}$ 或 $\begin{cases}m+1\leqslant 2m-1\\ 2m-1<-2\end{cases}$，}
\step{解得 $m>4$.}
\step{综上，有 $m<2$ 或 $m>4$，所以 $m$ 的范围为 $\{m\mid m<2\text{ 或 }m>4\}$.}
\solend

\fi

\ansspace[10]

% 压轴题：其后已无内容，故作答区全用可丢弃胶水（\ansspace*），并在题目前先量一下版面。
\ansneed{7cm}

\item 已知集合 $A=\{x\mid ax=1\}$，$B=\{x\mid x^2-ax+b=0\}$.

\begin{enumerate}[label=(\arabic*)]
\item 若 $a=2$，且 $A\sub B$，求 $b$ 的值；
\item 当 $a^2\geqslant 4b$ 时，若 $B\sub A$，求 $a$、$b$ 的值；
\item 若 $A\sub B$，讨论 $a$、$b$ 的取值范围.
\end{enumerate}

\ans{（1）$\dfrac{3}{4}$；（2）$a=\pm\sqrt2$，$b=\dfrac{1}{2}$；（3）$a=0$ 时 $b\in\R$，$a\neq 0$ 时 $b\in(-\infty,1)$}

\ifanswer
\solbegin
\stepm{（1）}{当 $a=2$ 时，$A=\left\{\dfrac{1}{2}\right\}$，}
\step{由 $A\sub B$ 得 $\dfrac{1}{2}$ 是方程 $x^2-2x+b=0$ 的解，代入得 $\left(\dfrac{1}{2}\right)^2-2\times\dfrac{1}{2}+b=0$，}
\step{解得 $b=\dfrac{3}{4}$.}
\stepm{（2）}{当 $a^2\geqslant 4b$ 时，方程 $x^2-ax+b=0$ 有实根.}
\step{若 $a=0$，则 $A=\e$，要 $B\sub A$ 需 $B=\e$，}
\step{但此时方程 $x^2+b=0$ 在 $a^2\geqslant 4b$（即 $b\leqslant 0$）时有实根，矛盾，舍去.}
\step{若 $a\neq 0$，则 $A=\left\{\dfrac{1}{a}\right\}$，$B\sub A$ 且方程有实根，故方程有且仅有一个根 $\dfrac{1}{a}$.}
\step{由韦达定理，根的和为 $a=\dfrac{1}{a}+\dfrac{1}{a}$，根的积为 $b=\dfrac{1}{a}\times\dfrac{1}{a}$.}
\step{由根的和得 $a=\dfrac{2}{a}$，即 $a^2=2$，$a=\pm\sqrt2$，$b=\dfrac{1}{2}$.}
\stepm{（3）}{当 $a=0$ 时，$A=\e$，故 $A\sub B$ 恒成立，此时 $b\in\R$.}
\step{当 $a\neq 0$ 时，$A=\left\{\dfrac{1}{a}\right\}$，将 $x=\dfrac{1}{a}$ 代入 $x^2-ax+b=0$，得 $b=1-\dfrac{1}{a^2}$.}
\step{此时方程判别式 $\Delta=a^2-4b=\dfrac{(a^2-2)^2}{a^2}\geqslant 0$，}
\step{若 $\Delta=0$，则 $a=\pm\sqrt2$，$b=\dfrac{1}{2}$；}
\step{若 $\Delta>0$，则 $a\neq 0$ 且 $a\neq\pm\sqrt2$，此时 $b=1-\dfrac{1}{a^2}<1$.}
\step{综上，当 $a=0$ 时，$b\in\R$；当 $a\neq 0$ 时，$b\in(-\infty,1)$.}
\solend

\fi

\ansspace*[12]

\end{enumerate}

\end{document}
"
% 紧凑版：xelatex "\def\iscompact{1}% 高一10_复附浦东_周末作业1.tex —— 高一上数学“2026-2027 年复附浦东高一上周末作业 1”（试卷版/答案版）
% 试卷版：xelatex 高一10_复附浦东_周末作业1.tex
% 答案版：xelatex "\def\isanswer{1}% 高一10_复附浦东_周末作业1.tex —— 高一上数学“2026-2027 年复附浦东高一上周末作业 1”（试卷版/答案版）
% 试卷版：xelatex 高一10_复附浦东_周末作业1.tex
% 答案版：xelatex "\def\isanswer{1}\input{高一10_复附浦东_周末作业1.tex}"
% 紧凑版：xelatex "\def\iscompact{1}\input{高一10_复附浦东_周末作业1.tex}"
%
% 原始材料是微信公众号图文（公众号“魔都数学加油站”），正文为 9 张 PNG 页；原文本身就是一份
% **讲评版**——题目下方直接印出【答案】或【解析】。试题文字、答案与解析均照原卷录入，未改内容。
%
% 与原卷的排版差异（均已在 README“说明”中交代）：
%   ① 原文自**第 3 题**起（第 1、2 题未在该文中发布），故本稿题号亦自 3 起；
%   ② 原卷本就印有“一、填空题”“二、选择题”“三、解答题”三节标题（分别在原图 01/04/06.png
%      上，逐页放大核对过），本稿照原卷分节，只把标题改排成全稿统一的“大节标题”样式；
%   ③ 原卷填空、选择的答案直接印在题目下方，本稿统一改为试卷版隐去、答案版以【答案】给出，
%      使两版彻底分离；
%   ④ 原卷的实数集符号时作正体 $R$、时作粗体 $\mathbf{R}$（第 21(3) 与 (3) 末句均为粗体），
%      本稿按全稿惯例统一写作 $\R$.
%
% 原卷笔误一处（本稿已补，见该处上方的 `%` 注）：
%   ① 第 10 题【解析】首句原卷印作“这个元不是 $1$ 就是 $m$”，“元素”漏了“素”字——
%      本稿按 高一b 第 13 题“原卷漏排、本稿补上”的先例补出“素”字，其余照原卷。
% 除此一处外，本卷 9 页逐页放大核对，未发现别的笔误（题 16 的 $1995=15\times 133$ 与
% “$15^2$ 的倍数有 8 个”两处最易被误判为笔误，实测 $133$、$8$ 均正确）。

% 本篇在公众号“魔都数学加油站”连载里的编号是“27学年高一10”，本地编号与之对齐（高一10）；
% 对应关系见 README“与公众号连载号的对照表”。
\documentclass[a4paper,11pt]{article}
\usepackage{common}

\begin{document}

\examtitlest[3]{2026--2027 年复附浦东高一上周末作业 1}{集合与常用逻辑用语}

\bigsec{一、填空题}

\begin{enumerate}
\setcounter{enumi}{2}

\item 已知集合 $M=\{1,2m+1\}$，$N=\{-1,m^2\}$，且 $M=N$，则 $m=$ \blankline[2.4em].

\ans{$-1$}

\ifanswer
\solbegin
\step{若 $M=N$，必有 $\begin{cases}2m+1=-1\\ m^2=1\end{cases}$，解得 $m=-1$.}
\solend

\fi

\item 集合 $\{(x,y)\mid x+2y=2,\ x,y\in\N\}$ 用列举法表示为 \blankline[4.4em].

\ans{$\{(0,1),(2,0)\}$}

\ifanswer
\solbegin
\step{由题意得 $x+2y=2$ 可化为 $y=1-\dfrac{x}{2}$，}
\step{由 $y\in\N$，有 $y=1-\dfrac{x}{2}\geqslant 0$，解得 $x\leqslant 2$，}
\step{又由 $x\in\N$，得 $x$ 可能取值为 $0$，$1$，$2$，}
\step{当 $x=0$ 时，得 $y=1$，满足条件，}
\step{当 $x=1$ 时，得 $y=\dfrac{1}{2}\notin\N$，不满足条件，}
\step{当 $x=2$ 时，得 $y=0$，满足条件，}
\step{综上，集合用列举法表示为 $\{(0,1),(2,0)\}$.}
\solend

\fi

\item 已知集合 $A=\{1\}$，集合 $B=\{0,1,2\}$，则满足关系 $A\sub P\sub B$ 的所有集合 $P$ 为 \blankline[4.4em].

\ans{$\{1\},\{0,1\},\{1,2\},\{0,1,2\}$}

\ifanswer
\solbegin
\step{由题意，集合 $P$ 可以是 $\{1\},\{0,1\},\{1,2\},\{0,1,2\}$.}
\solend

\fi

\item 已知集合 $A=\{x\mid a-1\leqslant x\leqslant a+2\}$，$B=\{x\mid 3<x<5\}$，则使 $A\supseteq B$ 成立的实数 $a$ 的取值范围是 \blankline[4.4em].

\ans{$[3,4]$}

\ifanswer
\solbegin
\step{因为 $A\supseteq B$，所以 $\begin{cases}a+2\geqslant 5\\ a-1\leqslant 3\end{cases}$，解得 $3\leqslant a\leqslant 4$.}
\solend

\fi

\item 若集合 $\{x\mid mx^2-4x+1=0\}=\{n\}$，则 $m+n=$ \blankline[2.4em].

\ans{$\dfrac{1}{4}$ 或 $\dfrac{9}{2}$}

\ifanswer
\solbegin
\step{若集合 $\{x\mid mx^2-4x+1=0\}=\{n\}$，则 $mx^2-4x+1=0$ 有唯一实数根 $n$，}
\step{当 $m=0$ 时，得 $n=\dfrac{1}{4}$，$m+n=\dfrac{1}{4}$，}
\step{当 $m\neq 0$ 时，$\Delta=16-4m=0$，即 $m=4$，}
\step{此时 $4x^2-4x+1=0$ 的根为 $n=\dfrac{1}{2}$，$m+n=\dfrac{9}{2}$.}
\step{故 $m+n=\dfrac{1}{4}$ 或 $\dfrac{9}{2}$.}
\solend

\fi

\item 集合 $M=\{-1,2,3,5,6\}$ 的全部非空子集的元素和等于 \blankline[4.4em].

\ans{$240$}

\ifanswer
\solbegin
\step{集合 $M=\{-1,2,3,5,6\}$ 中共有 $5$ 个元素，}
\step{在其所有子集中，元素 $-1$ 出现的次数为 $2^{5-1}=2^4=16$，}
\step{同理 $2$、$3$、$5$、$6$ 都出现 $16$ 次，}
\step{则集合 $M$ 的全部非空子集的元素和为 $(-1+2+3+5+6)\times 16=240$.}
\solend

\fi

\item 设 $A=\{x\mid x^2-3x-4=0\}$，$B=\{x\mid ax-1=0\}$，若 $A\cap B=B$，则实数 $a$ 的值可以为 \blankline[4.4em].

\ans{$0$、$-1$、$\dfrac{1}{4}$}

\ifanswer
\solbegin
\step{$A=\{-1,4\}$，因为 $A\cap B=B$，所以 $B\sub A$，}
\step{①$a=0$ 时，$B=\e$，满足 $B\sub A$；}
\step{②$a\neq 0$ 时，$B=\left\{\dfrac{1}{a}\right\}$，则 $\dfrac{1}{a}=-1$ 或 $\dfrac{1}{a}=4$，所以 $a=-1$ 或 $\dfrac{1}{4}$；}
\step{所以 $a$ 的值可以为 $0$、$-1$、$\dfrac{1}{4}$.}
\solend

\fi

\item 设集合 $A=\{1,m\}$，其中 $m$ 为实数. 令集合 $B=\{a+b\mid a\in A,b\in A\}$，若 $A\cap B$ 恰有一个元素，则 $A\cup B$ 的元素之和为 \blankline[2.4em].

\ans{$5$ 或 $10$}

\ifanswer
\solbegin
% 注：原卷本句印作“这个元不是 $1$ 就是 $m$”——“元素”漏了“素”字。按 高一b 第 13 题
%     “原卷漏排、本稿补上”的先例，本稿补出“素”字，其余照原卷，见文件头“原卷笔误”。
\step{由题意得，若 $A\cap B$ 恰有一个元素，这个元素不是 $1$ 就是 $m$，}
\stepm{（1）}{当 $1\in B$ 时，}
\step{①若 $1=1+m$，则 $m=0$，此时 $A=\{1,0\},B=\{2,1,0\},A\cap B=\{0,1\}$，不满足题意，舍去；}
\step{②若 $1=2m$，则 $m=\dfrac{1}{2}$，此时 $A=\left\{1,\dfrac{1}{2}\right\}$，$B=\left\{2,\dfrac{3}{2},1\right\}$，$A\cap B=\{1\}$，}
\step{满足题意，所以 $A\cup B=\left\{1,\dfrac{1}{2},2,\dfrac{3}{2}\right\}$，其元素之和为 $5$，}
\stepm{（2）}{当 $m\in B$ 时，}
\step{①若 $m=1+m$，则 $m$ 无解，}
\step{②若 $m=2m$，则 $m=0$，此时 $A=\{1,0\},B=\{2,1,0\},A\cap B=\{0,1\}$，不满足题意，}
\step{③若 $m=2$，则 $A=\{1,2\},B=\{2,3,4\},A\cap B=\{2\}$，满足题意，}
\step{所以 $A\cup B=\{1,2,3,4\}$，其元素之和为 $10$，}
\step{综上所述，$A\cup B$ 的元素之和为 $5$ 或 $10$.}
\solend

\fi

\item 设非空集合 $Q\sub M$，当 $Q$ 中所有元素和为偶数时（集合为单元素时和为元素本身），称 $Q$ 是 $M$ 的偶子集. 若集合 $M=\{1,2,3,4,5,6,7\}$，则其偶子集 $Q$ 的个数为 \blankline[2.4em].

\ans{$63$}

\ifanswer
\solbegin
\step{偶子集可以分为只有奇数、偶数及偶数个奇数和偶数的集合，}
\step{奇数有 $4$ 个，偶数有 $3$ 个，}
\step{只有奇数的有 $\mathrm{C}_4^2+1=7$ 个，只有偶数的有 $\mathrm{C}_3^1+\mathrm{C}_3^2+1=7$ 个，}
\step{偶数个奇数和偶数的：两奇一偶 $\mathrm{C}_4^2\cdot\mathrm{C}_3^1=18$ 个，两奇二偶 $\mathrm{C}_4^2\cdot\mathrm{C}_3^2=18$ 个，}
\step{两奇三偶 $\mathrm{C}_4^2\cdot\mathrm{C}_3^3=6$ 个，四奇一偶 $\mathrm{C}_4^4\cdot\mathrm{C}_3^1=3$ 个，四奇二偶 $\mathrm{C}_4^4\cdot\mathrm{C}_3^2=3$ 个，}
\step{四奇三偶 $\mathrm{C}_4^4\cdot\mathrm{C}_3^3=1$ 个，所以共有 $7+7+18+18+6+3+3+1=63$ 个.}
\solend

\fi

\item 已知非空集合 $S=\left\{x\mid -\dfrac{1}{2}\leqslant x\leqslant m\right\}$ 满足：当 $x\in S$ 时，有 $x^2\in S$，则实数 $m$ 的取值范围是 \blankline[4.4em].

\ans{$\left[\dfrac{1}{4},1\right]$}

\ifanswer
\solbegin
\step{因为集合 $S=\left\{x\mid -\dfrac{1}{2}\leqslant x\leqslant m\right\}$ 是非空集合，所以 $m\geqslant -\dfrac{1}{2}$，}
\step{又当 $x\in S$ 时，有 $x^2\in S$，所以 $m^2\leqslant m$ 且 $m\geqslant \left(-\dfrac{1}{2}\right)^2=\dfrac{1}{4}$，所以 $\dfrac{1}{4}\leqslant m\leqslant 1$.}
\solend

\fi

\end{enumerate}

\bigsec{二、选择题}

\begin{enumerate}
\setcounter{enumi}{12}

\item 已知集合 $U=\{1+x,1\}$，$x\in U$，则 $x=$\mbox{（\quad）}

\keepwithprev
A. $-1$\qquad B. $0$\qquad C. $1$\qquad D. $2$

\ans{C}

\ifanswer
\solbegin
\step{因为集合 $U=\{1+x,1\}$，$x\in U$，所以 $1+x=x$ 或 $x=1$，}
\step{当 $x=1+x$ 时，即 $0=1$，此时显然不成立，舍去，}
\step{当 $x=1$，此时集合 $U=\{1,2\}$，满足互异性，}
\step{综上所述，$x=1$. 故选 C.}
\solend

\fi

\item 已知集合 $M=\{m\mid m=a+b\sqrt2,\ a,b\in\Q\}$，则下列四个元素中属于 $M$ 的元素的个数是\mbox{（\quad）}

①$1+\sqrt2\pi$；②$\sqrt{11+6\sqrt2}$；③$\dfrac{1}{2+\sqrt2}$；④$\sqrt{2-\sqrt3}+\sqrt{2+\sqrt3}$

\keepwithprev
A. $4$\qquad B. $3$\qquad C. $2$\qquad D. $1$

\ans{C}

\ifanswer
\solbegin
\step{①$1+\sqrt2\pi$，由于 $\sqrt2$、$\pi$ 都为无理数，因此 $1+\sqrt2\pi\notin M$；}
\step{②$\sqrt{11+6\sqrt2}=\sqrt{11+2\sqrt{18}}=3+\sqrt2\in M$；}
\step{③$\dfrac{1}{2+\sqrt2}=1-\dfrac{1}{2}\sqrt2\in M$；}
\step{④$\sqrt{2-\sqrt3}+\sqrt{2+\sqrt3}=\sqrt{\dfrac{8-2\sqrt{12}}{4}}+\sqrt{\dfrac{8+2\sqrt{12}}{4}}=\dfrac{\sqrt6-\sqrt2}{2}+\dfrac{\sqrt6+\sqrt2}{2}=\sqrt6=\sqrt3\cdot\sqrt2\notin M$；}
\step{则上述四个元素中属于 $M$ 的元素的个数是 $2$. 故选 C.}
\solend

\fi

\item 集合 $A=\{2,3,7,8,9,11\}$，集合 $B=\{1,2,4,6,8\}$，下面说法正确的是\mbox{（\quad）}

\keepwithprev
A. 两个集合的交集是 $\{2,8\}$

B. 两个集合没有并集

C. $B$ 集合的元素无法用 $\{x\mid x=f(k),k\in D\}$ 这种形式表示出来

D. 集合 $A$ 有 $6$ 个子集

\ans{A}

\ifanswer
\solbegin
\step{对于 A，$A\cap B=\{2,8\}$，故 A 正确；}
\step{对于 B，$A\cup B=\{1,2,3,4,6,7,8,9,11\}$，故 B 错误；}
\step{对于 C，集合 $B$ 可以表示为 $B=\{x\mid x=f(k),k\in\{1,2,3,4,5\}\}$，}
\step{其中 $f(1)=1,f(2)=2,f(3)=4,f(4)=6,f(5)=8$，}
\step{即集合 $B$ 可描述为 $\{x\mid x=f(k),k\in D\}$ 形式，故 C 错误；}
\step{对于 D，集合 $A$ 有 $6$ 个元素，其子集总数为 $2^6=64\neq 6$，故 D 错误.}
\step{故选 A.}
\solend

\fi

\item 设 $M=\{1,2,\cdots,1995\}$，$A$ 是 $M$ 的子集且满足条件：当 $x\in A$ 时，$15x\notin A$，则 $A$ 中元素的个数的最大值为\mbox{（\quad）}

\keepwithprev
A. $1862$\qquad B. $1866$\qquad C. $1868$\qquad D. $1870$

\ans{D}

\ifanswer
\solbegin
\step{$1995=15\times 133$. 故取出所有不是 $15$ 的倍数的数，共 $1862$ 个，}
\step{这些数均符合要求.}
\step{在所有 $15$ 的倍数的数中，$15^2$ 的倍数有 $8$ 个，这些数又可以取出，}
\step{这样共取出了 $1870$ 个，即 $|A|\geqslant 1870$.}
\step{又 $\{k,15k\}(k=9,10,11,\cdots,133)$ 中的两个元素不能同时取出，}
\step{故 $|A|\leqslant 1995-133+8=1870$，故选 D.}
\solend

\fi

\end{enumerate}

\bigsec{三、解答题}

\begin{enumerate}
\setcounter{enumi}{16}

\item 已知集合 $A=\{x\mid a\leqslant x\leqslant a+4\}$，$B=\{x\mid x\geqslant 6\}$.

\begin{enumerate}[label=(\arabic*)]
\item 若 $a=3$，求 $A\cup B$；
\item 若 $A\cap B=A$，求实数 $a$ 的取值范围.
\end{enumerate}

\ans{（1）$\{x\mid x\geqslant 3\}$；（2）$[6,+\infty)$}

\ifanswer
\solbegin
\stepm{（1）}{当 $a=3$ 时，$A=\{x\mid 3\leqslant x\leqslant 7\}$，所以 $A\cup B=\{x\mid x\geqslant 3\}$；}
\stepm{（2）}{若 $A\cap B=A$，则 $A\sub B$，所以 $a\geqslant 6$，所以实数 $a$ 的取值范围是 $[6,+\infty)$.}
\solend

\fi

\ansspace[7]

\item 已知集合 $M=\{x\mid (x-a)(x^2-ax+a-1)=0\}$.

\begin{enumerate}[label=(\arabic*)]
\item 若 $a=3$，求集合 $M$；
\item 若集合 $M$ 中各元素之和等于 $3$，求实数 $a$ 的值，并用列举法表示集合 $M$.
\end{enumerate}

\ans{（1）$\{1,2,3\}$；（2）$a=2$ 或 $\dfrac{3}{2}$，此时 $M=\{1,2\}$ 或 $M=\left\{\dfrac{1}{2},1,\dfrac{3}{2}\right\}$}

\ifanswer
\solbegin
\stepm{（1）}{$\{1,2,3\}$；}
\stepm{（2）}{当 $(x-a)(x^2-ax+a-1)=0$ 有重根时，重根只能算作集合的一个元素，}
\step{$(x-a)(x^2-ax+a-1)=(x-a)(x-1)[x-(a-1)]=0$，}
\step{当 $a=1$ 时，得 $M=\{0,1\}$，不符合题意；}
\step{当 $a-1=1$ 时，即 $a=2$ 时，得 $M=\{1,2\}$，符合题意；}
\step{当 $a\neq 1$ 且 $a\neq 2$ 时，此时 $M=\{1,a,a-1\}$，得 $1+a+a-1=3$，}
\step{解得 $a=\dfrac{3}{2}$，此时 $M=\left\{\dfrac{1}{2},1,\dfrac{3}{2}\right\}$，符合题意，}
\step{综上，实数 $a$ 的值为 $2$ 或 $\dfrac{3}{2}$，}
\step{当 $a=2$ 时，$M=\{1,2\}$；当 $a=\dfrac{3}{2}$ 时，$M=\left\{\dfrac{1}{2},1,\dfrac{3}{2}\right\}$.}
\solend

\fi

\ansspace[8]

\item 已知集合 $A=\{x\mid x^2-2x-3=0\}$，集合 $B=\{x\mid ax-1=0\}$. 若 $B\sub A$，求实数 $a$ 的值.

\ans{$a=0$ 或 $a=-1$ 或 $a=\dfrac{1}{3}$}

\ifanswer
\solbegin
\step{当 $B=\e$ 时，$a=0$，满足题意；}
\step{当 $B\neq\e$，即 $a\neq 0$ 时，$B=\left\{\dfrac{1}{a}\right\}$，}
\step{又 $A=\{x\mid x^2-2x-3=0\}=\{x\mid x=-1\text{ 或 }x=3\}$，$B\sub A$，}
\step{所以 $\dfrac{1}{a}=-1$ 或 $\dfrac{1}{a}=3$，所以 $a=-1$ 或 $a=\dfrac{1}{3}$.}
\step{综上所述，$a=0$ 或 $a=-1$ 或 $a=\dfrac{1}{3}$.}
\solend

\fi

\ansspace[8]

\item 设集合 $A=\{x\mid -2\leqslant x\leqslant 5\}$，$B=\{x\mid m+1\leqslant x\leqslant 2m-1\}$.

\begin{enumerate}[label=(\arabic*)]
\item 若 $B\sub A$，求实数 $m$ 的取值范围；
\item 当 $x\in\Z$ 时，求 $A$ 的非空真子集个数；
\item 当 $x\in\R$ 时，不存在元素 $x$ 使 $x\in A$ 与 $x\in B$ 同时成立，求实数 $m$ 的取值范围.
\end{enumerate}

\ans{（1）$\{m\mid m\leqslant 3\}$；（2）$254$；（3）$\{m\mid m<2\text{ 或 }m>4\}$}

\ifanswer
\solbegin
\stepm{（1）}{当 $m+1>2m-1$，即 $m<2$ 时，$B=\e$，满足 $B\sub A$.}
\step{当 $m+1\leqslant 2m-1$，即 $m\geqslant 2$ 时，要使 $B\sub A$ 成立，}
\step{需 $\begin{cases}m+1\geqslant -2\\ 2m-1\leqslant 5\end{cases}$，得 $2\leqslant m\leqslant 3$，}
\step{综上，$m$ 的范围为 $\{m\mid m\leqslant 3\}$；}
\stepm{（2）}{$A=\{x\mid -2\leqslant x\leqslant 5,x\in\Z\}=\{-2,-1,0,1,2,3,4,5\}$，}
\step{则 $A$ 的非空真子集个数为 $2^8-1-1=254$；}
\stepm{（3）}{因为没有元素 $x$ 使 $x\in A$ 与 $x\in B$ 同时成立，所以 $A$ 与 $B$ 交集为空集.}
\step{①若 $B=\e$，即 $m+1>2m-1$，得 $m<2$ 时满足条件；}
\step{②若 $B\neq\e$，则要满足的条件是 $\begin{cases}m+1\leqslant 2m-1\\ m+1>5\end{cases}$ 或 $\begin{cases}m+1\leqslant 2m-1\\ 2m-1<-2\end{cases}$，}
\step{解得 $m>4$.}
\step{综上，有 $m<2$ 或 $m>4$，所以 $m$ 的范围为 $\{m\mid m<2\text{ 或 }m>4\}$.}
\solend

\fi

\ansspace[10]

% 压轴题：其后已无内容，故作答区全用可丢弃胶水（\ansspace*），并在题目前先量一下版面。
\ansneed{7cm}

\item 已知集合 $A=\{x\mid ax=1\}$，$B=\{x\mid x^2-ax+b=0\}$.

\begin{enumerate}[label=(\arabic*)]
\item 若 $a=2$，且 $A\sub B$，求 $b$ 的值；
\item 当 $a^2\geqslant 4b$ 时，若 $B\sub A$，求 $a$、$b$ 的值；
\item 若 $A\sub B$，讨论 $a$、$b$ 的取值范围.
\end{enumerate}

\ans{（1）$\dfrac{3}{4}$；（2）$a=\pm\sqrt2$，$b=\dfrac{1}{2}$；（3）$a=0$ 时 $b\in\R$，$a\neq 0$ 时 $b\in(-\infty,1)$}

\ifanswer
\solbegin
\stepm{（1）}{当 $a=2$ 时，$A=\left\{\dfrac{1}{2}\right\}$，}
\step{由 $A\sub B$ 得 $\dfrac{1}{2}$ 是方程 $x^2-2x+b=0$ 的解，代入得 $\left(\dfrac{1}{2}\right)^2-2\times\dfrac{1}{2}+b=0$，}
\step{解得 $b=\dfrac{3}{4}$.}
\stepm{（2）}{当 $a^2\geqslant 4b$ 时，方程 $x^2-ax+b=0$ 有实根.}
\step{若 $a=0$，则 $A=\e$，要 $B\sub A$ 需 $B=\e$，}
\step{但此时方程 $x^2+b=0$ 在 $a^2\geqslant 4b$（即 $b\leqslant 0$）时有实根，矛盾，舍去.}
\step{若 $a\neq 0$，则 $A=\left\{\dfrac{1}{a}\right\}$，$B\sub A$ 且方程有实根，故方程有且仅有一个根 $\dfrac{1}{a}$.}
\step{由韦达定理，根的和为 $a=\dfrac{1}{a}+\dfrac{1}{a}$，根的积为 $b=\dfrac{1}{a}\times\dfrac{1}{a}$.}
\step{由根的和得 $a=\dfrac{2}{a}$，即 $a^2=2$，$a=\pm\sqrt2$，$b=\dfrac{1}{2}$.}
\stepm{（3）}{当 $a=0$ 时，$A=\e$，故 $A\sub B$ 恒成立，此时 $b\in\R$.}
\step{当 $a\neq 0$ 时，$A=\left\{\dfrac{1}{a}\right\}$，将 $x=\dfrac{1}{a}$ 代入 $x^2-ax+b=0$，得 $b=1-\dfrac{1}{a^2}$.}
\step{此时方程判别式 $\Delta=a^2-4b=\dfrac{(a^2-2)^2}{a^2}\geqslant 0$，}
\step{若 $\Delta=0$，则 $a=\pm\sqrt2$，$b=\dfrac{1}{2}$；}
\step{若 $\Delta>0$，则 $a\neq 0$ 且 $a\neq\pm\sqrt2$，此时 $b=1-\dfrac{1}{a^2}<1$.}
\step{综上，当 $a=0$ 时，$b\in\R$；当 $a\neq 0$ 时，$b\in(-\infty,1)$.}
\solend

\fi

\ansspace*[12]

\end{enumerate}

\end{document}
"
% 紧凑版：xelatex "\def\iscompact{1}% 高一10_复附浦东_周末作业1.tex —— 高一上数学“2026-2027 年复附浦东高一上周末作业 1”（试卷版/答案版）
% 试卷版：xelatex 高一10_复附浦东_周末作业1.tex
% 答案版：xelatex "\def\isanswer{1}\input{高一10_复附浦东_周末作业1.tex}"
% 紧凑版：xelatex "\def\iscompact{1}\input{高一10_复附浦东_周末作业1.tex}"
%
% 原始材料是微信公众号图文（公众号“魔都数学加油站”），正文为 9 张 PNG 页；原文本身就是一份
% **讲评版**——题目下方直接印出【答案】或【解析】。试题文字、答案与解析均照原卷录入，未改内容。
%
% 与原卷的排版差异（均已在 README“说明”中交代）：
%   ① 原文自**第 3 题**起（第 1、2 题未在该文中发布），故本稿题号亦自 3 起；
%   ② 原卷本就印有“一、填空题”“二、选择题”“三、解答题”三节标题（分别在原图 01/04/06.png
%      上，逐页放大核对过），本稿照原卷分节，只把标题改排成全稿统一的“大节标题”样式；
%   ③ 原卷填空、选择的答案直接印在题目下方，本稿统一改为试卷版隐去、答案版以【答案】给出，
%      使两版彻底分离；
%   ④ 原卷的实数集符号时作正体 $R$、时作粗体 $\mathbf{R}$（第 21(3) 与 (3) 末句均为粗体），
%      本稿按全稿惯例统一写作 $\R$.
%
% 原卷笔误一处（本稿已补，见该处上方的 `%` 注）：
%   ① 第 10 题【解析】首句原卷印作“这个元不是 $1$ 就是 $m$”，“元素”漏了“素”字——
%      本稿按 高一b 第 13 题“原卷漏排、本稿补上”的先例补出“素”字，其余照原卷。
% 除此一处外，本卷 9 页逐页放大核对，未发现别的笔误（题 16 的 $1995=15\times 133$ 与
% “$15^2$ 的倍数有 8 个”两处最易被误判为笔误，实测 $133$、$8$ 均正确）。

% 本篇在公众号“魔都数学加油站”连载里的编号是“27学年高一10”，本地编号与之对齐（高一10）；
% 对应关系见 README“与公众号连载号的对照表”。
\documentclass[a4paper,11pt]{article}
\usepackage{common}

\begin{document}

\examtitlest[3]{2026--2027 年复附浦东高一上周末作业 1}{集合与常用逻辑用语}

\bigsec{一、填空题}

\begin{enumerate}
\setcounter{enumi}{2}

\item 已知集合 $M=\{1,2m+1\}$，$N=\{-1,m^2\}$，且 $M=N$，则 $m=$ \blankline[2.4em].

\ans{$-1$}

\ifanswer
\solbegin
\step{若 $M=N$，必有 $\begin{cases}2m+1=-1\\ m^2=1\end{cases}$，解得 $m=-1$.}
\solend

\fi

\item 集合 $\{(x,y)\mid x+2y=2,\ x,y\in\N\}$ 用列举法表示为 \blankline[4.4em].

\ans{$\{(0,1),(2,0)\}$}

\ifanswer
\solbegin
\step{由题意得 $x+2y=2$ 可化为 $y=1-\dfrac{x}{2}$，}
\step{由 $y\in\N$，有 $y=1-\dfrac{x}{2}\geqslant 0$，解得 $x\leqslant 2$，}
\step{又由 $x\in\N$，得 $x$ 可能取值为 $0$，$1$，$2$，}
\step{当 $x=0$ 时，得 $y=1$，满足条件，}
\step{当 $x=1$ 时，得 $y=\dfrac{1}{2}\notin\N$，不满足条件，}
\step{当 $x=2$ 时，得 $y=0$，满足条件，}
\step{综上，集合用列举法表示为 $\{(0,1),(2,0)\}$.}
\solend

\fi

\item 已知集合 $A=\{1\}$，集合 $B=\{0,1,2\}$，则满足关系 $A\sub P\sub B$ 的所有集合 $P$ 为 \blankline[4.4em].

\ans{$\{1\},\{0,1\},\{1,2\},\{0,1,2\}$}

\ifanswer
\solbegin
\step{由题意，集合 $P$ 可以是 $\{1\},\{0,1\},\{1,2\},\{0,1,2\}$.}
\solend

\fi

\item 已知集合 $A=\{x\mid a-1\leqslant x\leqslant a+2\}$，$B=\{x\mid 3<x<5\}$，则使 $A\supseteq B$ 成立的实数 $a$ 的取值范围是 \blankline[4.4em].

\ans{$[3,4]$}

\ifanswer
\solbegin
\step{因为 $A\supseteq B$，所以 $\begin{cases}a+2\geqslant 5\\ a-1\leqslant 3\end{cases}$，解得 $3\leqslant a\leqslant 4$.}
\solend

\fi

\item 若集合 $\{x\mid mx^2-4x+1=0\}=\{n\}$，则 $m+n=$ \blankline[2.4em].

\ans{$\dfrac{1}{4}$ 或 $\dfrac{9}{2}$}

\ifanswer
\solbegin
\step{若集合 $\{x\mid mx^2-4x+1=0\}=\{n\}$，则 $mx^2-4x+1=0$ 有唯一实数根 $n$，}
\step{当 $m=0$ 时，得 $n=\dfrac{1}{4}$，$m+n=\dfrac{1}{4}$，}
\step{当 $m\neq 0$ 时，$\Delta=16-4m=0$，即 $m=4$，}
\step{此时 $4x^2-4x+1=0$ 的根为 $n=\dfrac{1}{2}$，$m+n=\dfrac{9}{2}$.}
\step{故 $m+n=\dfrac{1}{4}$ 或 $\dfrac{9}{2}$.}
\solend

\fi

\item 集合 $M=\{-1,2,3,5,6\}$ 的全部非空子集的元素和等于 \blankline[4.4em].

\ans{$240$}

\ifanswer
\solbegin
\step{集合 $M=\{-1,2,3,5,6\}$ 中共有 $5$ 个元素，}
\step{在其所有子集中，元素 $-1$ 出现的次数为 $2^{5-1}=2^4=16$，}
\step{同理 $2$、$3$、$5$、$6$ 都出现 $16$ 次，}
\step{则集合 $M$ 的全部非空子集的元素和为 $(-1+2+3+5+6)\times 16=240$.}
\solend

\fi

\item 设 $A=\{x\mid x^2-3x-4=0\}$，$B=\{x\mid ax-1=0\}$，若 $A\cap B=B$，则实数 $a$ 的值可以为 \blankline[4.4em].

\ans{$0$、$-1$、$\dfrac{1}{4}$}

\ifanswer
\solbegin
\step{$A=\{-1,4\}$，因为 $A\cap B=B$，所以 $B\sub A$，}
\step{①$a=0$ 时，$B=\e$，满足 $B\sub A$；}
\step{②$a\neq 0$ 时，$B=\left\{\dfrac{1}{a}\right\}$，则 $\dfrac{1}{a}=-1$ 或 $\dfrac{1}{a}=4$，所以 $a=-1$ 或 $\dfrac{1}{4}$；}
\step{所以 $a$ 的值可以为 $0$、$-1$、$\dfrac{1}{4}$.}
\solend

\fi

\item 设集合 $A=\{1,m\}$，其中 $m$ 为实数. 令集合 $B=\{a+b\mid a\in A,b\in A\}$，若 $A\cap B$ 恰有一个元素，则 $A\cup B$ 的元素之和为 \blankline[2.4em].

\ans{$5$ 或 $10$}

\ifanswer
\solbegin
% 注：原卷本句印作“这个元不是 $1$ 就是 $m$”——“元素”漏了“素”字。按 高一b 第 13 题
%     “原卷漏排、本稿补上”的先例，本稿补出“素”字，其余照原卷，见文件头“原卷笔误”。
\step{由题意得，若 $A\cap B$ 恰有一个元素，这个元素不是 $1$ 就是 $m$，}
\stepm{（1）}{当 $1\in B$ 时，}
\step{①若 $1=1+m$，则 $m=0$，此时 $A=\{1,0\},B=\{2,1,0\},A\cap B=\{0,1\}$，不满足题意，舍去；}
\step{②若 $1=2m$，则 $m=\dfrac{1}{2}$，此时 $A=\left\{1,\dfrac{1}{2}\right\}$，$B=\left\{2,\dfrac{3}{2},1\right\}$，$A\cap B=\{1\}$，}
\step{满足题意，所以 $A\cup B=\left\{1,\dfrac{1}{2},2,\dfrac{3}{2}\right\}$，其元素之和为 $5$，}
\stepm{（2）}{当 $m\in B$ 时，}
\step{①若 $m=1+m$，则 $m$ 无解，}
\step{②若 $m=2m$，则 $m=0$，此时 $A=\{1,0\},B=\{2,1,0\},A\cap B=\{0,1\}$，不满足题意，}
\step{③若 $m=2$，则 $A=\{1,2\},B=\{2,3,4\},A\cap B=\{2\}$，满足题意，}
\step{所以 $A\cup B=\{1,2,3,4\}$，其元素之和为 $10$，}
\step{综上所述，$A\cup B$ 的元素之和为 $5$ 或 $10$.}
\solend

\fi

\item 设非空集合 $Q\sub M$，当 $Q$ 中所有元素和为偶数时（集合为单元素时和为元素本身），称 $Q$ 是 $M$ 的偶子集. 若集合 $M=\{1,2,3,4,5,6,7\}$，则其偶子集 $Q$ 的个数为 \blankline[2.4em].

\ans{$63$}

\ifanswer
\solbegin
\step{偶子集可以分为只有奇数、偶数及偶数个奇数和偶数的集合，}
\step{奇数有 $4$ 个，偶数有 $3$ 个，}
\step{只有奇数的有 $\mathrm{C}_4^2+1=7$ 个，只有偶数的有 $\mathrm{C}_3^1+\mathrm{C}_3^2+1=7$ 个，}
\step{偶数个奇数和偶数的：两奇一偶 $\mathrm{C}_4^2\cdot\mathrm{C}_3^1=18$ 个，两奇二偶 $\mathrm{C}_4^2\cdot\mathrm{C}_3^2=18$ 个，}
\step{两奇三偶 $\mathrm{C}_4^2\cdot\mathrm{C}_3^3=6$ 个，四奇一偶 $\mathrm{C}_4^4\cdot\mathrm{C}_3^1=3$ 个，四奇二偶 $\mathrm{C}_4^4\cdot\mathrm{C}_3^2=3$ 个，}
\step{四奇三偶 $\mathrm{C}_4^4\cdot\mathrm{C}_3^3=1$ 个，所以共有 $7+7+18+18+6+3+3+1=63$ 个.}
\solend

\fi

\item 已知非空集合 $S=\left\{x\mid -\dfrac{1}{2}\leqslant x\leqslant m\right\}$ 满足：当 $x\in S$ 时，有 $x^2\in S$，则实数 $m$ 的取值范围是 \blankline[4.4em].

\ans{$\left[\dfrac{1}{4},1\right]$}

\ifanswer
\solbegin
\step{因为集合 $S=\left\{x\mid -\dfrac{1}{2}\leqslant x\leqslant m\right\}$ 是非空集合，所以 $m\geqslant -\dfrac{1}{2}$，}
\step{又当 $x\in S$ 时，有 $x^2\in S$，所以 $m^2\leqslant m$ 且 $m\geqslant \left(-\dfrac{1}{2}\right)^2=\dfrac{1}{4}$，所以 $\dfrac{1}{4}\leqslant m\leqslant 1$.}
\solend

\fi

\end{enumerate}

\bigsec{二、选择题}

\begin{enumerate}
\setcounter{enumi}{12}

\item 已知集合 $U=\{1+x,1\}$，$x\in U$，则 $x=$\mbox{（\quad）}

\keepwithprev
A. $-1$\qquad B. $0$\qquad C. $1$\qquad D. $2$

\ans{C}

\ifanswer
\solbegin
\step{因为集合 $U=\{1+x,1\}$，$x\in U$，所以 $1+x=x$ 或 $x=1$，}
\step{当 $x=1+x$ 时，即 $0=1$，此时显然不成立，舍去，}
\step{当 $x=1$，此时集合 $U=\{1,2\}$，满足互异性，}
\step{综上所述，$x=1$. 故选 C.}
\solend

\fi

\item 已知集合 $M=\{m\mid m=a+b\sqrt2,\ a,b\in\Q\}$，则下列四个元素中属于 $M$ 的元素的个数是\mbox{（\quad）}

①$1+\sqrt2\pi$；②$\sqrt{11+6\sqrt2}$；③$\dfrac{1}{2+\sqrt2}$；④$\sqrt{2-\sqrt3}+\sqrt{2+\sqrt3}$

\keepwithprev
A. $4$\qquad B. $3$\qquad C. $2$\qquad D. $1$

\ans{C}

\ifanswer
\solbegin
\step{①$1+\sqrt2\pi$，由于 $\sqrt2$、$\pi$ 都为无理数，因此 $1+\sqrt2\pi\notin M$；}
\step{②$\sqrt{11+6\sqrt2}=\sqrt{11+2\sqrt{18}}=3+\sqrt2\in M$；}
\step{③$\dfrac{1}{2+\sqrt2}=1-\dfrac{1}{2}\sqrt2\in M$；}
\step{④$\sqrt{2-\sqrt3}+\sqrt{2+\sqrt3}=\sqrt{\dfrac{8-2\sqrt{12}}{4}}+\sqrt{\dfrac{8+2\sqrt{12}}{4}}=\dfrac{\sqrt6-\sqrt2}{2}+\dfrac{\sqrt6+\sqrt2}{2}=\sqrt6=\sqrt3\cdot\sqrt2\notin M$；}
\step{则上述四个元素中属于 $M$ 的元素的个数是 $2$. 故选 C.}
\solend

\fi

\item 集合 $A=\{2,3,7,8,9,11\}$，集合 $B=\{1,2,4,6,8\}$，下面说法正确的是\mbox{（\quad）}

\keepwithprev
A. 两个集合的交集是 $\{2,8\}$

B. 两个集合没有并集

C. $B$ 集合的元素无法用 $\{x\mid x=f(k),k\in D\}$ 这种形式表示出来

D. 集合 $A$ 有 $6$ 个子集

\ans{A}

\ifanswer
\solbegin
\step{对于 A，$A\cap B=\{2,8\}$，故 A 正确；}
\step{对于 B，$A\cup B=\{1,2,3,4,6,7,8,9,11\}$，故 B 错误；}
\step{对于 C，集合 $B$ 可以表示为 $B=\{x\mid x=f(k),k\in\{1,2,3,4,5\}\}$，}
\step{其中 $f(1)=1,f(2)=2,f(3)=4,f(4)=6,f(5)=8$，}
\step{即集合 $B$ 可描述为 $\{x\mid x=f(k),k\in D\}$ 形式，故 C 错误；}
\step{对于 D，集合 $A$ 有 $6$ 个元素，其子集总数为 $2^6=64\neq 6$，故 D 错误.}
\step{故选 A.}
\solend

\fi

\item 设 $M=\{1,2,\cdots,1995\}$，$A$ 是 $M$ 的子集且满足条件：当 $x\in A$ 时，$15x\notin A$，则 $A$ 中元素的个数的最大值为\mbox{（\quad）}

\keepwithprev
A. $1862$\qquad B. $1866$\qquad C. $1868$\qquad D. $1870$

\ans{D}

\ifanswer
\solbegin
\step{$1995=15\times 133$. 故取出所有不是 $15$ 的倍数的数，共 $1862$ 个，}
\step{这些数均符合要求.}
\step{在所有 $15$ 的倍数的数中，$15^2$ 的倍数有 $8$ 个，这些数又可以取出，}
\step{这样共取出了 $1870$ 个，即 $|A|\geqslant 1870$.}
\step{又 $\{k,15k\}(k=9,10,11,\cdots,133)$ 中的两个元素不能同时取出，}
\step{故 $|A|\leqslant 1995-133+8=1870$，故选 D.}
\solend

\fi

\end{enumerate}

\bigsec{三、解答题}

\begin{enumerate}
\setcounter{enumi}{16}

\item 已知集合 $A=\{x\mid a\leqslant x\leqslant a+4\}$，$B=\{x\mid x\geqslant 6\}$.

\begin{enumerate}[label=(\arabic*)]
\item 若 $a=3$，求 $A\cup B$；
\item 若 $A\cap B=A$，求实数 $a$ 的取值范围.
\end{enumerate}

\ans{（1）$\{x\mid x\geqslant 3\}$；（2）$[6,+\infty)$}

\ifanswer
\solbegin
\stepm{（1）}{当 $a=3$ 时，$A=\{x\mid 3\leqslant x\leqslant 7\}$，所以 $A\cup B=\{x\mid x\geqslant 3\}$；}
\stepm{（2）}{若 $A\cap B=A$，则 $A\sub B$，所以 $a\geqslant 6$，所以实数 $a$ 的取值范围是 $[6,+\infty)$.}
\solend

\fi

\ansspace[7]

\item 已知集合 $M=\{x\mid (x-a)(x^2-ax+a-1)=0\}$.

\begin{enumerate}[label=(\arabic*)]
\item 若 $a=3$，求集合 $M$；
\item 若集合 $M$ 中各元素之和等于 $3$，求实数 $a$ 的值，并用列举法表示集合 $M$.
\end{enumerate}

\ans{（1）$\{1,2,3\}$；（2）$a=2$ 或 $\dfrac{3}{2}$，此时 $M=\{1,2\}$ 或 $M=\left\{\dfrac{1}{2},1,\dfrac{3}{2}\right\}$}

\ifanswer
\solbegin
\stepm{（1）}{$\{1,2,3\}$；}
\stepm{（2）}{当 $(x-a)(x^2-ax+a-1)=0$ 有重根时，重根只能算作集合的一个元素，}
\step{$(x-a)(x^2-ax+a-1)=(x-a)(x-1)[x-(a-1)]=0$，}
\step{当 $a=1$ 时，得 $M=\{0,1\}$，不符合题意；}
\step{当 $a-1=1$ 时，即 $a=2$ 时，得 $M=\{1,2\}$，符合题意；}
\step{当 $a\neq 1$ 且 $a\neq 2$ 时，此时 $M=\{1,a,a-1\}$，得 $1+a+a-1=3$，}
\step{解得 $a=\dfrac{3}{2}$，此时 $M=\left\{\dfrac{1}{2},1,\dfrac{3}{2}\right\}$，符合题意，}
\step{综上，实数 $a$ 的值为 $2$ 或 $\dfrac{3}{2}$，}
\step{当 $a=2$ 时，$M=\{1,2\}$；当 $a=\dfrac{3}{2}$ 时，$M=\left\{\dfrac{1}{2},1,\dfrac{3}{2}\right\}$.}
\solend

\fi

\ansspace[8]

\item 已知集合 $A=\{x\mid x^2-2x-3=0\}$，集合 $B=\{x\mid ax-1=0\}$. 若 $B\sub A$，求实数 $a$ 的值.

\ans{$a=0$ 或 $a=-1$ 或 $a=\dfrac{1}{3}$}

\ifanswer
\solbegin
\step{当 $B=\e$ 时，$a=0$，满足题意；}
\step{当 $B\neq\e$，即 $a\neq 0$ 时，$B=\left\{\dfrac{1}{a}\right\}$，}
\step{又 $A=\{x\mid x^2-2x-3=0\}=\{x\mid x=-1\text{ 或 }x=3\}$，$B\sub A$，}
\step{所以 $\dfrac{1}{a}=-1$ 或 $\dfrac{1}{a}=3$，所以 $a=-1$ 或 $a=\dfrac{1}{3}$.}
\step{综上所述，$a=0$ 或 $a=-1$ 或 $a=\dfrac{1}{3}$.}
\solend

\fi

\ansspace[8]

\item 设集合 $A=\{x\mid -2\leqslant x\leqslant 5\}$，$B=\{x\mid m+1\leqslant x\leqslant 2m-1\}$.

\begin{enumerate}[label=(\arabic*)]
\item 若 $B\sub A$，求实数 $m$ 的取值范围；
\item 当 $x\in\Z$ 时，求 $A$ 的非空真子集个数；
\item 当 $x\in\R$ 时，不存在元素 $x$ 使 $x\in A$ 与 $x\in B$ 同时成立，求实数 $m$ 的取值范围.
\end{enumerate}

\ans{（1）$\{m\mid m\leqslant 3\}$；（2）$254$；（3）$\{m\mid m<2\text{ 或 }m>4\}$}

\ifanswer
\solbegin
\stepm{（1）}{当 $m+1>2m-1$，即 $m<2$ 时，$B=\e$，满足 $B\sub A$.}
\step{当 $m+1\leqslant 2m-1$，即 $m\geqslant 2$ 时，要使 $B\sub A$ 成立，}
\step{需 $\begin{cases}m+1\geqslant -2\\ 2m-1\leqslant 5\end{cases}$，得 $2\leqslant m\leqslant 3$，}
\step{综上，$m$ 的范围为 $\{m\mid m\leqslant 3\}$；}
\stepm{（2）}{$A=\{x\mid -2\leqslant x\leqslant 5,x\in\Z\}=\{-2,-1,0,1,2,3,4,5\}$，}
\step{则 $A$ 的非空真子集个数为 $2^8-1-1=254$；}
\stepm{（3）}{因为没有元素 $x$ 使 $x\in A$ 与 $x\in B$ 同时成立，所以 $A$ 与 $B$ 交集为空集.}
\step{①若 $B=\e$，即 $m+1>2m-1$，得 $m<2$ 时满足条件；}
\step{②若 $B\neq\e$，则要满足的条件是 $\begin{cases}m+1\leqslant 2m-1\\ m+1>5\end{cases}$ 或 $\begin{cases}m+1\leqslant 2m-1\\ 2m-1<-2\end{cases}$，}
\step{解得 $m>4$.}
\step{综上，有 $m<2$ 或 $m>4$，所以 $m$ 的范围为 $\{m\mid m<2\text{ 或 }m>4\}$.}
\solend

\fi

\ansspace[10]

% 压轴题：其后已无内容，故作答区全用可丢弃胶水（\ansspace*），并在题目前先量一下版面。
\ansneed{7cm}

\item 已知集合 $A=\{x\mid ax=1\}$，$B=\{x\mid x^2-ax+b=0\}$.

\begin{enumerate}[label=(\arabic*)]
\item 若 $a=2$，且 $A\sub B$，求 $b$ 的值；
\item 当 $a^2\geqslant 4b$ 时，若 $B\sub A$，求 $a$、$b$ 的值；
\item 若 $A\sub B$，讨论 $a$、$b$ 的取值范围.
\end{enumerate}

\ans{（1）$\dfrac{3}{4}$；（2）$a=\pm\sqrt2$，$b=\dfrac{1}{2}$；（3）$a=0$ 时 $b\in\R$，$a\neq 0$ 时 $b\in(-\infty,1)$}

\ifanswer
\solbegin
\stepm{（1）}{当 $a=2$ 时，$A=\left\{\dfrac{1}{2}\right\}$，}
\step{由 $A\sub B$ 得 $\dfrac{1}{2}$ 是方程 $x^2-2x+b=0$ 的解，代入得 $\left(\dfrac{1}{2}\right)^2-2\times\dfrac{1}{2}+b=0$，}
\step{解得 $b=\dfrac{3}{4}$.}
\stepm{（2）}{当 $a^2\geqslant 4b$ 时，方程 $x^2-ax+b=0$ 有实根.}
\step{若 $a=0$，则 $A=\e$，要 $B\sub A$ 需 $B=\e$，}
\step{但此时方程 $x^2+b=0$ 在 $a^2\geqslant 4b$（即 $b\leqslant 0$）时有实根，矛盾，舍去.}
\step{若 $a\neq 0$，则 $A=\left\{\dfrac{1}{a}\right\}$，$B\sub A$ 且方程有实根，故方程有且仅有一个根 $\dfrac{1}{a}$.}
\step{由韦达定理，根的和为 $a=\dfrac{1}{a}+\dfrac{1}{a}$，根的积为 $b=\dfrac{1}{a}\times\dfrac{1}{a}$.}
\step{由根的和得 $a=\dfrac{2}{a}$，即 $a^2=2$，$a=\pm\sqrt2$，$b=\dfrac{1}{2}$.}
\stepm{（3）}{当 $a=0$ 时，$A=\e$，故 $A\sub B$ 恒成立，此时 $b\in\R$.}
\step{当 $a\neq 0$ 时，$A=\left\{\dfrac{1}{a}\right\}$，将 $x=\dfrac{1}{a}$ 代入 $x^2-ax+b=0$，得 $b=1-\dfrac{1}{a^2}$.}
\step{此时方程判别式 $\Delta=a^2-4b=\dfrac{(a^2-2)^2}{a^2}\geqslant 0$，}
\step{若 $\Delta=0$，则 $a=\pm\sqrt2$，$b=\dfrac{1}{2}$；}
\step{若 $\Delta>0$，则 $a\neq 0$ 且 $a\neq\pm\sqrt2$，此时 $b=1-\dfrac{1}{a^2}<1$.}
\step{综上，当 $a=0$ 时，$b\in\R$；当 $a\neq 0$ 时，$b\in(-\infty,1)$.}
\solend

\fi

\ansspace*[12]

\end{enumerate}

\end{document}
"
%
% 原始材料是微信公众号图文（公众号“魔都数学加油站”），正文为 9 张 PNG 页；原文本身就是一份
% **讲评版**——题目下方直接印出【答案】或【解析】。试题文字、答案与解析均照原卷录入，未改内容。
%
% 与原卷的排版差异（均已在 README“说明”中交代）：
%   ① 原文自**第 3 题**起（第 1、2 题未在该文中发布），故本稿题号亦自 3 起；
%   ② 原卷本就印有“一、填空题”“二、选择题”“三、解答题”三节标题（分别在原图 01/04/06.png
%      上，逐页放大核对过），本稿照原卷分节，只把标题改排成全稿统一的“大节标题”样式；
%   ③ 原卷填空、选择的答案直接印在题目下方，本稿统一改为试卷版隐去、答案版以【答案】给出，
%      使两版彻底分离；
%   ④ 原卷的实数集符号时作正体 $R$、时作粗体 $\mathbf{R}$（第 21(3) 与 (3) 末句均为粗体），
%      本稿按全稿惯例统一写作 $\R$.
%
% 原卷笔误一处（本稿已补，见该处上方的 `%` 注）：
%   ① 第 10 题【解析】首句原卷印作“这个元不是 $1$ 就是 $m$”，“元素”漏了“素”字——
%      本稿按 高一b 第 13 题“原卷漏排、本稿补上”的先例补出“素”字，其余照原卷。
% 除此一处外，本卷 9 页逐页放大核对，未发现别的笔误（题 16 的 $1995=15\times 133$ 与
% “$15^2$ 的倍数有 8 个”两处最易被误判为笔误，实测 $133$、$8$ 均正确）。

% 本篇在公众号“魔都数学加油站”连载里的编号是“27学年高一10”，本地编号与之对齐（高一10）；
% 对应关系见 README“与公众号连载号的对照表”。
\documentclass[a4paper,11pt]{article}
\usepackage{common}

\begin{document}

\examtitlest[3]{2026--2027 年复附浦东高一上周末作业 1}{集合与常用逻辑用语}

\bigsec{一、填空题}

\begin{enumerate}
\setcounter{enumi}{2}

\item 已知集合 $M=\{1,2m+1\}$，$N=\{-1,m^2\}$，且 $M=N$，则 $m=$ \blankline[2.4em].

\ans{$-1$}

\ifanswer
\solbegin
\step{若 $M=N$，必有 $\begin{cases}2m+1=-1\\ m^2=1\end{cases}$，解得 $m=-1$.}
\solend

\fi

\item 集合 $\{(x,y)\mid x+2y=2,\ x,y\in\N\}$ 用列举法表示为 \blankline[4.4em].

\ans{$\{(0,1),(2,0)\}$}

\ifanswer
\solbegin
\step{由题意得 $x+2y=2$ 可化为 $y=1-\dfrac{x}{2}$，}
\step{由 $y\in\N$，有 $y=1-\dfrac{x}{2}\geqslant 0$，解得 $x\leqslant 2$，}
\step{又由 $x\in\N$，得 $x$ 可能取值为 $0$，$1$，$2$，}
\step{当 $x=0$ 时，得 $y=1$，满足条件，}
\step{当 $x=1$ 时，得 $y=\dfrac{1}{2}\notin\N$，不满足条件，}
\step{当 $x=2$ 时，得 $y=0$，满足条件，}
\step{综上，集合用列举法表示为 $\{(0,1),(2,0)\}$.}
\solend

\fi

\item 已知集合 $A=\{1\}$，集合 $B=\{0,1,2\}$，则满足关系 $A\sub P\sub B$ 的所有集合 $P$ 为 \blankline[4.4em].

\ans{$\{1\},\{0,1\},\{1,2\},\{0,1,2\}$}

\ifanswer
\solbegin
\step{由题意，集合 $P$ 可以是 $\{1\},\{0,1\},\{1,2\},\{0,1,2\}$.}
\solend

\fi

\item 已知集合 $A=\{x\mid a-1\leqslant x\leqslant a+2\}$，$B=\{x\mid 3<x<5\}$，则使 $A\supseteq B$ 成立的实数 $a$ 的取值范围是 \blankline[4.4em].

\ans{$[3,4]$}

\ifanswer
\solbegin
\step{因为 $A\supseteq B$，所以 $\begin{cases}a+2\geqslant 5\\ a-1\leqslant 3\end{cases}$，解得 $3\leqslant a\leqslant 4$.}
\solend

\fi

\item 若集合 $\{x\mid mx^2-4x+1=0\}=\{n\}$，则 $m+n=$ \blankline[2.4em].

\ans{$\dfrac{1}{4}$ 或 $\dfrac{9}{2}$}

\ifanswer
\solbegin
\step{若集合 $\{x\mid mx^2-4x+1=0\}=\{n\}$，则 $mx^2-4x+1=0$ 有唯一实数根 $n$，}
\step{当 $m=0$ 时，得 $n=\dfrac{1}{4}$，$m+n=\dfrac{1}{4}$，}
\step{当 $m\neq 0$ 时，$\Delta=16-4m=0$，即 $m=4$，}
\step{此时 $4x^2-4x+1=0$ 的根为 $n=\dfrac{1}{2}$，$m+n=\dfrac{9}{2}$.}
\step{故 $m+n=\dfrac{1}{4}$ 或 $\dfrac{9}{2}$.}
\solend

\fi

\item 集合 $M=\{-1,2,3,5,6\}$ 的全部非空子集的元素和等于 \blankline[4.4em].

\ans{$240$}

\ifanswer
\solbegin
\step{集合 $M=\{-1,2,3,5,6\}$ 中共有 $5$ 个元素，}
\step{在其所有子集中，元素 $-1$ 出现的次数为 $2^{5-1}=2^4=16$，}
\step{同理 $2$、$3$、$5$、$6$ 都出现 $16$ 次，}
\step{则集合 $M$ 的全部非空子集的元素和为 $(-1+2+3+5+6)\times 16=240$.}
\solend

\fi

\item 设 $A=\{x\mid x^2-3x-4=0\}$，$B=\{x\mid ax-1=0\}$，若 $A\cap B=B$，则实数 $a$ 的值可以为 \blankline[4.4em].

\ans{$0$、$-1$、$\dfrac{1}{4}$}

\ifanswer
\solbegin
\step{$A=\{-1,4\}$，因为 $A\cap B=B$，所以 $B\sub A$，}
\step{①$a=0$ 时，$B=\e$，满足 $B\sub A$；}
\step{②$a\neq 0$ 时，$B=\left\{\dfrac{1}{a}\right\}$，则 $\dfrac{1}{a}=-1$ 或 $\dfrac{1}{a}=4$，所以 $a=-1$ 或 $\dfrac{1}{4}$；}
\step{所以 $a$ 的值可以为 $0$、$-1$、$\dfrac{1}{4}$.}
\solend

\fi

\item 设集合 $A=\{1,m\}$，其中 $m$ 为实数. 令集合 $B=\{a+b\mid a\in A,b\in A\}$，若 $A\cap B$ 恰有一个元素，则 $A\cup B$ 的元素之和为 \blankline[2.4em].

\ans{$5$ 或 $10$}

\ifanswer
\solbegin
% 注：原卷本句印作“这个元不是 $1$ 就是 $m$”——“元素”漏了“素”字。按 高一b 第 13 题
%     “原卷漏排、本稿补上”的先例，本稿补出“素”字，其余照原卷，见文件头“原卷笔误”。
\step{由题意得，若 $A\cap B$ 恰有一个元素，这个元素不是 $1$ 就是 $m$，}
\stepm{（1）}{当 $1\in B$ 时，}
\step{①若 $1=1+m$，则 $m=0$，此时 $A=\{1,0\},B=\{2,1,0\},A\cap B=\{0,1\}$，不满足题意，舍去；}
\step{②若 $1=2m$，则 $m=\dfrac{1}{2}$，此时 $A=\left\{1,\dfrac{1}{2}\right\}$，$B=\left\{2,\dfrac{3}{2},1\right\}$，$A\cap B=\{1\}$，}
\step{满足题意，所以 $A\cup B=\left\{1,\dfrac{1}{2},2,\dfrac{3}{2}\right\}$，其元素之和为 $5$，}
\stepm{（2）}{当 $m\in B$ 时，}
\step{①若 $m=1+m$，则 $m$ 无解，}
\step{②若 $m=2m$，则 $m=0$，此时 $A=\{1,0\},B=\{2,1,0\},A\cap B=\{0,1\}$，不满足题意，}
\step{③若 $m=2$，则 $A=\{1,2\},B=\{2,3,4\},A\cap B=\{2\}$，满足题意，}
\step{所以 $A\cup B=\{1,2,3,4\}$，其元素之和为 $10$，}
\step{综上所述，$A\cup B$ 的元素之和为 $5$ 或 $10$.}
\solend

\fi

\item 设非空集合 $Q\sub M$，当 $Q$ 中所有元素和为偶数时（集合为单元素时和为元素本身），称 $Q$ 是 $M$ 的偶子集. 若集合 $M=\{1,2,3,4,5,6,7\}$，则其偶子集 $Q$ 的个数为 \blankline[2.4em].

\ans{$63$}

\ifanswer
\solbegin
\step{偶子集可以分为只有奇数、偶数及偶数个奇数和偶数的集合，}
\step{奇数有 $4$ 个，偶数有 $3$ 个，}
\step{只有奇数的有 $\mathrm{C}_4^2+1=7$ 个，只有偶数的有 $\mathrm{C}_3^1+\mathrm{C}_3^2+1=7$ 个，}
\step{偶数个奇数和偶数的：两奇一偶 $\mathrm{C}_4^2\cdot\mathrm{C}_3^1=18$ 个，两奇二偶 $\mathrm{C}_4^2\cdot\mathrm{C}_3^2=18$ 个，}
\step{两奇三偶 $\mathrm{C}_4^2\cdot\mathrm{C}_3^3=6$ 个，四奇一偶 $\mathrm{C}_4^4\cdot\mathrm{C}_3^1=3$ 个，四奇二偶 $\mathrm{C}_4^4\cdot\mathrm{C}_3^2=3$ 个，}
\step{四奇三偶 $\mathrm{C}_4^4\cdot\mathrm{C}_3^3=1$ 个，所以共有 $7+7+18+18+6+3+3+1=63$ 个.}
\solend

\fi

\item 已知非空集合 $S=\left\{x\mid -\dfrac{1}{2}\leqslant x\leqslant m\right\}$ 满足：当 $x\in S$ 时，有 $x^2\in S$，则实数 $m$ 的取值范围是 \blankline[4.4em].

\ans{$\left[\dfrac{1}{4},1\right]$}

\ifanswer
\solbegin
\step{因为集合 $S=\left\{x\mid -\dfrac{1}{2}\leqslant x\leqslant m\right\}$ 是非空集合，所以 $m\geqslant -\dfrac{1}{2}$，}
\step{又当 $x\in S$ 时，有 $x^2\in S$，所以 $m^2\leqslant m$ 且 $m\geqslant \left(-\dfrac{1}{2}\right)^2=\dfrac{1}{4}$，所以 $\dfrac{1}{4}\leqslant m\leqslant 1$.}
\solend

\fi

\end{enumerate}

\bigsec{二、选择题}

\begin{enumerate}
\setcounter{enumi}{12}

\item 已知集合 $U=\{1+x,1\}$，$x\in U$，则 $x=$\mbox{（\quad）}

\keepwithprev
A. $-1$\qquad B. $0$\qquad C. $1$\qquad D. $2$

\ans{C}

\ifanswer
\solbegin
\step{因为集合 $U=\{1+x,1\}$，$x\in U$，所以 $1+x=x$ 或 $x=1$，}
\step{当 $x=1+x$ 时，即 $0=1$，此时显然不成立，舍去，}
\step{当 $x=1$，此时集合 $U=\{1,2\}$，满足互异性，}
\step{综上所述，$x=1$. 故选 C.}
\solend

\fi

\item 已知集合 $M=\{m\mid m=a+b\sqrt2,\ a,b\in\Q\}$，则下列四个元素中属于 $M$ 的元素的个数是\mbox{（\quad）}

①$1+\sqrt2\pi$；②$\sqrt{11+6\sqrt2}$；③$\dfrac{1}{2+\sqrt2}$；④$\sqrt{2-\sqrt3}+\sqrt{2+\sqrt3}$

\keepwithprev
A. $4$\qquad B. $3$\qquad C. $2$\qquad D. $1$

\ans{C}

\ifanswer
\solbegin
\step{①$1+\sqrt2\pi$，由于 $\sqrt2$、$\pi$ 都为无理数，因此 $1+\sqrt2\pi\notin M$；}
\step{②$\sqrt{11+6\sqrt2}=\sqrt{11+2\sqrt{18}}=3+\sqrt2\in M$；}
\step{③$\dfrac{1}{2+\sqrt2}=1-\dfrac{1}{2}\sqrt2\in M$；}
\step{④$\sqrt{2-\sqrt3}+\sqrt{2+\sqrt3}=\sqrt{\dfrac{8-2\sqrt{12}}{4}}+\sqrt{\dfrac{8+2\sqrt{12}}{4}}=\dfrac{\sqrt6-\sqrt2}{2}+\dfrac{\sqrt6+\sqrt2}{2}=\sqrt6=\sqrt3\cdot\sqrt2\notin M$；}
\step{则上述四个元素中属于 $M$ 的元素的个数是 $2$. 故选 C.}
\solend

\fi

\item 集合 $A=\{2,3,7,8,9,11\}$，集合 $B=\{1,2,4,6,8\}$，下面说法正确的是\mbox{（\quad）}

\keepwithprev
A. 两个集合的交集是 $\{2,8\}$

B. 两个集合没有并集

C. $B$ 集合的元素无法用 $\{x\mid x=f(k),k\in D\}$ 这种形式表示出来

D. 集合 $A$ 有 $6$ 个子集

\ans{A}

\ifanswer
\solbegin
\step{对于 A，$A\cap B=\{2,8\}$，故 A 正确；}
\step{对于 B，$A\cup B=\{1,2,3,4,6,7,8,9,11\}$，故 B 错误；}
\step{对于 C，集合 $B$ 可以表示为 $B=\{x\mid x=f(k),k\in\{1,2,3,4,5\}\}$，}
\step{其中 $f(1)=1,f(2)=2,f(3)=4,f(4)=6,f(5)=8$，}
\step{即集合 $B$ 可描述为 $\{x\mid x=f(k),k\in D\}$ 形式，故 C 错误；}
\step{对于 D，集合 $A$ 有 $6$ 个元素，其子集总数为 $2^6=64\neq 6$，故 D 错误.}
\step{故选 A.}
\solend

\fi

\item 设 $M=\{1,2,\cdots,1995\}$，$A$ 是 $M$ 的子集且满足条件：当 $x\in A$ 时，$15x\notin A$，则 $A$ 中元素的个数的最大值为\mbox{（\quad）}

\keepwithprev
A. $1862$\qquad B. $1866$\qquad C. $1868$\qquad D. $1870$

\ans{D}

\ifanswer
\solbegin
\step{$1995=15\times 133$. 故取出所有不是 $15$ 的倍数的数，共 $1862$ 个，}
\step{这些数均符合要求.}
\step{在所有 $15$ 的倍数的数中，$15^2$ 的倍数有 $8$ 个，这些数又可以取出，}
\step{这样共取出了 $1870$ 个，即 $|A|\geqslant 1870$.}
\step{又 $\{k,15k\}(k=9,10,11,\cdots,133)$ 中的两个元素不能同时取出，}
\step{故 $|A|\leqslant 1995-133+8=1870$，故选 D.}
\solend

\fi

\end{enumerate}

\bigsec{三、解答题}

\begin{enumerate}
\setcounter{enumi}{16}

\item 已知集合 $A=\{x\mid a\leqslant x\leqslant a+4\}$，$B=\{x\mid x\geqslant 6\}$.

\begin{enumerate}[label=(\arabic*)]
\item 若 $a=3$，求 $A\cup B$；
\item 若 $A\cap B=A$，求实数 $a$ 的取值范围.
\end{enumerate}

\ans{（1）$\{x\mid x\geqslant 3\}$；（2）$[6,+\infty)$}

\ifanswer
\solbegin
\stepm{（1）}{当 $a=3$ 时，$A=\{x\mid 3\leqslant x\leqslant 7\}$，所以 $A\cup B=\{x\mid x\geqslant 3\}$；}
\stepm{（2）}{若 $A\cap B=A$，则 $A\sub B$，所以 $a\geqslant 6$，所以实数 $a$ 的取值范围是 $[6,+\infty)$.}
\solend

\fi

\ansspace[7]

\item 已知集合 $M=\{x\mid (x-a)(x^2-ax+a-1)=0\}$.

\begin{enumerate}[label=(\arabic*)]
\item 若 $a=3$，求集合 $M$；
\item 若集合 $M$ 中各元素之和等于 $3$，求实数 $a$ 的值，并用列举法表示集合 $M$.
\end{enumerate}

\ans{（1）$\{1,2,3\}$；（2）$a=2$ 或 $\dfrac{3}{2}$，此时 $M=\{1,2\}$ 或 $M=\left\{\dfrac{1}{2},1,\dfrac{3}{2}\right\}$}

\ifanswer
\solbegin
\stepm{（1）}{$\{1,2,3\}$；}
\stepm{（2）}{当 $(x-a)(x^2-ax+a-1)=0$ 有重根时，重根只能算作集合的一个元素，}
\step{$(x-a)(x^2-ax+a-1)=(x-a)(x-1)[x-(a-1)]=0$，}
\step{当 $a=1$ 时，得 $M=\{0,1\}$，不符合题意；}
\step{当 $a-1=1$ 时，即 $a=2$ 时，得 $M=\{1,2\}$，符合题意；}
\step{当 $a\neq 1$ 且 $a\neq 2$ 时，此时 $M=\{1,a,a-1\}$，得 $1+a+a-1=3$，}
\step{解得 $a=\dfrac{3}{2}$，此时 $M=\left\{\dfrac{1}{2},1,\dfrac{3}{2}\right\}$，符合题意，}
\step{综上，实数 $a$ 的值为 $2$ 或 $\dfrac{3}{2}$，}
\step{当 $a=2$ 时，$M=\{1,2\}$；当 $a=\dfrac{3}{2}$ 时，$M=\left\{\dfrac{1}{2},1,\dfrac{3}{2}\right\}$.}
\solend

\fi

\ansspace[8]

\item 已知集合 $A=\{x\mid x^2-2x-3=0\}$，集合 $B=\{x\mid ax-1=0\}$. 若 $B\sub A$，求实数 $a$ 的值.

\ans{$a=0$ 或 $a=-1$ 或 $a=\dfrac{1}{3}$}

\ifanswer
\solbegin
\step{当 $B=\e$ 时，$a=0$，满足题意；}
\step{当 $B\neq\e$，即 $a\neq 0$ 时，$B=\left\{\dfrac{1}{a}\right\}$，}
\step{又 $A=\{x\mid x^2-2x-3=0\}=\{x\mid x=-1\text{ 或 }x=3\}$，$B\sub A$，}
\step{所以 $\dfrac{1}{a}=-1$ 或 $\dfrac{1}{a}=3$，所以 $a=-1$ 或 $a=\dfrac{1}{3}$.}
\step{综上所述，$a=0$ 或 $a=-1$ 或 $a=\dfrac{1}{3}$.}
\solend

\fi

\ansspace[8]

\item 设集合 $A=\{x\mid -2\leqslant x\leqslant 5\}$，$B=\{x\mid m+1\leqslant x\leqslant 2m-1\}$.

\begin{enumerate}[label=(\arabic*)]
\item 若 $B\sub A$，求实数 $m$ 的取值范围；
\item 当 $x\in\Z$ 时，求 $A$ 的非空真子集个数；
\item 当 $x\in\R$ 时，不存在元素 $x$ 使 $x\in A$ 与 $x\in B$ 同时成立，求实数 $m$ 的取值范围.
\end{enumerate}

\ans{（1）$\{m\mid m\leqslant 3\}$；（2）$254$；（3）$\{m\mid m<2\text{ 或 }m>4\}$}

\ifanswer
\solbegin
\stepm{（1）}{当 $m+1>2m-1$，即 $m<2$ 时，$B=\e$，满足 $B\sub A$.}
\step{当 $m+1\leqslant 2m-1$，即 $m\geqslant 2$ 时，要使 $B\sub A$ 成立，}
\step{需 $\begin{cases}m+1\geqslant -2\\ 2m-1\leqslant 5\end{cases}$，得 $2\leqslant m\leqslant 3$，}
\step{综上，$m$ 的范围为 $\{m\mid m\leqslant 3\}$；}
\stepm{（2）}{$A=\{x\mid -2\leqslant x\leqslant 5,x\in\Z\}=\{-2,-1,0,1,2,3,4,5\}$，}
\step{则 $A$ 的非空真子集个数为 $2^8-1-1=254$；}
\stepm{（3）}{因为没有元素 $x$ 使 $x\in A$ 与 $x\in B$ 同时成立，所以 $A$ 与 $B$ 交集为空集.}
\step{①若 $B=\e$，即 $m+1>2m-1$，得 $m<2$ 时满足条件；}
\step{②若 $B\neq\e$，则要满足的条件是 $\begin{cases}m+1\leqslant 2m-1\\ m+1>5\end{cases}$ 或 $\begin{cases}m+1\leqslant 2m-1\\ 2m-1<-2\end{cases}$，}
\step{解得 $m>4$.}
\step{综上，有 $m<2$ 或 $m>4$，所以 $m$ 的范围为 $\{m\mid m<2\text{ 或 }m>4\}$.}
\solend

\fi

\ansspace[10]

% 压轴题：其后已无内容，故作答区全用可丢弃胶水（\ansspace*），并在题目前先量一下版面。
\ansneed{7cm}

\item 已知集合 $A=\{x\mid ax=1\}$，$B=\{x\mid x^2-ax+b=0\}$.

\begin{enumerate}[label=(\arabic*)]
\item 若 $a=2$，且 $A\sub B$，求 $b$ 的值；
\item 当 $a^2\geqslant 4b$ 时，若 $B\sub A$，求 $a$、$b$ 的值；
\item 若 $A\sub B$，讨论 $a$、$b$ 的取值范围.
\end{enumerate}

\ans{（1）$\dfrac{3}{4}$；（2）$a=\pm\sqrt2$，$b=\dfrac{1}{2}$；（3）$a=0$ 时 $b\in\R$，$a\neq 0$ 时 $b\in(-\infty,1)$}

\ifanswer
\solbegin
\stepm{（1）}{当 $a=2$ 时，$A=\left\{\dfrac{1}{2}\right\}$，}
\step{由 $A\sub B$ 得 $\dfrac{1}{2}$ 是方程 $x^2-2x+b=0$ 的解，代入得 $\left(\dfrac{1}{2}\right)^2-2\times\dfrac{1}{2}+b=0$，}
\step{解得 $b=\dfrac{3}{4}$.}
\stepm{（2）}{当 $a^2\geqslant 4b$ 时，方程 $x^2-ax+b=0$ 有实根.}
\step{若 $a=0$，则 $A=\e$，要 $B\sub A$ 需 $B=\e$，}
\step{但此时方程 $x^2+b=0$ 在 $a^2\geqslant 4b$（即 $b\leqslant 0$）时有实根，矛盾，舍去.}
\step{若 $a\neq 0$，则 $A=\left\{\dfrac{1}{a}\right\}$，$B\sub A$ 且方程有实根，故方程有且仅有一个根 $\dfrac{1}{a}$.}
\step{由韦达定理，根的和为 $a=\dfrac{1}{a}+\dfrac{1}{a}$，根的积为 $b=\dfrac{1}{a}\times\dfrac{1}{a}$.}
\step{由根的和得 $a=\dfrac{2}{a}$，即 $a^2=2$，$a=\pm\sqrt2$，$b=\dfrac{1}{2}$.}
\stepm{（3）}{当 $a=0$ 时，$A=\e$，故 $A\sub B$ 恒成立，此时 $b\in\R$.}
\step{当 $a\neq 0$ 时，$A=\left\{\dfrac{1}{a}\right\}$，将 $x=\dfrac{1}{a}$ 代入 $x^2-ax+b=0$，得 $b=1-\dfrac{1}{a^2}$.}
\step{此时方程判别式 $\Delta=a^2-4b=\dfrac{(a^2-2)^2}{a^2}\geqslant 0$，}
\step{若 $\Delta=0$，则 $a=\pm\sqrt2$，$b=\dfrac{1}{2}$；}
\step{若 $\Delta>0$，则 $a\neq 0$ 且 $a\neq\pm\sqrt2$，此时 $b=1-\dfrac{1}{a^2}<1$.}
\step{综上，当 $a=0$ 时，$b\in\R$；当 $a\neq 0$ 时，$b\in(-\infty,1)$.}
\solend

\fi

\ansspace*[12]

\end{enumerate}

\end{document}
"
%
% 原始材料是微信公众号图文（公众号“魔都数学加油站”），正文为 9 张 PNG 页；原文本身就是一份
% **讲评版**——题目下方直接印出【答案】或【解析】。试题文字、答案与解析均照原卷录入，未改内容。
%
% 与原卷的排版差异（均已在 README“说明”中交代）：
%   ① 原文自**第 3 题**起（第 1、2 题未在该文中发布），故本稿题号亦自 3 起；
%   ② 原卷本就印有“一、填空题”“二、选择题”“三、解答题”三节标题（分别在原图 01/04/06.png
%      上，逐页放大核对过），本稿照原卷分节，只把标题改排成全稿统一的“大节标题”样式；
%   ③ 原卷填空、选择的答案直接印在题目下方，本稿统一改为试卷版隐去、答案版以【答案】给出，
%      使两版彻底分离；
%   ④ 原卷的实数集符号时作正体 $R$、时作粗体 $\mathbf{R}$（第 21(3) 与 (3) 末句均为粗体），
%      本稿按全稿惯例统一写作 $\R$.
%
% 原卷笔误一处（本稿已补，见该处上方的 `%` 注）：
%   ① 第 10 题【解析】首句原卷印作“这个元不是 $1$ 就是 $m$”，“元素”漏了“素”字——
%      本稿按 高一b 第 13 题“原卷漏排、本稿补上”的先例补出“素”字，其余照原卷。
% 除此一处外，本卷 9 页逐页放大核对，未发现别的笔误（题 16 的 $1995=15\times 133$ 与
% “$15^2$ 的倍数有 8 个”两处最易被误判为笔误，实测 $133$、$8$ 均正确）。

% 本篇在公众号“魔都数学加油站”连载里的编号是“27学年高一10”，本地编号与之对齐（高一10）；
% 对应关系见 README“与公众号连载号的对照表”。
\documentclass[a4paper,11pt]{article}
\usepackage{common}

\begin{document}

\examtitlest[3]{2026--2027 年复附浦东高一上周末作业 1}{集合与常用逻辑用语}

\bigsec{一、填空题}

\begin{enumerate}
\setcounter{enumi}{2}

\item 已知集合 $M=\{1,2m+1\}$，$N=\{-1,m^2\}$，且 $M=N$，则 $m=$ \blankline[2.4em].

\ans{$-1$}

\ifanswer
\solbegin
\step{若 $M=N$，必有 $\begin{cases}2m+1=-1\\ m^2=1\end{cases}$，解得 $m=-1$.}
\solend

\fi

\item 集合 $\{(x,y)\mid x+2y=2,\ x,y\in\N\}$ 用列举法表示为 \blankline[4.4em].

\ans{$\{(0,1),(2,0)\}$}

\ifanswer
\solbegin
\step{由题意得 $x+2y=2$ 可化为 $y=1-\dfrac{x}{2}$，}
\step{由 $y\in\N$，有 $y=1-\dfrac{x}{2}\geqslant 0$，解得 $x\leqslant 2$，}
\step{又由 $x\in\N$，得 $x$ 可能取值为 $0$，$1$，$2$，}
\step{当 $x=0$ 时，得 $y=1$，满足条件，}
\step{当 $x=1$ 时，得 $y=\dfrac{1}{2}\notin\N$，不满足条件，}
\step{当 $x=2$ 时，得 $y=0$，满足条件，}
\step{综上，集合用列举法表示为 $\{(0,1),(2,0)\}$.}
\solend

\fi

\item 已知集合 $A=\{1\}$，集合 $B=\{0,1,2\}$，则满足关系 $A\sub P\sub B$ 的所有集合 $P$ 为 \blankline[4.4em].

\ans{$\{1\},\{0,1\},\{1,2\},\{0,1,2\}$}

\ifanswer
\solbegin
\step{由题意，集合 $P$ 可以是 $\{1\},\{0,1\},\{1,2\},\{0,1,2\}$.}
\solend

\fi

\item 已知集合 $A=\{x\mid a-1\leqslant x\leqslant a+2\}$，$B=\{x\mid 3<x<5\}$，则使 $A\supseteq B$ 成立的实数 $a$ 的取值范围是 \blankline[4.4em].

\ans{$[3,4]$}

\ifanswer
\solbegin
\step{因为 $A\supseteq B$，所以 $\begin{cases}a+2\geqslant 5\\ a-1\leqslant 3\end{cases}$，解得 $3\leqslant a\leqslant 4$.}
\solend

\fi

\item 若集合 $\{x\mid mx^2-4x+1=0\}=\{n\}$，则 $m+n=$ \blankline[2.4em].

\ans{$\dfrac{1}{4}$ 或 $\dfrac{9}{2}$}

\ifanswer
\solbegin
\step{若集合 $\{x\mid mx^2-4x+1=0\}=\{n\}$，则 $mx^2-4x+1=0$ 有唯一实数根 $n$，}
\step{当 $m=0$ 时，得 $n=\dfrac{1}{4}$，$m+n=\dfrac{1}{4}$，}
\step{当 $m\neq 0$ 时，$\Delta=16-4m=0$，即 $m=4$，}
\step{此时 $4x^2-4x+1=0$ 的根为 $n=\dfrac{1}{2}$，$m+n=\dfrac{9}{2}$.}
\step{故 $m+n=\dfrac{1}{4}$ 或 $\dfrac{9}{2}$.}
\solend

\fi

\item 集合 $M=\{-1,2,3,5,6\}$ 的全部非空子集的元素和等于 \blankline[4.4em].

\ans{$240$}

\ifanswer
\solbegin
\step{集合 $M=\{-1,2,3,5,6\}$ 中共有 $5$ 个元素，}
\step{在其所有子集中，元素 $-1$ 出现的次数为 $2^{5-1}=2^4=16$，}
\step{同理 $2$、$3$、$5$、$6$ 都出现 $16$ 次，}
\step{则集合 $M$ 的全部非空子集的元素和为 $(-1+2+3+5+6)\times 16=240$.}
\solend

\fi

\item 设 $A=\{x\mid x^2-3x-4=0\}$，$B=\{x\mid ax-1=0\}$，若 $A\cap B=B$，则实数 $a$ 的值可以为 \blankline[4.4em].

\ans{$0$、$-1$、$\dfrac{1}{4}$}

\ifanswer
\solbegin
\step{$A=\{-1,4\}$，因为 $A\cap B=B$，所以 $B\sub A$，}
\step{①$a=0$ 时，$B=\e$，满足 $B\sub A$；}
\step{②$a\neq 0$ 时，$B=\left\{\dfrac{1}{a}\right\}$，则 $\dfrac{1}{a}=-1$ 或 $\dfrac{1}{a}=4$，所以 $a=-1$ 或 $\dfrac{1}{4}$；}
\step{所以 $a$ 的值可以为 $0$、$-1$、$\dfrac{1}{4}$.}
\solend

\fi

\item 设集合 $A=\{1,m\}$，其中 $m$ 为实数. 令集合 $B=\{a+b\mid a\in A,b\in A\}$，若 $A\cap B$ 恰有一个元素，则 $A\cup B$ 的元素之和为 \blankline[2.4em].

\ans{$5$ 或 $10$}

\ifanswer
\solbegin
% 注：原卷本句印作“这个元不是 $1$ 就是 $m$”——“元素”漏了“素”字。按 高一b 第 13 题
%     “原卷漏排、本稿补上”的先例，本稿补出“素”字，其余照原卷，见文件头“原卷笔误”。
\step{由题意得，若 $A\cap B$ 恰有一个元素，这个元素不是 $1$ 就是 $m$，}
\stepm{（1）}{当 $1\in B$ 时，}
\step{①若 $1=1+m$，则 $m=0$，此时 $A=\{1,0\},B=\{2,1,0\},A\cap B=\{0,1\}$，不满足题意，舍去；}
\step{②若 $1=2m$，则 $m=\dfrac{1}{2}$，此时 $A=\left\{1,\dfrac{1}{2}\right\}$，$B=\left\{2,\dfrac{3}{2},1\right\}$，$A\cap B=\{1\}$，}
\step{满足题意，所以 $A\cup B=\left\{1,\dfrac{1}{2},2,\dfrac{3}{2}\right\}$，其元素之和为 $5$，}
\stepm{（2）}{当 $m\in B$ 时，}
\step{①若 $m=1+m$，则 $m$ 无解，}
\step{②若 $m=2m$，则 $m=0$，此时 $A=\{1,0\},B=\{2,1,0\},A\cap B=\{0,1\}$，不满足题意，}
\step{③若 $m=2$，则 $A=\{1,2\},B=\{2,3,4\},A\cap B=\{2\}$，满足题意，}
\step{所以 $A\cup B=\{1,2,3,4\}$，其元素之和为 $10$，}
\step{综上所述，$A\cup B$ 的元素之和为 $5$ 或 $10$.}
\solend

\fi

\item 设非空集合 $Q\sub M$，当 $Q$ 中所有元素和为偶数时（集合为单元素时和为元素本身），称 $Q$ 是 $M$ 的偶子集. 若集合 $M=\{1,2,3,4,5,6,7\}$，则其偶子集 $Q$ 的个数为 \blankline[2.4em].

\ans{$63$}

\ifanswer
\solbegin
\step{偶子集可以分为只有奇数、偶数及偶数个奇数和偶数的集合，}
\step{奇数有 $4$ 个，偶数有 $3$ 个，}
\step{只有奇数的有 $\mathrm{C}_4^2+1=7$ 个，只有偶数的有 $\mathrm{C}_3^1+\mathrm{C}_3^2+1=7$ 个，}
\step{偶数个奇数和偶数的：两奇一偶 $\mathrm{C}_4^2\cdot\mathrm{C}_3^1=18$ 个，两奇二偶 $\mathrm{C}_4^2\cdot\mathrm{C}_3^2=18$ 个，}
\step{两奇三偶 $\mathrm{C}_4^2\cdot\mathrm{C}_3^3=6$ 个，四奇一偶 $\mathrm{C}_4^4\cdot\mathrm{C}_3^1=3$ 个，四奇二偶 $\mathrm{C}_4^4\cdot\mathrm{C}_3^2=3$ 个，}
\step{四奇三偶 $\mathrm{C}_4^4\cdot\mathrm{C}_3^3=1$ 个，所以共有 $7+7+18+18+6+3+3+1=63$ 个.}
\solend

\fi

\item 已知非空集合 $S=\left\{x\mid -\dfrac{1}{2}\leqslant x\leqslant m\right\}$ 满足：当 $x\in S$ 时，有 $x^2\in S$，则实数 $m$ 的取值范围是 \blankline[4.4em].

\ans{$\left[\dfrac{1}{4},1\right]$}

\ifanswer
\solbegin
\step{因为集合 $S=\left\{x\mid -\dfrac{1}{2}\leqslant x\leqslant m\right\}$ 是非空集合，所以 $m\geqslant -\dfrac{1}{2}$，}
\step{又当 $x\in S$ 时，有 $x^2\in S$，所以 $m^2\leqslant m$ 且 $m\geqslant \left(-\dfrac{1}{2}\right)^2=\dfrac{1}{4}$，所以 $\dfrac{1}{4}\leqslant m\leqslant 1$.}
\solend

\fi

\end{enumerate}

\bigsec{二、选择题}

\begin{enumerate}
\setcounter{enumi}{12}

\item 已知集合 $U=\{1+x,1\}$，$x\in U$，则 $x=$\mbox{（\quad）}

\keepwithprev
A. $-1$\qquad B. $0$\qquad C. $1$\qquad D. $2$

\ans{C}

\ifanswer
\solbegin
\step{因为集合 $U=\{1+x,1\}$，$x\in U$，所以 $1+x=x$ 或 $x=1$，}
\step{当 $x=1+x$ 时，即 $0=1$，此时显然不成立，舍去，}
\step{当 $x=1$，此时集合 $U=\{1,2\}$，满足互异性，}
\step{综上所述，$x=1$. 故选 C.}
\solend

\fi

\item 已知集合 $M=\{m\mid m=a+b\sqrt2,\ a,b\in\Q\}$，则下列四个元素中属于 $M$ 的元素的个数是\mbox{（\quad）}

①$1+\sqrt2\pi$；②$\sqrt{11+6\sqrt2}$；③$\dfrac{1}{2+\sqrt2}$；④$\sqrt{2-\sqrt3}+\sqrt{2+\sqrt3}$

\keepwithprev
A. $4$\qquad B. $3$\qquad C. $2$\qquad D. $1$

\ans{C}

\ifanswer
\solbegin
\step{①$1+\sqrt2\pi$，由于 $\sqrt2$、$\pi$ 都为无理数，因此 $1+\sqrt2\pi\notin M$；}
\step{②$\sqrt{11+6\sqrt2}=\sqrt{11+2\sqrt{18}}=3+\sqrt2\in M$；}
\step{③$\dfrac{1}{2+\sqrt2}=1-\dfrac{1}{2}\sqrt2\in M$；}
\step{④$\sqrt{2-\sqrt3}+\sqrt{2+\sqrt3}=\sqrt{\dfrac{8-2\sqrt{12}}{4}}+\sqrt{\dfrac{8+2\sqrt{12}}{4}}=\dfrac{\sqrt6-\sqrt2}{2}+\dfrac{\sqrt6+\sqrt2}{2}=\sqrt6=\sqrt3\cdot\sqrt2\notin M$；}
\step{则上述四个元素中属于 $M$ 的元素的个数是 $2$. 故选 C.}
\solend

\fi

\item 集合 $A=\{2,3,7,8,9,11\}$，集合 $B=\{1,2,4,6,8\}$，下面说法正确的是\mbox{（\quad）}

\keepwithprev
A. 两个集合的交集是 $\{2,8\}$

B. 两个集合没有并集

C. $B$ 集合的元素无法用 $\{x\mid x=f(k),k\in D\}$ 这种形式表示出来

D. 集合 $A$ 有 $6$ 个子集

\ans{A}

\ifanswer
\solbegin
\step{对于 A，$A\cap B=\{2,8\}$，故 A 正确；}
\step{对于 B，$A\cup B=\{1,2,3,4,6,7,8,9,11\}$，故 B 错误；}
\step{对于 C，集合 $B$ 可以表示为 $B=\{x\mid x=f(k),k\in\{1,2,3,4,5\}\}$，}
\step{其中 $f(1)=1,f(2)=2,f(3)=4,f(4)=6,f(5)=8$，}
\step{即集合 $B$ 可描述为 $\{x\mid x=f(k),k\in D\}$ 形式，故 C 错误；}
\step{对于 D，集合 $A$ 有 $6$ 个元素，其子集总数为 $2^6=64\neq 6$，故 D 错误.}
\step{故选 A.}
\solend

\fi

\item 设 $M=\{1,2,\cdots,1995\}$，$A$ 是 $M$ 的子集且满足条件：当 $x\in A$ 时，$15x\notin A$，则 $A$ 中元素的个数的最大值为\mbox{（\quad）}

\keepwithprev
A. $1862$\qquad B. $1866$\qquad C. $1868$\qquad D. $1870$

\ans{D}

\ifanswer
\solbegin
\step{$1995=15\times 133$. 故取出所有不是 $15$ 的倍数的数，共 $1862$ 个，}
\step{这些数均符合要求.}
\step{在所有 $15$ 的倍数的数中，$15^2$ 的倍数有 $8$ 个，这些数又可以取出，}
\step{这样共取出了 $1870$ 个，即 $|A|\geqslant 1870$.}
\step{又 $\{k,15k\}(k=9,10,11,\cdots,133)$ 中的两个元素不能同时取出，}
\step{故 $|A|\leqslant 1995-133+8=1870$，故选 D.}
\solend

\fi

\end{enumerate}

\bigsec{三、解答题}

\begin{enumerate}
\setcounter{enumi}{16}

\item 已知集合 $A=\{x\mid a\leqslant x\leqslant a+4\}$，$B=\{x\mid x\geqslant 6\}$.

\begin{enumerate}[label=(\arabic*)]
\item 若 $a=3$，求 $A\cup B$；
\item 若 $A\cap B=A$，求实数 $a$ 的取值范围.
\end{enumerate}

\ans{（1）$\{x\mid x\geqslant 3\}$；（2）$[6,+\infty)$}

\ifanswer
\solbegin
\stepm{（1）}{当 $a=3$ 时，$A=\{x\mid 3\leqslant x\leqslant 7\}$，所以 $A\cup B=\{x\mid x\geqslant 3\}$；}
\stepm{（2）}{若 $A\cap B=A$，则 $A\sub B$，所以 $a\geqslant 6$，所以实数 $a$ 的取值范围是 $[6,+\infty)$.}
\solend

\fi

\ansspace[7]

\item 已知集合 $M=\{x\mid (x-a)(x^2-ax+a-1)=0\}$.

\begin{enumerate}[label=(\arabic*)]
\item 若 $a=3$，求集合 $M$；
\item 若集合 $M$ 中各元素之和等于 $3$，求实数 $a$ 的值，并用列举法表示集合 $M$.
\end{enumerate}

\ans{（1）$\{1,2,3\}$；（2）$a=2$ 或 $\dfrac{3}{2}$，此时 $M=\{1,2\}$ 或 $M=\left\{\dfrac{1}{2},1,\dfrac{3}{2}\right\}$}

\ifanswer
\solbegin
\stepm{（1）}{$\{1,2,3\}$；}
\stepm{（2）}{当 $(x-a)(x^2-ax+a-1)=0$ 有重根时，重根只能算作集合的一个元素，}
\step{$(x-a)(x^2-ax+a-1)=(x-a)(x-1)[x-(a-1)]=0$，}
\step{当 $a=1$ 时，得 $M=\{0,1\}$，不符合题意；}
\step{当 $a-1=1$ 时，即 $a=2$ 时，得 $M=\{1,2\}$，符合题意；}
\step{当 $a\neq 1$ 且 $a\neq 2$ 时，此时 $M=\{1,a,a-1\}$，得 $1+a+a-1=3$，}
\step{解得 $a=\dfrac{3}{2}$，此时 $M=\left\{\dfrac{1}{2},1,\dfrac{3}{2}\right\}$，符合题意，}
\step{综上，实数 $a$ 的值为 $2$ 或 $\dfrac{3}{2}$，}
\step{当 $a=2$ 时，$M=\{1,2\}$；当 $a=\dfrac{3}{2}$ 时，$M=\left\{\dfrac{1}{2},1,\dfrac{3}{2}\right\}$.}
\solend

\fi

\ansspace[8]

\item 已知集合 $A=\{x\mid x^2-2x-3=0\}$，集合 $B=\{x\mid ax-1=0\}$. 若 $B\sub A$，求实数 $a$ 的值.

\ans{$a=0$ 或 $a=-1$ 或 $a=\dfrac{1}{3}$}

\ifanswer
\solbegin
\step{当 $B=\e$ 时，$a=0$，满足题意；}
\step{当 $B\neq\e$，即 $a\neq 0$ 时，$B=\left\{\dfrac{1}{a}\right\}$，}
\step{又 $A=\{x\mid x^2-2x-3=0\}=\{x\mid x=-1\text{ 或 }x=3\}$，$B\sub A$，}
\step{所以 $\dfrac{1}{a}=-1$ 或 $\dfrac{1}{a}=3$，所以 $a=-1$ 或 $a=\dfrac{1}{3}$.}
\step{综上所述，$a=0$ 或 $a=-1$ 或 $a=\dfrac{1}{3}$.}
\solend

\fi

\ansspace[8]

\item 设集合 $A=\{x\mid -2\leqslant x\leqslant 5\}$，$B=\{x\mid m+1\leqslant x\leqslant 2m-1\}$.

\begin{enumerate}[label=(\arabic*)]
\item 若 $B\sub A$，求实数 $m$ 的取值范围；
\item 当 $x\in\Z$ 时，求 $A$ 的非空真子集个数；
\item 当 $x\in\R$ 时，不存在元素 $x$ 使 $x\in A$ 与 $x\in B$ 同时成立，求实数 $m$ 的取值范围.
\end{enumerate}

\ans{（1）$\{m\mid m\leqslant 3\}$；（2）$254$；（3）$\{m\mid m<2\text{ 或 }m>4\}$}

\ifanswer
\solbegin
\stepm{（1）}{当 $m+1>2m-1$，即 $m<2$ 时，$B=\e$，满足 $B\sub A$.}
\step{当 $m+1\leqslant 2m-1$，即 $m\geqslant 2$ 时，要使 $B\sub A$ 成立，}
\step{需 $\begin{cases}m+1\geqslant -2\\ 2m-1\leqslant 5\end{cases}$，得 $2\leqslant m\leqslant 3$，}
\step{综上，$m$ 的范围为 $\{m\mid m\leqslant 3\}$；}
\stepm{（2）}{$A=\{x\mid -2\leqslant x\leqslant 5,x\in\Z\}=\{-2,-1,0,1,2,3,4,5\}$，}
\step{则 $A$ 的非空真子集个数为 $2^8-1-1=254$；}
\stepm{（3）}{因为没有元素 $x$ 使 $x\in A$ 与 $x\in B$ 同时成立，所以 $A$ 与 $B$ 交集为空集.}
\step{①若 $B=\e$，即 $m+1>2m-1$，得 $m<2$ 时满足条件；}
\step{②若 $B\neq\e$，则要满足的条件是 $\begin{cases}m+1\leqslant 2m-1\\ m+1>5\end{cases}$ 或 $\begin{cases}m+1\leqslant 2m-1\\ 2m-1<-2\end{cases}$，}
\step{解得 $m>4$.}
\step{综上，有 $m<2$ 或 $m>4$，所以 $m$ 的范围为 $\{m\mid m<2\text{ 或 }m>4\}$.}
\solend

\fi

\ansspace[10]

% 压轴题：其后已无内容，故作答区全用可丢弃胶水（\ansspace*），并在题目前先量一下版面。
\ansneed{7cm}

\item 已知集合 $A=\{x\mid ax=1\}$，$B=\{x\mid x^2-ax+b=0\}$.

\begin{enumerate}[label=(\arabic*)]
\item 若 $a=2$，且 $A\sub B$，求 $b$ 的值；
\item 当 $a^2\geqslant 4b$ 时，若 $B\sub A$，求 $a$、$b$ 的值；
\item 若 $A\sub B$，讨论 $a$、$b$ 的取值范围.
\end{enumerate}

\ans{（1）$\dfrac{3}{4}$；（2）$a=\pm\sqrt2$，$b=\dfrac{1}{2}$；（3）$a=0$ 时 $b\in\R$，$a\neq 0$ 时 $b\in(-\infty,1)$}

\ifanswer
\solbegin
\stepm{（1）}{当 $a=2$ 时，$A=\left\{\dfrac{1}{2}\right\}$，}
\step{由 $A\sub B$ 得 $\dfrac{1}{2}$ 是方程 $x^2-2x+b=0$ 的解，代入得 $\left(\dfrac{1}{2}\right)^2-2\times\dfrac{1}{2}+b=0$，}
\step{解得 $b=\dfrac{3}{4}$.}
\stepm{（2）}{当 $a^2\geqslant 4b$ 时，方程 $x^2-ax+b=0$ 有实根.}
\step{若 $a=0$，则 $A=\e$，要 $B\sub A$ 需 $B=\e$，}
\step{但此时方程 $x^2+b=0$ 在 $a^2\geqslant 4b$（即 $b\leqslant 0$）时有实根，矛盾，舍去.}
\step{若 $a\neq 0$，则 $A=\left\{\dfrac{1}{a}\right\}$，$B\sub A$ 且方程有实根，故方程有且仅有一个根 $\dfrac{1}{a}$.}
\step{由韦达定理，根的和为 $a=\dfrac{1}{a}+\dfrac{1}{a}$，根的积为 $b=\dfrac{1}{a}\times\dfrac{1}{a}$.}
\step{由根的和得 $a=\dfrac{2}{a}$，即 $a^2=2$，$a=\pm\sqrt2$，$b=\dfrac{1}{2}$.}
\stepm{（3）}{当 $a=0$ 时，$A=\e$，故 $A\sub B$ 恒成立，此时 $b\in\R$.}
\step{当 $a\neq 0$ 时，$A=\left\{\dfrac{1}{a}\right\}$，将 $x=\dfrac{1}{a}$ 代入 $x^2-ax+b=0$，得 $b=1-\dfrac{1}{a^2}$.}
\step{此时方程判别式 $\Delta=a^2-4b=\dfrac{(a^2-2)^2}{a^2}\geqslant 0$，}
\step{若 $\Delta=0$，则 $a=\pm\sqrt2$，$b=\dfrac{1}{2}$；}
\step{若 $\Delta>0$，则 $a\neq 0$ 且 $a\neq\pm\sqrt2$，此时 $b=1-\dfrac{1}{a^2}<1$.}
\step{综上，当 $a=0$ 时，$b\in\R$；当 $a\neq 0$ 时，$b\in(-\infty,1)$.}
\solend

\fi

\ansspace*[12]

\end{enumerate}

\end{document}
"
%
% 原始材料是微信公众号图文（公众号“魔都数学加油站”），正文为 9 张 PNG 页；原文本身就是一份
% **讲评版**——题目下方直接印出【答案】或【解析】。试题文字、答案与解析均照原卷录入，未改内容。
%
% 与原卷的排版差异（均已在 README“说明”中交代）：
%   ① 原文自**第 3 题**起（第 1、2 题未在该文中发布），故本稿题号亦自 3 起；
%   ② 原卷本就印有“一、填空题”“二、选择题”“三、解答题”三节标题（分别在原图 01/04/06.png
%      上，逐页放大核对过），本稿照原卷分节，只把标题改排成全稿统一的“大节标题”样式；
%   ③ 原卷填空、选择的答案直接印在题目下方，本稿统一改为试卷版隐去、答案版以【答案】给出，
%      使两版彻底分离；
%   ④ 原卷的实数集符号时作正体 $R$、时作粗体 $\mathbf{R}$（第 21(3) 与 (3) 末句均为粗体），
%      本稿按全稿惯例统一写作 $\R$.
%
% 原卷笔误一处（本稿已补，见该处上方的 `%` 注）：
%   ① 第 10 题【解析】首句原卷印作“这个元不是 $1$ 就是 $m$”，“元素”漏了“素”字——
%      本稿按 高一b 第 13 题“原卷漏排、本稿补上”的先例补出“素”字，其余照原卷。
% 除此一处外，本卷 9 页逐页放大核对，未发现别的笔误（题 16 的 $1995=15\times 133$ 与
% “$15^2$ 的倍数有 8 个”两处最易被误判为笔误，实测 $133$、$8$ 均正确）。

% 本篇在公众号“魔都数学加油站”连载里的编号是“27学年高一10”，本地编号与之对齐（高一10）；
% 对应关系见 README“与公众号连载号的对照表”。
\documentclass[a4paper,11pt]{article}
\usepackage{common}

\begin{document}

\examtitlest[3]{2026--2027 年复附浦东高一上周末作业 1}{集合与常用逻辑用语}

\bigsec{一、填空题}

\begin{enumerate}
\setcounter{enumi}{2}

\item 已知集合 $M=\{1,2m+1\}$，$N=\{-1,m^2\}$，且 $M=N$，则 $m=$ \blankline[2.4em].

\ans{$-1$}

\ifanswer
\solbegin
\step{若 $M=N$，必有 $\begin{cases}2m+1=-1\\ m^2=1\end{cases}$，解得 $m=-1$.}
\solend

\fi

\item 集合 $\{(x,y)\mid x+2y=2,\ x,y\in\N\}$ 用列举法表示为 \blankline[4.4em].

\ans{$\{(0,1),(2,0)\}$}

\ifanswer
\solbegin
\step{由题意得 $x+2y=2$ 可化为 $y=1-\dfrac{x}{2}$，}
\step{由 $y\in\N$，有 $y=1-\dfrac{x}{2}\geqslant 0$，解得 $x\leqslant 2$，}
\step{又由 $x\in\N$，得 $x$ 可能取值为 $0$，$1$，$2$，}
\step{当 $x=0$ 时，得 $y=1$，满足条件，}
\step{当 $x=1$ 时，得 $y=\dfrac{1}{2}\notin\N$，不满足条件，}
\step{当 $x=2$ 时，得 $y=0$，满足条件，}
\step{综上，集合用列举法表示为 $\{(0,1),(2,0)\}$.}
\solend

\fi

\item 已知集合 $A=\{1\}$，集合 $B=\{0,1,2\}$，则满足关系 $A\sub P\sub B$ 的所有集合 $P$ 为 \blankline[4.4em].

\ans{$\{1\},\{0,1\},\{1,2\},\{0,1,2\}$}

\ifanswer
\solbegin
\step{由题意，集合 $P$ 可以是 $\{1\},\{0,1\},\{1,2\},\{0,1,2\}$.}
\solend

\fi

\item 已知集合 $A=\{x\mid a-1\leqslant x\leqslant a+2\}$，$B=\{x\mid 3<x<5\}$，则使 $A\supseteq B$ 成立的实数 $a$ 的取值范围是 \blankline[4.4em].

\ans{$[3,4]$}

\ifanswer
\solbegin
\step{因为 $A\supseteq B$，所以 $\begin{cases}a+2\geqslant 5\\ a-1\leqslant 3\end{cases}$，解得 $3\leqslant a\leqslant 4$.}
\solend

\fi

\item 若集合 $\{x\mid mx^2-4x+1=0\}=\{n\}$，则 $m+n=$ \blankline[2.4em].

\ans{$\dfrac{1}{4}$ 或 $\dfrac{9}{2}$}

\ifanswer
\solbegin
\step{若集合 $\{x\mid mx^2-4x+1=0\}=\{n\}$，则 $mx^2-4x+1=0$ 有唯一实数根 $n$，}
\step{当 $m=0$ 时，得 $n=\dfrac{1}{4}$，$m+n=\dfrac{1}{4}$，}
\step{当 $m\neq 0$ 时，$\Delta=16-4m=0$，即 $m=4$，}
\step{此时 $4x^2-4x+1=0$ 的根为 $n=\dfrac{1}{2}$，$m+n=\dfrac{9}{2}$.}
\step{故 $m+n=\dfrac{1}{4}$ 或 $\dfrac{9}{2}$.}
\solend

\fi

\item 集合 $M=\{-1,2,3,5,6\}$ 的全部非空子集的元素和等于 \blankline[4.4em].

\ans{$240$}

\ifanswer
\solbegin
\step{集合 $M=\{-1,2,3,5,6\}$ 中共有 $5$ 个元素，}
\step{在其所有子集中，元素 $-1$ 出现的次数为 $2^{5-1}=2^4=16$，}
\step{同理 $2$、$3$、$5$、$6$ 都出现 $16$ 次，}
\step{则集合 $M$ 的全部非空子集的元素和为 $(-1+2+3+5+6)\times 16=240$.}
\solend

\fi

\item 设 $A=\{x\mid x^2-3x-4=0\}$，$B=\{x\mid ax-1=0\}$，若 $A\cap B=B$，则实数 $a$ 的值可以为 \blankline[4.4em].

\ans{$0$、$-1$、$\dfrac{1}{4}$}

\ifanswer
\solbegin
\step{$A=\{-1,4\}$，因为 $A\cap B=B$，所以 $B\sub A$，}
\step{①$a=0$ 时，$B=\e$，满足 $B\sub A$；}
\step{②$a\neq 0$ 时，$B=\left\{\dfrac{1}{a}\right\}$，则 $\dfrac{1}{a}=-1$ 或 $\dfrac{1}{a}=4$，所以 $a=-1$ 或 $\dfrac{1}{4}$；}
\step{所以 $a$ 的值可以为 $0$、$-1$、$\dfrac{1}{4}$.}
\solend

\fi

\item 设集合 $A=\{1,m\}$，其中 $m$ 为实数. 令集合 $B=\{a+b\mid a\in A,b\in A\}$，若 $A\cap B$ 恰有一个元素，则 $A\cup B$ 的元素之和为 \blankline[2.4em].

\ans{$5$ 或 $10$}

\ifanswer
\solbegin
% 注：原卷本句印作“这个元不是 $1$ 就是 $m$”——“元素”漏了“素”字。按 高一b 第 13 题
%     “原卷漏排、本稿补上”的先例，本稿补出“素”字，其余照原卷，见文件头“原卷笔误”。
\step{由题意得，若 $A\cap B$ 恰有一个元素，这个元素不是 $1$ 就是 $m$，}
\stepm{（1）}{当 $1\in B$ 时，}
\step{①若 $1=1+m$，则 $m=0$，此时 $A=\{1,0\},B=\{2,1,0\},A\cap B=\{0,1\}$，不满足题意，舍去；}
\step{②若 $1=2m$，则 $m=\dfrac{1}{2}$，此时 $A=\left\{1,\dfrac{1}{2}\right\}$，$B=\left\{2,\dfrac{3}{2},1\right\}$，$A\cap B=\{1\}$，}
\step{满足题意，所以 $A\cup B=\left\{1,\dfrac{1}{2},2,\dfrac{3}{2}\right\}$，其元素之和为 $5$，}
\stepm{（2）}{当 $m\in B$ 时，}
\step{①若 $m=1+m$，则 $m$ 无解，}
\step{②若 $m=2m$，则 $m=0$，此时 $A=\{1,0\},B=\{2,1,0\},A\cap B=\{0,1\}$，不满足题意，}
\step{③若 $m=2$，则 $A=\{1,2\},B=\{2,3,4\},A\cap B=\{2\}$，满足题意，}
\step{所以 $A\cup B=\{1,2,3,4\}$，其元素之和为 $10$，}
\step{综上所述，$A\cup B$ 的元素之和为 $5$ 或 $10$.}
\solend

\fi

\item 设非空集合 $Q\sub M$，当 $Q$ 中所有元素和为偶数时（集合为单元素时和为元素本身），称 $Q$ 是 $M$ 的偶子集. 若集合 $M=\{1,2,3,4,5,6,7\}$，则其偶子集 $Q$ 的个数为 \blankline[2.4em].

\ans{$63$}

\ifanswer
\solbegin
\step{偶子集可以分为只有奇数、偶数及偶数个奇数和偶数的集合，}
\step{奇数有 $4$ 个，偶数有 $3$ 个，}
\step{只有奇数的有 $\mathrm{C}_4^2+1=7$ 个，只有偶数的有 $\mathrm{C}_3^1+\mathrm{C}_3^2+1=7$ 个，}
\step{偶数个奇数和偶数的：两奇一偶 $\mathrm{C}_4^2\cdot\mathrm{C}_3^1=18$ 个，两奇二偶 $\mathrm{C}_4^2\cdot\mathrm{C}_3^2=18$ 个，}
\step{两奇三偶 $\mathrm{C}_4^2\cdot\mathrm{C}_3^3=6$ 个，四奇一偶 $\mathrm{C}_4^4\cdot\mathrm{C}_3^1=3$ 个，四奇二偶 $\mathrm{C}_4^4\cdot\mathrm{C}_3^2=3$ 个，}
\step{四奇三偶 $\mathrm{C}_4^4\cdot\mathrm{C}_3^3=1$ 个，所以共有 $7+7+18+18+6+3+3+1=63$ 个.}
\solend

\fi

\item 已知非空集合 $S=\left\{x\mid -\dfrac{1}{2}\leqslant x\leqslant m\right\}$ 满足：当 $x\in S$ 时，有 $x^2\in S$，则实数 $m$ 的取值范围是 \blankline[4.4em].

\ans{$\left[\dfrac{1}{4},1\right]$}

\ifanswer
\solbegin
\step{因为集合 $S=\left\{x\mid -\dfrac{1}{2}\leqslant x\leqslant m\right\}$ 是非空集合，所以 $m\geqslant -\dfrac{1}{2}$，}
\step{又当 $x\in S$ 时，有 $x^2\in S$，所以 $m^2\leqslant m$ 且 $m\geqslant \left(-\dfrac{1}{2}\right)^2=\dfrac{1}{4}$，所以 $\dfrac{1}{4}\leqslant m\leqslant 1$.}
\solend

\fi

\end{enumerate}

\bigsec{二、选择题}

\begin{enumerate}
\setcounter{enumi}{12}

\item 已知集合 $U=\{1+x,1\}$，$x\in U$，则 $x=$\mbox{（\quad）}

\keepwithprev
A. $-1$\qquad B. $0$\qquad C. $1$\qquad D. $2$

\ans{C}

\ifanswer
\solbegin
\step{因为集合 $U=\{1+x,1\}$，$x\in U$，所以 $1+x=x$ 或 $x=1$，}
\step{当 $x=1+x$ 时，即 $0=1$，此时显然不成立，舍去，}
\step{当 $x=1$，此时集合 $U=\{1,2\}$，满足互异性，}
\step{综上所述，$x=1$. 故选 C.}
\solend

\fi

\item 已知集合 $M=\{m\mid m=a+b\sqrt2,\ a,b\in\Q\}$，则下列四个元素中属于 $M$ 的元素的个数是\mbox{（\quad）}

①$1+\sqrt2\pi$；②$\sqrt{11+6\sqrt2}$；③$\dfrac{1}{2+\sqrt2}$；④$\sqrt{2-\sqrt3}+\sqrt{2+\sqrt3}$

\keepwithprev
A. $4$\qquad B. $3$\qquad C. $2$\qquad D. $1$

\ans{C}

\ifanswer
\solbegin
\step{①$1+\sqrt2\pi$，由于 $\sqrt2$、$\pi$ 都为无理数，因此 $1+\sqrt2\pi\notin M$；}
\step{②$\sqrt{11+6\sqrt2}=\sqrt{11+2\sqrt{18}}=3+\sqrt2\in M$；}
\step{③$\dfrac{1}{2+\sqrt2}=1-\dfrac{1}{2}\sqrt2\in M$；}
\step{④$\sqrt{2-\sqrt3}+\sqrt{2+\sqrt3}=\sqrt{\dfrac{8-2\sqrt{12}}{4}}+\sqrt{\dfrac{8+2\sqrt{12}}{4}}=\dfrac{\sqrt6-\sqrt2}{2}+\dfrac{\sqrt6+\sqrt2}{2}=\sqrt6=\sqrt3\cdot\sqrt2\notin M$；}
\step{则上述四个元素中属于 $M$ 的元素的个数是 $2$. 故选 C.}
\solend

\fi

\item 集合 $A=\{2,3,7,8,9,11\}$，集合 $B=\{1,2,4,6,8\}$，下面说法正确的是\mbox{（\quad）}

\keepwithprev
A. 两个集合的交集是 $\{2,8\}$

B. 两个集合没有并集

C. $B$ 集合的元素无法用 $\{x\mid x=f(k),k\in D\}$ 这种形式表示出来

D. 集合 $A$ 有 $6$ 个子集

\ans{A}

\ifanswer
\solbegin
\step{对于 A，$A\cap B=\{2,8\}$，故 A 正确；}
\step{对于 B，$A\cup B=\{1,2,3,4,6,7,8,9,11\}$，故 B 错误；}
\step{对于 C，集合 $B$ 可以表示为 $B=\{x\mid x=f(k),k\in\{1,2,3,4,5\}\}$，}
\step{其中 $f(1)=1,f(2)=2,f(3)=4,f(4)=6,f(5)=8$，}
\step{即集合 $B$ 可描述为 $\{x\mid x=f(k),k\in D\}$ 形式，故 C 错误；}
\step{对于 D，集合 $A$ 有 $6$ 个元素，其子集总数为 $2^6=64\neq 6$，故 D 错误.}
\step{故选 A.}
\solend

\fi

\item 设 $M=\{1,2,\cdots,1995\}$，$A$ 是 $M$ 的子集且满足条件：当 $x\in A$ 时，$15x\notin A$，则 $A$ 中元素的个数的最大值为\mbox{（\quad）}

\keepwithprev
A. $1862$\qquad B. $1866$\qquad C. $1868$\qquad D. $1870$

\ans{D}

\ifanswer
\solbegin
\step{$1995=15\times 133$. 故取出所有不是 $15$ 的倍数的数，共 $1862$ 个，}
\step{这些数均符合要求.}
\step{在所有 $15$ 的倍数的数中，$15^2$ 的倍数有 $8$ 个，这些数又可以取出，}
\step{这样共取出了 $1870$ 个，即 $|A|\geqslant 1870$.}
\step{又 $\{k,15k\}(k=9,10,11,\cdots,133)$ 中的两个元素不能同时取出，}
\step{故 $|A|\leqslant 1995-133+8=1870$，故选 D.}
\solend

\fi

\end{enumerate}

\bigsec{三、解答题}

\begin{enumerate}
\setcounter{enumi}{16}

\item 已知集合 $A=\{x\mid a\leqslant x\leqslant a+4\}$，$B=\{x\mid x\geqslant 6\}$.

\begin{enumerate}[label=(\arabic*)]
\item 若 $a=3$，求 $A\cup B$；
\item 若 $A\cap B=A$，求实数 $a$ 的取值范围.
\end{enumerate}

\ans{（1）$\{x\mid x\geqslant 3\}$；（2）$[6,+\infty)$}

\ifanswer
\solbegin
\stepm{（1）}{当 $a=3$ 时，$A=\{x\mid 3\leqslant x\leqslant 7\}$，所以 $A\cup B=\{x\mid x\geqslant 3\}$；}
\stepm{（2）}{若 $A\cap B=A$，则 $A\sub B$，所以 $a\geqslant 6$，所以实数 $a$ 的取值范围是 $[6,+\infty)$.}
\solend

\fi

\ansspace[7]

\item 已知集合 $M=\{x\mid (x-a)(x^2-ax+a-1)=0\}$.

\begin{enumerate}[label=(\arabic*)]
\item 若 $a=3$，求集合 $M$；
\item 若集合 $M$ 中各元素之和等于 $3$，求实数 $a$ 的值，并用列举法表示集合 $M$.
\end{enumerate}

\ans{（1）$\{1,2,3\}$；（2）$a=2$ 或 $\dfrac{3}{2}$，此时 $M=\{1,2\}$ 或 $M=\left\{\dfrac{1}{2},1,\dfrac{3}{2}\right\}$}

\ifanswer
\solbegin
\stepm{（1）}{$\{1,2,3\}$；}
\stepm{（2）}{当 $(x-a)(x^2-ax+a-1)=0$ 有重根时，重根只能算作集合的一个元素，}
\step{$(x-a)(x^2-ax+a-1)=(x-a)(x-1)[x-(a-1)]=0$，}
\step{当 $a=1$ 时，得 $M=\{0,1\}$，不符合题意；}
\step{当 $a-1=1$ 时，即 $a=2$ 时，得 $M=\{1,2\}$，符合题意；}
\step{当 $a\neq 1$ 且 $a\neq 2$ 时，此时 $M=\{1,a,a-1\}$，得 $1+a+a-1=3$，}
\step{解得 $a=\dfrac{3}{2}$，此时 $M=\left\{\dfrac{1}{2},1,\dfrac{3}{2}\right\}$，符合题意，}
\step{综上，实数 $a$ 的值为 $2$ 或 $\dfrac{3}{2}$，}
\step{当 $a=2$ 时，$M=\{1,2\}$；当 $a=\dfrac{3}{2}$ 时，$M=\left\{\dfrac{1}{2},1,\dfrac{3}{2}\right\}$.}
\solend

\fi

\ansspace[8]

\item 已知集合 $A=\{x\mid x^2-2x-3=0\}$，集合 $B=\{x\mid ax-1=0\}$. 若 $B\sub A$，求实数 $a$ 的值.

\ans{$a=0$ 或 $a=-1$ 或 $a=\dfrac{1}{3}$}

\ifanswer
\solbegin
\step{当 $B=\e$ 时，$a=0$，满足题意；}
\step{当 $B\neq\e$，即 $a\neq 0$ 时，$B=\left\{\dfrac{1}{a}\right\}$，}
\step{又 $A=\{x\mid x^2-2x-3=0\}=\{x\mid x=-1\text{ 或 }x=3\}$，$B\sub A$，}
\step{所以 $\dfrac{1}{a}=-1$ 或 $\dfrac{1}{a}=3$，所以 $a=-1$ 或 $a=\dfrac{1}{3}$.}
\step{综上所述，$a=0$ 或 $a=-1$ 或 $a=\dfrac{1}{3}$.}
\solend

\fi

\ansspace[8]

\item 设集合 $A=\{x\mid -2\leqslant x\leqslant 5\}$，$B=\{x\mid m+1\leqslant x\leqslant 2m-1\}$.

\begin{enumerate}[label=(\arabic*)]
\item 若 $B\sub A$，求实数 $m$ 的取值范围；
\item 当 $x\in\Z$ 时，求 $A$ 的非空真子集个数；
\item 当 $x\in\R$ 时，不存在元素 $x$ 使 $x\in A$ 与 $x\in B$ 同时成立，求实数 $m$ 的取值范围.
\end{enumerate}

\ans{（1）$\{m\mid m\leqslant 3\}$；（2）$254$；（3）$\{m\mid m<2\text{ 或 }m>4\}$}

\ifanswer
\solbegin
\stepm{（1）}{当 $m+1>2m-1$，即 $m<2$ 时，$B=\e$，满足 $B\sub A$.}
\step{当 $m+1\leqslant 2m-1$，即 $m\geqslant 2$ 时，要使 $B\sub A$ 成立，}
\step{需 $\begin{cases}m+1\geqslant -2\\ 2m-1\leqslant 5\end{cases}$，得 $2\leqslant m\leqslant 3$，}
\step{综上，$m$ 的范围为 $\{m\mid m\leqslant 3\}$；}
\stepm{（2）}{$A=\{x\mid -2\leqslant x\leqslant 5,x\in\Z\}=\{-2,-1,0,1,2,3,4,5\}$，}
\step{则 $A$ 的非空真子集个数为 $2^8-1-1=254$；}
\stepm{（3）}{因为没有元素 $x$ 使 $x\in A$ 与 $x\in B$ 同时成立，所以 $A$ 与 $B$ 交集为空集.}
\step{①若 $B=\e$，即 $m+1>2m-1$，得 $m<2$ 时满足条件；}
\step{②若 $B\neq\e$，则要满足的条件是 $\begin{cases}m+1\leqslant 2m-1\\ m+1>5\end{cases}$ 或 $\begin{cases}m+1\leqslant 2m-1\\ 2m-1<-2\end{cases}$，}
\step{解得 $m>4$.}
\step{综上，有 $m<2$ 或 $m>4$，所以 $m$ 的范围为 $\{m\mid m<2\text{ 或 }m>4\}$.}
\solend

\fi

\ansspace[10]

% 压轴题：其后已无内容，故作答区全用可丢弃胶水（\ansspace*），并在题目前先量一下版面。
\ansneed{7cm}

\item 已知集合 $A=\{x\mid ax=1\}$，$B=\{x\mid x^2-ax+b=0\}$.

\begin{enumerate}[label=(\arabic*)]
\item 若 $a=2$，且 $A\sub B$，求 $b$ 的值；
\item 当 $a^2\geqslant 4b$ 时，若 $B\sub A$，求 $a$、$b$ 的值；
\item 若 $A\sub B$，讨论 $a$、$b$ 的取值范围.
\end{enumerate}

\ans{（1）$\dfrac{3}{4}$；（2）$a=\pm\sqrt2$，$b=\dfrac{1}{2}$；（3）$a=0$ 时 $b\in\R$，$a\neq 0$ 时 $b\in(-\infty,1)$}

\ifanswer
\solbegin
\stepm{（1）}{当 $a=2$ 时，$A=\left\{\dfrac{1}{2}\right\}$，}
\step{由 $A\sub B$ 得 $\dfrac{1}{2}$ 是方程 $x^2-2x+b=0$ 的解，代入得 $\left(\dfrac{1}{2}\right)^2-2\times\dfrac{1}{2}+b=0$，}
\step{解得 $b=\dfrac{3}{4}$.}
\stepm{（2）}{当 $a^2\geqslant 4b$ 时，方程 $x^2-ax+b=0$ 有实根.}
\step{若 $a=0$，则 $A=\e$，要 $B\sub A$ 需 $B=\e$，}
\step{但此时方程 $x^2+b=0$ 在 $a^2\geqslant 4b$（即 $b\leqslant 0$）时有实根，矛盾，舍去.}
\step{若 $a\neq 0$，则 $A=\left\{\dfrac{1}{a}\right\}$，$B\sub A$ 且方程有实根，故方程有且仅有一个根 $\dfrac{1}{a}$.}
\step{由韦达定理，根的和为 $a=\dfrac{1}{a}+\dfrac{1}{a}$，根的积为 $b=\dfrac{1}{a}\times\dfrac{1}{a}$.}
\step{由根的和得 $a=\dfrac{2}{a}$，即 $a^2=2$，$a=\pm\sqrt2$，$b=\dfrac{1}{2}$.}
\stepm{（3）}{当 $a=0$ 时，$A=\e$，故 $A\sub B$ 恒成立，此时 $b\in\R$.}
\step{当 $a\neq 0$ 时，$A=\left\{\dfrac{1}{a}\right\}$，将 $x=\dfrac{1}{a}$ 代入 $x^2-ax+b=0$，得 $b=1-\dfrac{1}{a^2}$.}
\step{此时方程判别式 $\Delta=a^2-4b=\dfrac{(a^2-2)^2}{a^2}\geqslant 0$，}
\step{若 $\Delta=0$，则 $a=\pm\sqrt2$，$b=\dfrac{1}{2}$；}
\step{若 $\Delta>0$，则 $a\neq 0$ 且 $a\neq\pm\sqrt2$，此时 $b=1-\dfrac{1}{a^2}<1$.}
\step{综上，当 $a=0$ 时，$b\in\R$；当 $a\neq 0$ 时，$b\in(-\infty,1)$.}
\solend

\fi

\ansspace*[12]

\end{enumerate}

\end{document}
